\documentclass[11pt]{article}

\usepackage[margin=1.15in]{geometry}
\usepackage[protrusion=true,expansion=false]{microtype}
\usepackage{setspace}
\usepackage{amsmath,amssymb}
\usepackage{amsthm}
\newtheorem{proposition}{Proposition}
\usepackage{booktabs}
\usepackage{csquotes}
\usepackage[hidelinks]{hyperref}
\usepackage{natbib}
\bibpunct{(}{)}{;}{a}{}{,}

\setstretch{1.08}

\title{\textbf{The Contraction of the Possible}\\[0.5em]
\large Enshittification, Artificial Scarcity, and the Manufacture of Replaceability}
\author{Flyxion}
\date{2026}

\begin{document}

\maketitle

\begin{abstract}
Cory Doctorow's theory of \emph{enshittification} describes a recurrent trajectory in which digital platforms first subsidize value for users in order to acquire and retain them, subsequently redirect value toward business customers once users have become difficult to dislodge, and finally extract value from both groups after each has become dependent upon the platform as an intermediary. This essay interprets that trajectory as an instance of a more general structural process: the contraction of the possible. Concentrated platforms acquire power not merely by determining prices, wages, rankings, recommendations, or access to audiences, but by progressively altering the set of alternatives against which those arrangements can meaningfully be evaluated. An alternative may continue to exist formally while becoming economically, socially, or institutionally unreachable. Exit therefore cannot be identified with the mere existence of an exit operation, preference cannot be inferred straightforwardly from choices made within an endogenously constrained choice set, and competition cannot be measured simply by counting nominal alternatives whose use would require abandoning accumulated relationships, reputations, audiences, histories, or livelihoods. Enshittification, on this account, is not merely deterioration under monopoly. It is deterioration accompanied by the progressive removal of the counterfactual discipline that would ordinarily make deterioration costly to the institution producing it.

The argument extends this analysis from consumers to creative workers. Platform opacity produces an extreme asymmetry of legibility in which workers and creators expose increasingly detailed information about their preferences, effort, audiences, responsiveness, economic circumstances, and reservation thresholds while the mechanisms governing ranking, compensation, and distribution remain inaccessible to them. Attempts to infer those mechanisms generate what Doctorow calls deliberately induced superstition: participants construct causal explanations from local observations while acting within an environment whose hidden decision function may itself change in response to their behaviour. Exceptional successes then perform a second institutional function. Doctorow's ``giant teddy bear gambit'' converts rare winners into advertisements for the process that selected them, allowing possibility to substitute for probability and survivorship to masquerade as evidence of merit. Combined with vocational awe and an abundant reserve of people willing to create despite poor expected remuneration, this structure permits platforms to capture value generated not only by labour but by the intrinsic importance that workers attach to the activity itself.

The essay argues further that contemporary artificial intelligence enters this process after a consequential transformation has already occurred. Digital platforms increasingly constrain human expression toward forms that are measurable, rankable, repeatable, and responsive to optimization: retention, engagement, frequency, format, duration, trend conformity, and other machine-legible proxies progressively narrow the space in which creative work is rewarded or even encountered. The subsequent claim that artificial systems can replace human creators therefore risks employing a corrupted counterfactual. Machine production is compared not with the wider space of human creative possibility but with human activity after that activity has already undergone prolonged selection for compatibility with computational measurement. Replaceability can consequently be manufactured before it is discovered. Against this trajectory, the essay develops a conception of refusal, interoperability, competition, modification, collective organization, and antimonopoly policy as mechanisms for preserving reachable alternatives. The opposite of enshittification, on this view, is not simply quality, benevolence, or better corporate governance. It is contestability: the maintenance of a sufficiently rich space of reachable continuations that no institution can unilaterally redefine degradation as success.
\end{abstract}

\section*{Preamble}

There is something peculiar about a system that becomes more profitable as it becomes worse at the activity it ostensibly exists to perform. Ordinary accounts of technological progress encourage us to imagine a rough correspondence between improvement and success: a better search engine should defeat a worse search engine, a service that gives listeners more of the music they actually requested should outperform one that quietly substitutes cheaper material, and a labour market that matches workers and employers efficiently should have little reason to conceal the principles by which compensation is determined. Yet many of the characteristic pathologies of contemporary digital platforms arise precisely where this correspondence has broken down. A search engine may profit when a user must reformulate a query that could have been answered immediately; a streaming platform may improve its margins when the requested work is displaced by something cheaper to license; a social network may benefit from withholding material that a user explicitly chose to follow in favour of material selected according to a different objective; and a labour platform may gain bargaining power by learning how desperate a worker is while ensuring that the worker cannot learn the corresponding reservation threshold of the platform. In each case, what appears from outside as malfunction can be perfectly intelligible from within the objective function of the institution producing it. The question is therefore not simply why technological systems sometimes become worse. It is how institutions acquire the freedom to make them worse without paying the costs that deterioration would ordinarily impose.

Cory Doctorow's theory of \emph{enshittification} \citep{doctorow2023enshittification,doctorow2024enshittification,doctorowgiblin2022chokepoint} provides a political-economic answer. Platforms begin by allocating substantial value to users, often subsidizing participation and tolerating forms of interoperability, openness, generosity, or inefficiency that facilitate rapid adoption. As users accumulate relationships, histories, audiences, reputations, purchases, habits, dependencies, and complementary practices around the platform, leaving becomes progressively more expensive. The platform can then transfer value away from those users and toward the businesses that require access to them. Once those businesses have themselves become dependent upon the intermediary, the direction of extraction changes again: value can be withdrawn from both sides and redirected toward the platform's owners. What remains need not be a good system in any ordinary sense. It need only retain enough residual value, and impose sufficiently large costs upon departure, that each population remains available to the other through the intermediary that is extracting from both. Enshittification is therefore not adequately described as corporate greed, poor product design, or the moral corruption of successful founders. Doctorow's more important claim is structural. Firms have always contained conflicting interests and competing conceptions of what their products should become; what changes is the presence or absence of external disciplines capable of making some choices costly. Competition, regulation, labour organization, interoperability, modification, and the credible possibility of exit constrain what an institution can profitably do. When those constraints disappear, conduct that previously would have destroyed the firm can become the rational means by which it extracts the value accumulated during an earlier and more generous phase.

This essay begins from Doctorow's account but treats enshittification as one instance of a broader transformation in the structure of choice. The decisive transition occurs when an institution acquires influence not merely over which option a person chooses, but over which alternatives remain sufficiently reachable to function as genuine counterfactuals to the option presented. This distinction matters because formal availability is a remarkably weak conception of freedom. A person may technically be able to leave a social network while being unable to transport the social graph through which family relationships, community organizations, customers, professional contacts, and friendships have become coordinated. A musician may technically be free not to participate in a dominant streaming service while discovering that absence from it amounts to absence from the ordinary means by which listeners encounter music. A gig worker may technically reject an offered task while the platform observes those rejections, infers the worker's economic circumstances, and modifies subsequent offers accordingly. A creator may technically ignore the recommendations of an algorithm while access to the audience required to continue creating professionally depends upon satisfying criteria that the algorithm neither discloses nor holds constant. None of these circumstances abolishes choice in the elementary sense. Instead, they change its topology. Paths remain imaginable after they cease to be traversable at tolerable cost.

The resulting distinction between \emph{existence} and \emph{reachability} will be central throughout what follows. An alternative does not discipline an institution merely because it can be named. It disciplines the institution when sufficiently many participants can actually move toward it without surrendering the accumulated goods whose production made their present participation valuable in the first place. This is why lock-in is more than inconvenience and why switching costs are more than another price. A switching cost changes the counterfactual against which the current arrangement is judged. Once departure entails forfeiting an audience, archive, social graph, reputation, livelihood, collection, workflow, or community, the relevant comparison quietly shifts from ``Is this service better than the alternatives?'' to ``Has this service become bad enough that abandoning what I have accumulated is finally worth the loss?'' The institution no longer needs to remain superior to its competitors. It needs only to remain preferable to the destruction or stranding of the dependencies that its earlier success helped to create. Degradation can then proceed through a long interval in which dissatisfaction rises while departure remains irrational. The proverbial frog does not merely fail to notice that the water is becoming hotter. The frog has acquired property, relationships, obligations, and a livelihood inside the pot.

Creative labour makes this structure unusually visible because artists enter platform markets with motivations that cannot be reduced to ordinary wage maximization. People make music, stories, paintings, films, jokes, criticism, scholarship, games, and innumerable other cultural objects even where the expected monetary return is negligible. Doctorow connects this circumstance to Fobazi Ettarh's concept of \emph{vocational awe} \citep{ettarh2018vocationalawe}: work experienced as socially or intrinsically important can permit institutions to impose conditions that workers performing less meaningful labour would reject. The artist's desire to continue making art consequently becomes an exploitable input to the market in which art is distributed. A platform controlling access to audiences need not manufacture the desire to create; it inherits that desire for free. Nor does it need every participant to succeed. Indeed, an enormous population of aspirants, marginal participants, intermittent successes, and uncompensated producers may strengthen the platform's position, provided that enough spectacular successes remain visible to sustain belief in the possibility of ascent. Doctorow's ``giant teddy bear gambit'' \citep{doctorowgiblin2022chokepoint} captures the resulting asymmetry with unusual precision: the exceptional winner is not merely the fortunate outcome of the game but part of the apparatus that persuades everyone else to continue playing it. The existence of a winner is transformed into evidence about the fairness of the mechanism, while the distribution that made the winner exceptional recedes from view.

Artificial intelligence now enters this already distorted environment as both a technological development and an ideological comparison class. Claims that machines can replace writers, illustrators, musicians, performers, programmers, customer-service workers, or other skilled practitioners are commonly framed as though two independently developed capacities had simply converged: human beings perform an activity, machines learn to perform sufficiently similar activity, and substitution follows from the comparison. That description neglects the transformations imposed upon human activity before the comparison is made. Platform-mediated work has for years been subjected to pressures toward regularity, quantification, frequency, responsiveness, stylistic conformity, standardized duration, searchable categorization, retention optimization, engagement maximization, and machine-readable evaluation. Human beings learn to produce what the platform can measure because what the platform can measure increasingly determines what other human beings are permitted to encounter. The worker consequently becomes more legible to the machine at the same time that the machine becomes less legible to the worker. By the time automated production is introduced as a substitute, the human activity against which it is evaluated may already have been narrowed toward precisely those dimensions on which computational substitution is easiest to demonstrate.

The problem, then, is larger than the deterioration of websites and more consequential than whether a particular recommendation algorithm occasionally produces irritating results. It concerns the institutional capacity to contract the space of possible continuation while leaving the vocabulary of choice intact. It concerns the conversion of abundance into scarcity by controlling access rather than production, the conversion of opacity into bargaining advantage, the conversion of exceptional survivors into causal explanations of their own success, and the conversion of intrinsically motivated activity into a reservoir from which economic value can be extracted without corresponding bargaining power. Most importantly, it concerns the difference between a world containing alternatives and a world in which those alternatives remain reachable. If enshittification is understood at this level, then its opposite cannot simply be a platform that behaves better. A benevolent monopolist remains a monopolist, and a temporarily generous intermediary can use its generosity to construct the dependencies that make later extraction possible. The more durable objective is therefore to preserve the conditions under which deterioration encounters resistance: multiple reachable continuations, meaningful capacities for refusal, limits on asymmetric legibility, collective forms of bargaining, and institutions unable to make themselves indispensable merely by accumulating the dependencies of those they mediate. The question running through this essay is consequently not whether digital platforms can be persuaded to give more value back. It is what kinds of social, technical, legal, and representational structure prevent any platform from acquiring the ability to determine, by itself, how little value everyone else must learn to accept.

\section{The Platform That Gets Worse on Purpose}

The most revealing cases of enshittification are not those in which a platform deteriorates because its engineers have become incompetent, its infrastructure has failed to keep pace with demand, or its managers have simply lost sight of what users want. They are cases in which deterioration is itself productive from the standpoint of the institution that causes it. Doctorow's discussion of Google Search \citep{doctorowfox2026} provides an unusually valuable example because the internal conflict that normally has to be reconstructed from external behaviour became partially visible through documents disclosed during antitrust litigation \citep{usgoogleantitrust2024,googlequalityrevenue}. As Doctorow presents the episode, Google's overwhelming share of the search market created an elementary problem of growth. Once nearly everyone who could plausibly be induced to use the service was already using it, increasing revenue by acquiring additional users became progressively more difficult. One response was to improve search further, but successful search contains an inconvenient property from the standpoint of an advertising business: when the correct result is delivered immediately, the interaction ends. A user who finds what they need after one query generates fewer opportunities for advertising than a user who must reformulate the query several times. Search quality and revenue therefore cease to be naturally aligned. Under conditions in which dissatisfied users can readily move elsewhere, deliberately degrading search would be self-defeating. Under conditions of sufficiently strong monopoly, however, the same degradation can increase revenue because the additional searches produced by failure remain inside the same advertising system.

This distinction is more important than the familiar observation that a corporation may sacrifice quality for profit. What changes is the sign of the feedback that ordinarily connects quality to institutional survival. In a competitive environment, a deterioration in the service provided to users can generate a corresponding deterioration in the firm's position: users leave, competitors acquire them, revenue falls, and the organization encounters pressure to correct the defect. The system contains a negative feedback mechanism. Monopoly and lock-in weaken that mechanism by reducing the extent to which dissatisfaction can be translated into departure. Once the probability of exit becomes sufficiently small, the same deterioration may generate additional revenue without generating a proportionate loss of users. What was previously an error becomes an optimization strategy. The platform has not ceased to optimize; it has acquired enough control over the environment that it can optimize against the interests of the population whose behaviour supplies its revenue. This is one reason moral explanations of enshittification are inadequate. The crucial event is not the arrival of unusually malicious executives. It is the alteration of the constraint structure within which executives, engineers, investors, advertisers, workers, and users interact. Conduct that would have been punished under one structure becomes rewarded under another.

The Google example also exposes a problem of measurement. Search quality can be evaluated according to the objective of the person searching: how reliably, quickly, and accurately does the system allow a user to locate what they sought? Yet the same interaction can be evaluated according to an entirely different coordinate system: number of queries, advertisements displayed, sponsored results encountered, time spent within the service, or revenue generated per session. An intervention can therefore register simultaneously as deterioration along one axis and improvement along another. If the institution controlling the system also controls which coordinate determines investment, promotion, compensation, and strategic success, it acquires an evaluative advantage that precedes any particular manipulation of the product. The question ``Did search improve?'' becomes underspecified until one asks \emph{with respect to whose objective and measured along which coordinate?} A system that requires three searches to accomplish what previously required one has deteriorated with respect to task completion while improving with respect to monetizable interaction. Calling the latter an improvement does not require falsifying the measurements. It requires selecting the coordinate along which the institution benefits and allowing that coordinate to stand in for the function the system publicly appears to perform.

This substitution recurs throughout Doctorow's examples. A social network ostensibly mediates relationships among people who have chosen to connect with one another, yet a feed optimized for monetizable attention may suppress material from those relationships in favour of recommendations predicted to increase engagement. A streaming service ostensibly supplies requested music, yet its financial incentives can favour playlists in which cheaper, promotional, or otherwise advantageous recordings displace the works for which listeners initially arrived. A platform for creative work may ostensibly help artists reach audiences, yet its internal objective may favour a continuous supply of inexpensive, rapidly produced material capable of retaining attention regardless of whether that material advances the artistic purposes of the people producing or consuming it. In each case, the service can continue to satisfy enough of its nominal function to preserve participation while increasingly optimizing a second function that would have been difficult to pursue had participants retained inexpensive and interoperable alternatives. Enshittification is therefore not simply a decline in quality. It is a divergence between objective functions made sustainable by an asymmetry of power.

This divergence also explains why satisfaction and engagement cannot be treated as interchangeable measures. A satisfied search can terminate quickly; a satisfying piece of music can be heard without generating an indefinite sequence of further interactions; a meaningful exchange between friends may last only a few minutes; a good book may cause its reader to close every application on the device for several hours. From the standpoint of a system whose revenue is proportional to opportunities for intervention, these can be commercially inferior outcomes precisely because they succeed at what the participant wanted. Satisfaction contains a natural stopping condition. Attention maximization does not. Once the platform's objective shifts from enabling an activity to extending the measurable interaction surrounding that activity, completion itself can become a defect. The user's base case---``I found it,'' ``I heard it,'' ``I spoke to them,'' ``I am finished''---interrupts the recursive production of opportunities to advertise, recommend, rank, measure, and extract. The commercial system therefore acquires an incentive to eliminate or defer the stopping condition while preserving the user's belief that the original activity remains the reason for the interaction.

There is an instructive connection here to the general problem of proxy optimization. A proxy is useful only insofar as its relationship to the underlying quantity remains sufficiently stable. Search frequency can indicate usefulness when people voluntarily return to a search engine because previous searches succeeded; time spent can indicate interest when continued participation is freely chosen; engagement can indicate relevance when the system is not itself manipulating the conditions under which engagement occurs. But once the institution begins optimizing the proxy directly, it can alter the causal relationship that made the proxy informative. More searches may cease to indicate greater usefulness because poor search quality itself produces additional searches. Longer sessions may cease to indicate greater satisfaction because friction, outrage, confusion, or deliberately incomplete resolution can extend them. Increased creator output may cease to indicate a flourishing creative culture because creators have learned that constant production is necessary to avoid algorithmic disappearance. The proxy does not merely become noisy. The optimization process changes what the proxy means.

This is one reason the familiar language of ``user experience'' can obscure more than it reveals. It suggests that the platform and user jointly inhabit a system whose success can be represented by a sufficiently comprehensive measure of satisfaction. Doctorow's account instead describes an environment in which the relevant parties can have systematically opposed objectives while remaining mutually dependent. The user wants to accomplish something; the platform may benefit from prolonging the process by which it is accomplished. The musician wants a listener to encounter a particular work; the intermediary may benefit from redirecting that listener toward a cheaper substitute. The driver wants the highest compensation available for a trip; the labour platform wants to discover the lowest compensation that particular driver will accept. The advertiser wants an advertisement to reach a relevant human customer; the advertising platform can profit from transactions whose effectiveness the advertiser is poorly positioned to verify. These are not merely imperfectly aligned incentives. They are opportunities for one participant to improve its own outcome by learning how much deterioration another participant will tolerate before leaving.

The relevant optimization problem can therefore be expressed in a deliberately simple form. Let \(Q\) denote the value of a service to the participant, \(E\) the extractable value captured by the intermediary, and \(C_{\mathrm{exit}}\) the participant's cost of leaving. A competitive intuition assumes that reducing \(Q\) eventually reduces \(E\), because participants can transfer their activity elsewhere. Lock-in changes this relationship. The platform can lower \(Q\) while increasing \(E\) so long as the loss imposed upon the participant remains smaller than the effective cost of departure. The operative constraint is no longer that the service remain good in an absolute sense, or even that it remain better than technically existing competitors. It need only satisfy something like
\[
Q_{\mathrm{remain}} - Q_{\mathrm{exit}} > -C_{\mathrm{exit}},
\]
where \(Q_{\mathrm{exit}}\) includes whatever value remains available after leaving and \(C_{\mathrm{exit}}\) includes the relationships, data, audiences, histories, reputations, habits, complementary tools, and institutional dependencies that departure would strand. The expression is not intended as a complete economic model. Its purpose is to make explicit what ordinary descriptions of consumer choice often suppress: the platform's room for deterioration grows with the cost of making an alternative trajectory real.

Doctorow's three-stage account can consequently be understood as a sequence in which the intermediary first creates value and then converts accumulated value into switching costs. Generosity during the initial phase is not incidental to later extraction. The things that make a platform worth joining are often the same things that make it costly to leave once they have accumulated. A social graph is valuable because one's friends are present; that value becomes lock-in when the graph cannot travel. An audience is valuable because it allows a creator to reach people; it becomes lock-in when the audience relationship belongs operationally to the platform rather than to the creator. A reputation system is valuable because it permits strangers to transact; it becomes lock-in when years of accumulated reputation disappear upon migration. A purchased library is valuable because it makes media readily accessible; it becomes lock-in when access depends upon remaining within a proprietary environment. A workflow becomes more valuable as complementary tools and institutional practices grow around it; the same complementarities later increase the cost of substitution. Value and dependency are therefore not opposites. Under conditions of enclosure, value can be transformed into dependency through accumulation.

This transformation complicates the common narrative in which users foolishly trade freedom for convenience. The relevant platforms frequently are convenient, useful, innovative, or even transformative during the periods in which dependence forms. There need be no initial mistake for lock-in to occur. Joining can be rational at the time of entry, continued participation can remain rational during subsequent deterioration, and the resulting population can nevertheless arrive collectively at a state that none of its members would have chosen from the beginning. This is a path-dependent process rather than a single bad bargain. At each local transition the cost of remaining may be smaller than the cost of leaving, even while the sequence of such transitions moves participants from a highly favourable initial state to an extractive terminal state. Evaluating only the endpoint obscures how the endpoint became reachable; evaluating only each local decision obscures how individually rational continuations can accumulate into collective dependence. Enshittification therefore requires a trajectory-first analysis. The question is not simply whether users consented to each change, but how previous changes altered the feasible alternatives available when later consent was obtained.

The same trajectory explains why competition cannot be inferred from the mere presence of nominal competitors. A technically available substitute does not constrain a platform if moving to it requires reconstructing the very network effects, reputational capital, accumulated data, or audience relationships that make the incumbent useful. The relevant competitive variable is not the number of alternatives that exist but the cost of reaching them from the participant's present state. This is a stronger criterion than market variety because it treats migration as part of the object being evaluated. Ten isolated social networks do not necessarily provide meaningful competition to an incumbent containing the relationships one actually needs. A new marketplace does not immediately discipline an incumbent if seller reputation and buyer histories cannot move with the seller. A new creative platform does not provide an equivalent opportunity if the creator must rebuild an audience from zero while continuing to produce for the incumbent in order to survive. What matters is not whether another state can be described but whether there exists an admissible path from here to there.

This suggests a first generalization of Doctorow's theory. Enshittification can be understood as a process in which an intermediary progressively increases the distance between dissatisfaction and effective refusal. Early in the trajectory, a small deterioration may produce departure because alternatives are nearby. Later, a much larger deterioration can be absorbed because accumulated dependencies have moved those alternatives farther away. The platform thereby acquires what might be called an \emph{extraction interval}: the difference between the quality participants currently receive and the lowest quality they will tolerate before paying the cost of exit. Increasing switching costs widens that interval. Extraction consists partly in discovering how much of it can be consumed without triggering migration.

Once formulated this way, the central phenomenon is no longer adequately captured by the image of a product becoming worse. The deeper change is occurring around the product, in the space of possible responses to its deterioration. A platform becomes powerful when it can modify not only what happens within the interaction but the consequences of refusing the interaction, the portability of what has already been accumulated, the visibility of competing possibilities, and the informational conditions under which participants estimate whether departure is worthwhile. The result is a contraction of the possible that can coexist with an apparent proliferation of options. There may be more buttons, more content, more creators, more products, more recommendations, and more nominal services than ever before while the number of practically reachable trajectories decreases. Abundance inside the enclosure is compatible with scarcity at its boundaries.

The Google case is therefore an appropriate beginning because it makes the inversion unmistakable. A worse search engine is not necessarily evidence that optimization has failed. Under sufficiently concentrated conditions, it may be evidence that the object being optimized has changed. Once that possibility is admitted, the analytical problem shifts. We can no longer infer the purpose of a system from the activity its interface appears to facilitate, nor can we infer the welfare of its participants from the intensity with which they continue to use it. We must instead ask what objective is actually being optimized, what constraints prevent participants from rejecting that objective, what information each party possesses about the others, and which alternatives remain reachable when the resulting arrangement becomes intolerable. Those questions lead from enshittification as a theory of deteriorating platforms toward the more general problem considered in the next section: how a system can preserve formal choice while progressively contracting the counterfactual space in which choice has practical meaning.

\section{Enshittification as Counterfactual Contraction}

A counterfactual acquires practical force only when it can discipline the actual state. It is easy to imagine alternatives that possess no such force: a worker could theoretically abandon a profession, a musician could theoretically withdraw from digital distribution, a family could theoretically persuade every relative and friend to migrate simultaneously to another communications system, and a small business could theoretically reconstruct its customer relationships on an independent infrastructure. The mere grammatical availability of ``could,'' however, establishes almost nothing about the conditions under which these possibilities can be realized. An alternative whose realization requires sacrificing the goods that make the current arrangement valuable may remain logically coherent while becoming institutionally inert. It can still be invoked rhetorically---``nobody is forcing you to use the platform''---while no longer imposing meaningful discipline upon the platform's conduct. This is the point at which the distinction between possibility and reachability becomes indispensable.

Doctorow's account of social-media lock-in illustrates the problem particularly clearly because the source of dependence is distributed among the users themselves. The platform need not physically prevent departure. It need only mediate relationships whose value depends upon coordinated participation. One remains because one's friends remain; one's friends remain because their families, communities, customers, audiences, schools, clubs, or professional contacts remain; those people remain because still other networks remain. The resulting collective-action problem cannot be reduced to the sum of independent consumer preferences because each person's payoff depends upon the choices of others. An isolated decision to leave therefore destroys part of the value that motivated participation, while a coordinated migration would require agreement among populations whose reasons for using the platform are heterogeneous. The platform's strongest asset is consequently not the software in isolation. It is the accumulated difficulty of coordinating an alternative.

The same structure appears in less obvious forms wherever an intermediary succeeds in making accumulated value nonportable. A creator's audience can be numerically enormous while remaining institutionally inaccessible except through the platform that assembled it. A driver's history of reliable work can be economically valuable while ceasing to exist as reputation when the driver enters another marketplace. A seller's positive reviews can lower transaction costs for thousands of prospective buyers while remaining attached to an account that cannot carry those reviews elsewhere. The participant has produced part of the asset, often over many years, but the intermediary controls the conditions under which the asset remains usable. The resulting dependence is especially powerful because what is threatened by exit is not merely something the platform supplied. It includes value created by the participant and by other participants through their own prior activity. Lock-in thus permits the intermediary to appropriate bargaining power from accumulated cooperation without necessarily owning the underlying relationships in any ordinary social sense.

This is why the rhetoric of revealed preference becomes unreliable in platform environments. Continued participation is observable, but the counterfactual conditions necessary to interpret that observation are not. If a person remains on a deteriorating service because leaving would strand an audience built over a decade, the observation ``this person continues to use the service'' cannot discriminate between satisfaction with the service and unwillingness to surrender the audience. If a worker accepts a low-paying task because declining it risks immediate financial hardship, acceptance does not establish that the wage reflects the worker's unconstrained valuation of the task. If a listener tolerates unwanted recommendations because the desired catalogue, playlists, social features, devices, and habits have been consolidated within one service, continued listening does not establish approval of the recommendation system. Behaviour reveals a choice from the available set; it does not by itself reveal why that set has its present shape, what alternatives were removed before the choice occurred, or what costs have been attached to refusing the options that remain.

The methodological error is subtle because the observation itself may be perfectly accurate. People really do remain. Workers really do accept offers. Creators really do alter their work. Users really do click recommendations. The mistake occurs when an observed action is treated as though it were generated by preference alone rather than by preference operating within a constrained and historically produced field of possibilities. Power enters the analysis not by making behaviour unreal but by helping determine which behaviours are feasible. A theory that observes only the selected action therefore risks attributing to desire what should partly be attributed to constraint. In the limiting case, coercion itself can be redescribed as preference whenever the coerced person selects the least damaging option still available.

Counterfactual discipline requires more. To say that a platform's users prefer remaining to leaving is informative only relative to the conditions under which leaving occurs. Would they remain if their social graph were portable? Would creators remain under the same terms if they could carry their audience relationships elsewhere? Would sellers accept the same fees if reputation were interoperable? Would drivers accept the same compensation if competing labour platforms could observe their work histories but not exploit proprietary behavioural profiles of their desperation? Would listeners accept the same substitutions if an independent client could recover precisely the recordings they requested? These questions do not assume that the answer must favour intervention. They identify the missing comparisons required before observed behaviour can bear the interpretive weight placed upon it. A choice cannot establish the desirability of the environment that structured the choice when plausible modifications to that environment would change the alternatives themselves.

Enshittification can therefore be redescribed as a progressive weakening of counterfactual discipline. During the initial phase, users can respond to deterioration by moving, modifying, substituting, refusing, or coordinating around alternatives. The platform must consequently surrender some potential extraction because attempting to capture it would provoke a response. As dependencies accumulate and alternative paths become more costly, the range of profitable deterioration expands. Eventually, the institution may be able to change prices, increase advertising, alter rankings, suppress interoperability, reduce compensation, insert unwanted material, or modify access without losing enough participation to make the intervention unprofitable. What has contracted is not necessarily the number of imaginable alternatives. It is their capacity to act upon the present.

This distinction between an alternative's formal existence and its causal efficacy will recur throughout the argument. A possible world matters politically when some path toward it remains open enough that institutions in the actual world must take its accessibility into account. Competition disciplines firms because customers might leave; labour organization disciplines employers because workers might coordinate; interoperability disciplines platforms because complementary services might route around undesirable design choices; modification disciplines manufacturers because users or competitors might remove profitable restrictions; democratic institutions discipline officeholders because organized publics might replace them. In every case, the alternative need not actually occur in order to exert force. Its credible reachability changes behaviour in advance. The counterfactual acts upon the actual precisely because the actors controlling the actual know that another trajectory remains available.

The contraction of the possible occurs when those credible trajectories are converted into merely hypothetical ones. A platform need not abolish exit if it can make exit ruinous. It need not prohibit competition if it can ensure that competitors begin without the accumulated network required to become viable. It need not forbid creators from working elsewhere if it can make their existing audience nonportable. It need not explicitly dictate artistic form if its ranking mechanisms make other forms difficult to encounter. It need not command workers to accept low wages if it can infer which workers cannot afford refusal. The characteristic operation is therefore often not prohibition but \emph{distance production}: alternatives remain visible at the edge of the system while the cost of reaching them increases.

Understood this way, the contraction of the possible is compatible with enormous apparent abundance. Indeed, abundance can conceal it. A streaming platform can contain tens of millions of recordings while concentrating control over the pathways by which listeners encounter them. A social platform can contain billions of posts while narrowing the fraction of them that any particular person can realistically see. A labour platform can display a constant stream of nominal opportunities while using individualized information to determine which opportunities and wages become available to each worker. A creator economy can contain millions of participants while concentrating remunerated attention among a small and strategically visible fraction. The relevant scarcity is not necessarily scarcity of objects. It is scarcity of \emph{routes}: scarcity of independent means by which objects, people, audiences, opportunities, and institutions can encounter one another without passing through the same controlling intermediary.

This is where enshittification intersects with artificial scarcity. Digital objects are often inexpensive to reproduce, and networks can in principle permit extraordinary multiplicity of distribution. Yet abundance at the level of reproduction does not imply abundance at the level of access, visibility, bargaining power, or coordination. An intermediary positioned at a chokepoint can manufacture scarcity by controlling traversal rather than supply. The song exists, but access to listeners is ranked. The creator exists, but access to the creator is mediated. The job exists, but its wage is individualized. The friend exists, but the feed decides whether the friend's post appears. The alternative client can be written, but anti-circumvention rules may make its operation legally hazardous. The scarce object is increasingly not information itself but permission to establish a relationship across an institutional boundary.

The resulting form of power is easy to underestimate because it often leaves the underlying entities untouched. Nothing has to be destroyed. The friend remains a friend, the musician retains the recording, the driver retains the capacity to drive, the user retains the nominal right to leave, and the programmer retains the technical capacity to write alternative software. What changes is the connective tissue among them. The institution controls routing, ranking, compatibility, reputation, discoverability, payment, or legal admissibility, and thereby determines which potential relations become actual. A system organized around such chokepoints can contract possibility without visibly reducing the inventory of things from which possibilities are supposedly composed.

The appropriate unit of analysis is therefore not the isolated object or isolated choice but the reachable relation. This is also why the political response cannot be exhausted by demanding better outputs from the existing intermediary. A platform can improve recommendations tomorrow while retaining the capacity to worsen them the day after; it can increase creator compensation while retaining unilateral authority to reduce it later; it can voluntarily expose an interface while preserving the legal or technical power to withdraw it. Benevolent exercise of concentrated power does not remove the underlying dependency. If the objective is to prevent enshittification rather than temporarily reverse one of its manifestations, the relevant question is whether participants possess alternatives whose continued availability does not depend upon the permission of the institution those alternatives are meant to discipline.

The distinction becomes especially important when considering Doctorow's emphasis on interoperability and user modification. Their significance is not merely that they allow people to customize products according to taste. They alter the geometry of departure and therefore the bargaining relationship even for people who never exercise them. A compatible third-party client, a portable social graph, an independent recommendation layer, a transferable reputation, or a generic replacement component can reduce the distance between dissatisfaction and response. Once such routes exist, the platform must anticipate that excessive degradation may be routed around. The credible possibility of modification thereby constrains behaviour before modification occurs. Conversely, laws and technical measures that prohibit circumvention do more than prevent particular hacks. They remove possible continuations from the strategic environment and thereby enlarge the range of conduct the incumbent can impose without provoking an effective response.

The contraction of the possible should therefore not be confused with determinism. The argument is not that people trapped within platform systems possess no agency, nor that every local decision has been dictated in advance. On the contrary, the importance of reachability arises precisely because agency operates through constrained trajectories rather than outside constraint altogether. People improvise, coordinate, refuse, migrate, build alternative tools, organize collectively, and sometimes transform the institutions that constrained them. But the cost and likelihood of those actions depend upon the structure inherited from previous states. Freedom is not well represented by counting the number of abstractly describable choices at a single instant. It is better approached through the continuations that remain available after a choice is made, the resources preserved along those continuations, and the extent to which taking one path forecloses others.

This trajectory-sensitive conception also clarifies why apparently small acts of enclosure can have disproportionately large later consequences. Preventing data export, closing an interface, making a reputation score proprietary, restricting alternative clients, weakening labour organization, or acquiring a nascent competitor may each appear modest when evaluated at the moment it occurs. Yet each can remove edges from the graph of future transitions. The immediate welfare effect may be small while the long-run effect on reachability is substantial. By the time extraction becomes obvious, the routes that would once have constrained it may already have disappeared. Diagnosis at the endpoint then encounters a system whose historical alternatives cannot simply be recreated by reversing the most recent policy change.

The central claim can now be stated more precisely. Enshittification is a trajectory in which accumulated dependence progressively converts alternatives from reachable counterfactuals into merely formal possibilities, thereby weakening the external disciplines that once aligned institutional success with participant welfare. The platform's power grows not only because it possesses more users, data, capital, or market share, but because those accumulations change the cost of traversing the surrounding possibility space. Once this process is recognized, a second question immediately follows. If participants cannot directly observe the mechanisms that determine which paths are rewarded, punished, displayed, or withheld, how do they learn to act within the contracted space? It is here that platform power becomes epistemic as well as economic, and where Doctorow's discussion of algorithmic wage discrimination reveals one of the most disturbing consequences of the entire structure: a system can make its subjects increasingly legible to itself while making itself increasingly illegible to them.

\section{Algorithmic Superstition and the Asymmetry of Legibility}

A person acting within an opaque platform does not cease trying to understand the system merely because the system refuses to explain itself. On the contrary, opacity can intensify the demand for explanation. A driver notices that some trips pay better than others, a musician observes that one release enters a recommendation stream while another disappears, a comedian discovers that one form of video reaches hundreds of thousands of viewers while a formally similar performance reaches almost nobody, and a seller watches the position of a listing rise and fall without knowing which variables produced the change. Each participant possesses a sequence of observations and must decide what to do next. Because continued participation may depend upon making that decision successfully, the participant begins constructing an informal causal model: perhaps posting at a particular hour matters, perhaps accepting every available task improves one's standing, perhaps a particular duration is rewarded, perhaps the first few seconds of a video determine distribution, perhaps a certain vocabulary triggers suppression, perhaps consistency is rewarded, perhaps refusing undesirable work damages future opportunities. Some of these hypotheses may be locally accurate. Others may be false. The deeper problem is that the participant usually lacks the information required to distinguish between them, while the platform possesses both a much larger observational field and the capacity to alter the mechanism being inferred.

Doctorow describes the resulting condition as superstition \citep{doctorowfox2026}. The term is stronger than ordinary uncertainty because superstition does not consist merely in failing to know why an event occurred. It consists in constructing a causal relation from repeated but inadequately controlled observations and then reorganizing behaviour around that inferred relation. The familiar lucky hat is harmless because the sporting event does not generally observe the hat and modify its rules in response. Platform superstition is more consequential because the hidden mechanism may be adaptive, individualized, experimental, and strategically responsive. Doctorow's discussion of algorithmic wage discrimination in gig work \citep{dubal2023algorithmic,calo2015digitalmarket} makes the difference especially clear. If a platform infers a worker's economic desperation from the worker's willingness to accept undesirable jobs, then the worker's attempt to learn the system becomes one of the variables through which the system learns the worker. Accepting a low-paying task is not merely an observation of the wage function. It can supply evidence used to modify subsequent offers. The subject and object of inference are therefore coupled: the worker updates a model of the platform while the platform simultaneously updates a model of the worker.

This coupling creates an epistemic environment fundamentally different from one in which a stable but unknown rule must merely be discovered. Suppose a driver notices that accepting almost every offered trip coincides with declining average compensation. Several causal explanations are compatible with the observation. The available trips may simply have become worse; demand may have changed; location or time may explain the difference; the worker may have selected an unusual sample; the platform may reward selectivity; the platform may infer willingness to accept lower compensation from previous acceptance; or the relationship may result from some combination of these mechanisms. Repetition does not necessarily resolve the ambiguity because the act of repeating the behaviour may change the system's estimate of the worker. The experiment contaminates its own control condition. What appears to the participant as a problem of discovering the correct strategy may therefore be a problem for which no stationary strategy exists.

Creative workers face a structurally similar problem when attempting to ``learn the algorithm.'' The phrase itself implies that somewhere behind the interface there exists a discoverable mapping between creative behaviour and distribution, and that sufficient experimentation will reveal the rules of successful participation. Yet Doctorow's discussion of TikTok's internal ``heating'' mechanism illustrates why this assumption can fail. A platform capable of selectively amplifying particular creators can manufacture observations that those creators, their competitors, and their audiences will naturally interpret as evidence about underlying demand. A creator whose work is artificially elevated may conclude that a particular format, subject, style, or strategy caused the success. Other creators observe the apparent success and imitate it. The initial creator may then become an evangelist for the platform precisely because the creator mistakes discretionary amplification for evidence of an independently discovered audience. When the amplification moves elsewhere, the resulting decline invites another round of causal speculation. The platform can therefore produce not merely outcomes but beliefs about the causes of outcomes.

The giant teddy bear appears here in epistemic form. The conspicuous winner does not merely attract additional participants by demonstrating that a large reward is possible. The winner becomes an apparent experiment from which others infer how the game works. If a particular comedian, musician, commentator, illustrator, or influencer becomes suddenly ubiquitous, observers naturally search for differentiating features that might explain the result. They examine posting frequency, tone, format, subject matter, editing, presentation, authenticity, aggressiveness, duration, timing, or willingness to follow trends. The resulting explanations can become a substantial secondary industry of tutorials, consultants, courses, optimization services, and creator folklore. Yet if distribution includes hidden discretionary interventions, individualized ranking, changing objectives, or experiments conducted by the platform itself, the visible success cannot support the causal inference being drawn from it. The winner's behaviour is observable; the selection mechanism is not.

This is an especially severe instance of a more general problem of causal identification. To infer that an intervention \(I\) produced an outcome \(O\), one requires some warranted comparison with what would have happened in the relevant absence of \(I\). Merely observing \(I\) followed by \(O\) is insufficient when multiple variables changed, selection into \(I\) was nonrandom, the environment responded to \(I\), or the mechanism assigning exposure was itself correlated with expected \(O\). Platform environments can contain all of these difficulties simultaneously. A creator changes format and receives more distribution, but the platform may already have selected that creator for additional distribution. A driver becomes selective and receives better offers, but selection may correlate with location, schedule, experience, or an inferred willingness to wait. A seller purchases advertising and experiences increased sales, but the platform controls both the visibility purchased through advertising and the baseline visibility against which the increase is measured. The institution is not merely another variable inside the experiment. It often controls the assignment mechanism through which the experiment becomes observable.

The asymmetry becomes still sharper when platforms possess information unavailable to the participants whose behaviour they govern. Doctorow's examples extend beyond creative labour and gig driving to contract nursing, where platforms may use information about workers' financial circumstances to estimate how little compensation particular workers are likely to accept. The principle is more general than any specific implementation. A platform situated between many buyers and sellers, employers and workers, advertisers and audiences, or creators and viewers can observe interactions at a population scale unavailable to any individual participant. It can compare responses across prices, rankings, messages, delays, recommendations, and interface designs. It can estimate elasticities and reservation thresholds from behavioural traces. It can conduct experiments whose subjects do not know they are participating in an experiment. The participant, by contrast, typically observes only a narrow local history and receives little information about the population against which that history is being evaluated.

The platform consequently approaches an ideal of unilateral legibility. It seeks to make the participant increasingly predictable while preserving its own unpredictability. The worker's acceptance threshold should be inferable; the wage-setting rule need not be disclosed. The creator's responsiveness to changes in distribution should be measurable; the distribution rule need not be stable. The listener's probability of skipping a substituted recording should be estimated; the economic reason for the substitution need not be visible. The advertiser's willingness to continue purchasing impressions should be known; the actual causal effect of those impressions can remain difficult to audit. The user becomes an object of measurement while the measuring institution remains partially withdrawn from reciprocal measurement.

This asymmetry should be distinguished from ordinary privacy concerns. Privacy asks, among other things, whether an institution should possess particular information about a person. Legibility asks a broader structural question: what can each participant infer about the other, and how does the asymmetry in inference alter their respective capacities to act? A worker may voluntarily disclose every individual datum used by a platform and still occupy a radically asymmetric relationship if the platform can aggregate those data across millions of interactions while refusing to disclose the rule by which they affect compensation. Conversely, preventing collection of a particular datum may do little to restore bargaining equality if the same latent variable can be inferred from behaviour. Desperation need not be entered into a form labelled \emph{desperation}; it can be estimated from response times, acceptance rates, willingness to travel, working hours, frequency of refusal, account balances where accessible, or correlated commercial data. The relevant object is therefore not merely data possession but inferential capacity.

This is where the old social virtue of \emph{sprezzatura} \citep{castiglione1528courtier} acquires an unexpected economic relevance. In its traditional setting, sprezzatura names the cultivated appearance of effortlessness, a refusal to expose completely the labour, preparation, or technique underlying accomplished performance. Transposed cautiously into a platform environment, it points toward something more general: the strategic value of preserving capacities that have not yet been rendered fully legible to the evaluator. Bargaining has always involved incomplete disclosure. A buyer does not normally begin negotiation by announcing the highest price they would pay, nor a seller by announcing the lowest price they would accept. Each retains some private information because the undisclosed interval is part of the space in which negotiation remains possible. A platform capable of estimating one party's reservation threshold while concealing its own transforms that ordinary asymmetry into a systematic advantage.

The significance of ``not showing all one's cards'' is therefore not that ignorance is intrinsically liberating or that workers should cultivate mystique as a substitute for institutional reform. Individual opacity is fragile against organizations possessing enormous datasets and sophisticated inference systems. The important point is structural: a bargaining relationship changes when one participant can estimate the other's minimum acceptable outcome with substantially greater accuracy than the other can estimate its maximum obtainable outcome. If a labour platform knows that one worker will accept \$18 for a task while another will refuse anything below \$27, a uniform market price becomes unnecessary. The platform can approach each worker's reservation threshold individually, capturing a portion of the surplus that would otherwise remain with workers who happened to value the task differently. Personalized pricing and personalized wages are mirror images of the same possibility: infer the maximum extraction compatible with continued participation.

The contraction of the possible therefore occurs not only by removing alternatives but by learning where the boundary of refusal lies for each participant. This produces a more dynamic form of enshittification than a uniform decline in service. Instead of lowering quality until an average user threatens to leave, the institution can attempt to discover how much deterioration each user will tolerate. Instead of lowering wages across an entire labour market, it can estimate which workers can withstand refusal and which cannot. Instead of exposing every creator to the same distribution regime, it can determine how much visibility is required to keep each creator producing. The extraction interval described earlier becomes individualized. A platform with sufficient information need not ask how badly it can treat its population. It can ask how badly it can treat \emph{this person, at this moment, given the alternatives presently available to them}.

Such a system also complicates conventional ideas about discrimination. Traditional discriminatory practices often classify people into relatively stable groups and impose different terms upon those groups. Algorithmic discrimination can operate at much finer resolution. The relevant category may contain a single person in a transient state: this worker, with this recent history, at this location, carrying this inferred level of financial pressure, at this hour. The discriminatory rule can change before an external observer accumulates enough cases to characterize it. The resulting system is difficult to contest not only because its operation is secret but because the object one wishes to contest may not possess a stable public form. There may be no posted wage schedule to challenge, no explicit rule stating that indebted workers receive less, and no durable category corresponding to the affected population. The discrimination resides in the function that maps an evolving representation of the individual onto an offer.

Opacity under these conditions does more than conceal exploitation after it occurs. It interferes with collective knowledge formation. Workers who receive different offers may interpret differences as evidence of individual competence. Creators who receive different levels of distribution may attribute those differences to talent, discipline, authenticity, or mastery of the platform. Sellers who occupy different positions in rankings may assume that competitors have discovered superior tactics. Because each participant observes a personalized slice of the system, common conditions can be experienced as private successes and failures. The platform's informational advantage thereby produces a social consequence: people who might otherwise recognize themselves as occupying a shared structural position instead encounter one another as competitors possessing mysterious differences in skill.

Doctorow's Uber example captures this fragmentation particularly well. A worker who is selective about accepting jobs may earn more and conclude that they are ``good at Uber,'' while another who accepts nearly everything may earn less and conclude that they are ``bad at Uber.'' The interpretation converts an institutional wage-setting mechanism into a personal characteristic. Once this occurs, the worker who benefits from the mechanism can become evidence for the fairness of the mechanism, while the worker harmed by it can internalize the outcome as incompetence. Structural asymmetry has been translated into merit. The platform does not need to announce this explanation. Participants can generate it themselves because individualized outcomes arrive without sufficient information about the process producing them.

Algorithmic superstition therefore prepares the ground for meritocratic attribution. When the causal mechanism is inaccessible, visible differences among participants become unusually attractive explanations for differences in outcome. The successful creator appears more disciplined, the well-paid driver more strategic, the highly ranked seller more entrepreneurial, the viral performer more authentic. Some of these differences may be real and causally relevant. The error lies in inferring their causal importance from a selection process whose assignment mechanism remains hidden. The less that participants know about the platform's interventions, the easier it becomes to explain the resulting distribution through the characteristics of those who survived it.

\section{The Giant Teddy Bear and the Meritocracy of Survivors}

Doctorow's ``giant teddy bear gambit'' is useful because it captures a mechanism that ordinary discussions of survivorship bias often leave incomplete. Survivorship bias describes an epistemic error: successful cases are disproportionately visible, and observers consequently infer properties of the underlying process from a sample conditioned upon survival. The carnival operator adds an institutional actor who benefits from producing precisely that error. The giant teddy bear carried through the fairground is not simply a naturally occurring outlier that observers happen to misunderstand. It is an advertisement. Its visibility changes the behaviour of prospective players by supplying an emotionally vivid demonstration that the improbable reward belongs to the space of possible outcomes. The relevant deception need not consist in falsely claiming that nobody can win. It can operate through the much simpler substitution of one proposition for another: because winning is possible, playing must be worthwhile.

The distinction between possibility and expected outcome is elementary but extraordinarily easy to suppress when the distribution contains spectacular extremes. A platform containing millions or billions of participants will almost inevitably produce exceptional cases even when the typical participant receives little benefit. Indeed, the scale of the participant population makes extreme observations more likely. If enough musicians release recordings, enough comedians post clips, enough writers publish newsletters, enough streamers broadcast themselves, and enough aspiring entrepreneurs attempt to build audiences, some individuals will experience trajectories so unusual that they appear to demand individual explanation. Their biographies can then be searched for the hidden variable that supposedly caused the result: unusual discipline, relentless consistency, willingness to take risks, authentic self-expression, superior networking, rapid adaptation, an exceptional morning routine, or an intuitive understanding of what ``the algorithm wants.'' The fact that millions of less successful participants may exhibit the same characteristics is harder to perceive because failure does not generate an equally visible archive.

The statistical problem can be stated without denying the existence of talent. Suppose some characteristic \(T\) genuinely increases the probability of success. Observing that successful people disproportionately possess \(T\) does not establish that \(T\) explains the magnitude of their success, that possession of \(T\) makes participation advantageous, or that people lacking the observed outcome generally lacked \(T\). Those claims require information about the distribution outside the selected sample. If a million highly disciplined creators compete for a thousand positions of unusual visibility, the thousand survivors may all be disciplined without discipline being sufficient to explain why those particular thousand survived. The platform's allocation mechanism, network effects, timing, accumulated advantage, discretionary promotion, random variation, and numerous unobserved variables may determine the residual selection. Retrospective biography tends to erase that residual by constructing a coherent causal narrative around the cases that remain visible.

The giant teddy bear gambit adds another layer because the platform has reason to amplify those narratives. Exceptional payouts, enormous audiences, lucrative exclusivity agreements, and celebrity creators establish an aspirational upper tail whose symbolic importance can greatly exceed its cost. Paying one creator an extraordinary sum may be cheaper than improving the terms offered to millions of ordinary creators, particularly if the extraordinary payment persuades those millions that the platform is where extraordinary careers happen. The winner therefore performs productive labour even when merely being seen to have won. The giant teddy bear is compensation and recruitment expenditure at the same time.

This structure helps explain why highly unequal creator economies can sustain intense participation without offering attractive median outcomes. Participants are not necessarily irrational. Creative work often supplies intrinsic benefits independent of income, and the possibility of reaching an audience may itself possess value. But the institutional presentation of exceptional success can distort the counterfactual against which continued participation is evaluated. Instead of comparing the expected return from platform participation with other available uses of time and labour, participants are encouraged to compare their current state with the spectacular state occupied by a highly visible survivor. The relevant question quietly changes from ``Is this arrangement worthwhile relative to my alternatives?'' to ``What must I change about myself in order to become one of them?'' The platform disappears from the causal model while the participant becomes the variable under continuous revision.

This is where algorithmic superstition and meritocratic attribution reinforce one another. An opaque selection mechanism produces uneven outcomes; uneven outcomes generate attempts to infer successful strategies; survivors become the principal evidence from which those strategies are inferred; aspiring participants imitate the survivors; and the resulting imitation supplies the platform with a large population willing to adapt itself continuously to whatever behaviours appear to be rewarded. When the strategy stops working, failure can be attributed to incorrect execution rather than to a changing allocation mechanism. The lucky hat was worn incorrectly. The hook was insufficiently immediate, the posting schedule insufficiently consistent, the niche insufficiently specific, the personal brand insufficiently authentic, or the creator insufficiently willing to pivot. Because the system never supplies a stable causal specification against which these hypotheses can be falsified, optimization becomes potentially endless.

The result resembles a meritocracy while lacking one of the evidential conditions required to establish meritocratic selection. Merit requires more than differential outcomes correlated with desirable characteristics. It requires some warranted account of the selection process connecting those characteristics to the outcomes. An opaque mechanism capable of discretionary amplification, individualized ranking, endogenous experimentation, and continual rule changes cannot be inferred to reward merit merely because its winners possess traits we admire. Nor does the existence of genuine skill eliminate the problem. A rigged carnival game can still require skill; a labour market can reward competence while simultaneously exploiting desperation; an algorithm can favour high-quality work while also favouring strategically subsidized participants. The relevant question is not whether merit exists but how much explanatory weight it can bear once selection effects and institutional interventions have been accounted for.

This distinction matters because meritocratic narratives redistribute responsibility. If success demonstrates superior adaptation, failure suggests deficient adaptation. Structural conditions are translated into personal shortcomings, and interventions directed at the institution are displaced by interventions directed at the self. The creator buys another course, changes format, posts more frequently, studies analytics more obsessively, narrows the subject matter, imitates successful competitors, or produces additional work in the hope of supplying the missing variable. The worker accepts more jobs, works longer hours, or attempts to discover the hidden strategy separating high earners from low earners. The platform benefits from both responses because each increases participation while leaving the selection mechanism itself uncontested.

There is also a counterfactual error in the way exceptional success is commonly defended. Pointing to a billionaire musician, enormously successful podcaster, or globally famous creator answers the question ``Can someone become extraordinarily successful within this system?'' It does not answer ``Does this system distribute the value produced by creative workers fairly?'' or ``Are participants better off under this arrangement than under plausible alternatives?'' These are different hypotheses with different reference classes. The existence of an extreme success can coexist with declining compensation for the median creator, increasing concentration among intermediaries, or worsening bargaining conditions across the profession. Indeed, extreme winners can become more valuable to the system as typical outcomes deteriorate because their visibility supplies an increasingly necessary justification for continued participation.

The same error appears in philanthropic and meritocratic reasoning more broadly. An exceptional positive outcome is often evaluated against the weakest counterfactual available: without the benefactor, entrepreneur, platform, or patron, the recipient would have received nothing. Yet the relevant comparison may be an institutional arrangement in which the resources were distributed differently, bargaining power was less concentrated, or the intermediary never acquired unilateral discretion over allocation in the first place. Comparing a giant teddy bear with an empty hand makes the prize appear maximally generous. Comparing the entire carnival's intake with the distribution of prizes asks a different question. The actor who controls which comparison is salient therefore possesses evaluative power in addition to material power.

The creator economy makes this especially visible because spectacular success and widespread precarity are not contradictory observations. They can be complementary components of the same institutional arrangement. The spectacular success supplies aspiration, the precarity supplies competition, and the enormous population between them supplies experimentation. Millions of participants continuously test styles, formats, topics, schedules, identities, editing techniques, and modes of presentation at their own expense. The platform observes the results and can amplify whatever proves useful to its objectives. From the platform's perspective, this resembles an enormous distributed research-and-development system in which the experimental subjects finance much of the experimentation themselves.

The unsuccessful participant is therefore not necessarily economically useless to the platform. This point leads beyond the ordinary language of survivorship bias. If only winners mattered, the remainder of the population would constitute waste. But platforms can derive substantial value from the population that never achieves stable success. Their work expands the content inventory, increases the probability that something will capture attention, supplies data about audience response, disciplines established creators through the threat of substitution, generates advertising inventory, and continually replenishes the pool from which new winners can be selected. The relevant population is better understood not as a collection of failed aspirants but as a \emph{talent reserve}: a standing surplus of creative capacity whose abundance strengthens the intermediary even when most of that capacity is never remunerated at a level sufficient to sustain its producers.

Van Vogt's \emph{The World of Null-A} literalizes this arrangement in a way worth pausing on \citep{vanvogt1945worldofnulla}. The novel's Earth is governed in part by the Games Machine, a vast computing apparatus that periodically tests the population's semantic and cognitive integration and, on the basis of that testing, allocates the small number of emigration slots to Venus, the novel's explicitly utopian world of abundance and self-realization. Citizens are told, and mostly believe, that the Machine's selection is simply a measurement of readiness: those who pass have achieved the mental discipline the better world requires, and those who fail have not yet done so. The protagonist, Gilbert Gosseyn, eventually discovers that his own memories, and indeed his continued existence across several duplicate bodies, have been arranged by interests operating behind the Machine's ostensibly neutral test. The novel is thus doing in miniature what this essay has argued platforms do at population scale: a small number of successful passages through an opaque evaluative apparatus are held up as proof that the apparatus rewards genuine merit, while the apparatus itself, and the vastly larger population it screens out, recedes from the citizens' account of their own condition. Venus functions here exactly as the giant teddy bear functions on the platform: not primarily as a reward, but as a standing demonstration that the game is worth playing, sufficient to keep the tested population continually working to optimize itself against a test whose actual criteria, and actual administrators, remain hidden from the people it evaluates.

\section{Talent Reserve and the Production of Artificial Scarcity}

The familiar concept of brain drain describes a movement of skilled people away from a community, institution, or region that consequently loses access to their capacities. Talent reserve names the inverse structural possibility. A system may contain far more skill, creativity, knowledge, and willingness to work than it can or intends to remunerate, while nevertheless benefiting from keeping that surplus capacity available. The reserve is not valuable because every member will eventually be selected. Its value lies partly in the fact that selection remains scarce while eligibility remains abundant. A platform with millions of aspiring creators does not need to provide millions of sustainable creative careers in order to benefit from their presence. It needs enough participation to maintain a continuous supply of material, enough visible mobility to preserve aspiration, and enough scarcity of remunerated attention that participants remain willing to compete for access.

This reverses an intuitive account of scarcity. Creative talent itself may be abundant. Digital reproduction may be effectively abundant. The capacity to distribute files may be abundant. What remains scarce is the institutional recognition that converts creative activity into a sustainable livelihood: recommendation, placement, audience access, contractual opportunity, monetization, sponsorship, inclusion in prominent playlists, or other privileged routes through which attention becomes income. The platform can therefore operate simultaneously as an engine of abundance and an administrator of scarcity. It encourages enormous production while rationing the channels through which production becomes economically consequential.

Artificial scarcity need not mean that the scarce resource is wholly fictitious. Human attention is finite, prominent positions in a feed are finite, and listeners cannot hear every recording. The artificial component concerns the institutional organization of those constraints. A platform determines which scarcity becomes salient, how access to it is allocated, what information participants receive about that allocation, and whether alternative routes around the bottleneck are permitted to develop. The fact that attention is finite does not entail that a single intermediary must control its allocation. Nor does the fact that not every creator can become famous entail that creators must surrender bargaining power over the conditions under which their work reaches audiences. Natural limits can become the raw material from which institutional chokepoints are constructed.

The talent reserve strengthens those chokepoints because abundance on the supply side weakens the bargaining position of any individual supplier \citep{azarmarinescu2022monopsony,naidu2024monopsony}. A successful creator who demands better terms can be reminded, implicitly or explicitly, that innumerable others desire access to the same audience. Even where particular creators possess substantial leverage, the platform can cultivate replacements, diversify attention, alter recommendation weights, or subsidize emerging participants until dependence upon any one supplier decreases. The giant teddy bear gambit assists this process by ensuring that entry into the reserve continues. The spectacular survivor tells prospective participants that the queue is not merely a queue; it is a path to exceptional success.

This is one reason platform capitalism can extract value from aspirations it does not satisfy. Traditional accounts of employment tend to focus on completed exchanges: a worker supplies labour and receives compensation. Creator platforms often benefit from labour performed before any remunerative exchange occurs. The aspiring musician records and uploads music, the comedian produces clips, the writer posts essays, the illustrator supplies images, and the streamer broadcasts for hours while attempting to become sufficiently visible to qualify for meaningful compensation. The cost of discovering which participants can attract attention is thereby externalized onto the participants themselves. The platform receives a vast field of candidate production and pays substantially only after selection has occurred, if it pays at all.

The reserve also has an informational function. Because participants vary their behaviour while attempting to discover what succeeds, the platform receives a continuous stream of natural experiments. One creator changes format, another changes duration, another adopts a trend, another specializes, another becomes more provocative, another increases posting frequency. Audiences respond, and the intermediary observes those responses at scale. What appears from the creator's perspective as individual experimentation with artistic or professional strategy appears from the platform's perspective as population-level exploration of an optimization landscape. The platform can harvest successful innovations without bearing the full cost of generating them.

This arrangement creates a peculiar asymmetry between exploration and reward. The population explores broadly, but rewards can be concentrated narrowly. Millions of unsuccessful experiments are financed by creators through uncompensated time, equipment, training, foregone alternatives, and emotional investment. A comparatively small subset of successful experiments can then be amplified, imitated, standardized, and eventually incorporated into the platform's own recommendation or generation systems. The talent reserve therefore does not merely supply replaceable labour. It supplies the variation from which the system learns what replacement should look like.

The relationship to artificial intelligence is already visible. A platform that has spent years observing millions of creators adapting themselves to measurable audience responses possesses not simply a collection of cultural objects but an enormous record of selection. Which hooks retained attention, which visual compositions produced clicks, which phrases generated replies, which musical transitions prevented skipping, which comedic structures were replayed, which voices were abandoned, which thumbnails attracted viewers, and which forms of novelty remained close enough to established patterns to be tolerated can all become features of a machine-legible environment. Human creators perform the exploratory search through cultural possibility while the platform accumulates the behavioural traces required to model the resulting terrain.

The talent reserve is therefore more than a labour-market phenomenon. It is part of the representational infrastructure through which creative practice becomes computationally tractable. An abundant population of creators continuously translates difficult-to-formalize cultural judgments into observable behavioural outcomes. Audiences click, skip, linger, replay, subscribe, abandon, share, purchase, or ignore. These actions are imperfect proxies for aesthetic value, but they are measurable, and measurement allows the platform to substitute an operational objective for a concept it cannot directly represent. The system need not know what makes a joke funny, a song meaningful, an essay illuminating, or a painting beautiful. It needs only enough data to predict which measurable behaviours follow exposure to particular objects.

The danger is not merely that the proxy is crude. As argued earlier, optimization changes the meaning of the proxy itself. Once creators know that retention is rewarded, they alter work to maximize retention. Once particular durations are favoured, work migrates toward those durations. Once frequency affects visibility, production schedules accelerate. Once certain subjects receive disproportionate distribution, participants move toward those subjects. The measured population therefore ceases to be independent of the measurement regime. Cultural production bends toward the coordinate system through which it is evaluated.

The reserve population accelerates this bending because participants who refuse adaptation compete with enormous numbers who will attempt it. An established creator may dislike a new format, but aspiring creators have strong incentives to adopt it. Once enough participants do so, audience expectations can themselves change, giving the platform's initially arbitrary constraint the appearance of a natural property of the medium. A historically contingent ranking preference becomes a convention; the convention becomes a skill; the skill becomes a criterion of professionalism; and eventually the resulting form can be treated as what audiences ``want.'' The path by which the preference was produced disappears from the explanation.

This is a form of path dependence with unusually strong consequences for cultural memory. If one observes only the resulting equilibrium, it may appear that creators converged upon a format because that format efficiently satisfied audience preferences. Yet the same equilibrium could have emerged because the platform temporarily privileged the format, creators adapted to obtain distribution, audiences became accustomed to what was distributed, and subsequent engagement data then confirmed the preference that the original intervention helped create. The system has learned from a world partly produced by its own previous learning. Its data are endogenous to its interventions.

Talent reserve therefore connects artificial scarcity to representational narrowing. The platform need not prohibit alternative creative forms. It can simply make access to remunerated attention sufficiently scarce that creators voluntarily concentrate around the forms most likely to traverse the bottleneck. Formal expressive freedom remains enormous while economically sustainable expressive freedom contracts. This is precisely the kind of structure in which the vocabulary of choice can coexist with a shrinking space of continuation.

The next step in the argument concerns why so many people remain willing to enter and remain within this reserve despite poor expected returns. The giant teddy bear supplies one answer by making exceptional success visible. But it is not enough. People also continue because the activity itself matters to them. Creative labour therefore exposes an additional resource available for extraction, one that cannot be understood solely through wages, probabilities, or platform incentives. The system encounters people who would continue making art even if the market disappeared tomorrow and discovers that their devotion can be converted into bargaining weakness. Doctorow's invocation of vocational awe identifies this mechanism, but its implications extend further: when work is bound to identity, care, vocation, curiosity, or the intrinsic need to create, the threat of withholding the conditions under which that work can continue becomes more powerful precisely because the worker values something beyond compensation.

\section{Vocational Awe and the Capture of Meaning}

The talent reserve cannot be explained entirely by mistaken expectations about future success. If every participant entered creative work solely because they had miscalculated the probability of becoming wealthy or famous, then correcting those probabilities should substantially empty the reserve. Yet musicians continue making music after learning how little streaming pays, writers continue writing despite knowing that most books sell modestly, comedians continue performing before small audiences, scholars continue researching questions that offer little prospect of financial reward, and teachers, librarians, nurses, caregivers, and innumerable other workers remain attached to occupations whose social or personal importance exceeds the compensation attached to them. Doctorow invokes Fobazi Ettarh's concept of \emph{vocational awe} \citep{ettarh2018vocationalawe} to describe one consequence of this attachment: institutions can exploit the belief that certain forms of work are intrinsically virtuous, socially indispensable, or personally fulfilling in order to normalize conditions that would otherwise be recognized as unacceptable, a dynamic documented across gig and platform labour more broadly \citep{rosenblat2018uberland,gray2019ghostwork,wood2020goodgig}. The worker's devotion does not stand outside the economic relationship. Once it becomes predictable, it becomes one of the resources around which the relationship can be organized.

This is easiest to see by comparing two superficially similar statements: ``I would do this work even if nobody paid me'' and ``therefore nobody needs to pay me adequately for doing it.'' The first can express something admirable about the activity. The second does not follow from it. Yet markets organized around vocational labour repeatedly behave as though the transition were legitimate. A person's willingness to create in the absence of adequate compensation becomes evidence that compensation can safely be reduced; the intrinsic reward of the work is implicitly counted as part of the wage. What the worker experiences as meaning appears to the institution as elasticity. If musicians continue uploading despite declining returns, if academics continue researching despite precarious employment, or if teachers continue supplying unpaid labour because abandoning students would violate their own sense of responsibility, then devotion functions economically as a subsidy transferred from the worker to the institution.

The subsidy is especially powerful because it is generated internally by the person being exploited. A factory owner who wishes to reduce wages encounters a worker's desire not to perform labour for insufficient compensation. An institution governing vocational labour can encounter a different structure: the worker may actively desire to continue performing the labour even while rejecting the conditions under which it is performed. Refusal therefore becomes psychologically and materially complicated. To refuse the employer, publisher, platform, school, hospital, streaming service, or cultural intermediary may also mean refusing access to the activity, audience, community, patient, student, reader, or practice that gave the work meaning. The institution need not manufacture this conflict. It need only position itself between the worker and the valued continuation.

This distinguishes vocational lock-in from ordinary switching costs while connecting it to the same general theory of reachability. A creator can possess the technical capacity to leave a platform yet remain because leaving threatens the audience through which the creative activity has become socially meaningful. A teacher can possess the formal right to withhold unpaid labour yet know that doing so imposes immediate costs upon students rather than administrators. A nurse can refuse an additional burden while knowing that the burden does not disappear but is redistributed among colleagues or patients. In each case, the institution externalizes part of the cost of refusal onto something the worker values. The worker is not merely deciding whether the wage justifies the labour. The worker is deciding whether resisting the institution justifies harming, abandoning, or losing access to the object of care that motivated the labour in the first place.

Vocational awe can consequently be understood as a mechanism for altering the worker's effective exit cost. The more deeply an activity is bound to identity, responsibility, curiosity, community, or care, the more costly it can become to abandon the institutional channel through which that activity is presently realized. This does not imply that strong identification with one's work is pathological. The pathology lies in an institutional arrangement capable of capturing the difference between what the activity means to the worker and what the institution must pay to keep the worker performing it. Meaning produces surplus value when the institution can appropriate the worker's willingness to continue despite deteriorating terms.

Creative platforms possess an especially advantageous version of this structure because they often do not need to employ creators at all. The artist arrives already motivated, frequently supplies equipment and training, bears much of the cost of production, assumes the risk that the work will attract no audience, and may continue producing even when compensation is negligible. The platform's task is not to induce production from an otherwise unwilling workforce but to mediate the relationship between a pre-existing desire to create and a pre-existing desire to encounter creative work. This helps explain why monopsony is as important as monopoly in Doctorow's account. The platform's power does not consist only in controlling what consumers can buy. It consists in occupying a privileged position as purchaser, distributor, gatekeeper, or allocator of attention for people whose willingness to supply greatly exceeds their willingness to abandon the activity altogether.

The distinction between monopoly and monopsony is particularly important because popular discussion often treats consumer abundance as evidence against concentration. A listener may encounter more recordings than any person in history could possibly hear, a reader may have access to more writing than could be consumed in a lifetime, and a viewer may face effectively unlimited audiovisual material. None of this establishes that the producers of those materials possess meaningful bargaining power. A market can be extraordinarily abundant from the consumer's perspective while highly concentrated from the supplier's perspective. Indeed, the two conditions can reinforce one another. An enormous talent reserve supplies the consumer with apparent abundance while the intermediary's control over access to that abundance places creators in intense competition with one another. The consumer sees millions of choices; the creator sees a small number of gateways through which economically consequential attention can be reached.

This asymmetry clarifies why the familiar injunction to ``do what you love'' can become economically dangerous when detached from institutional analysis. The phrase treats intrinsic motivation as though it automatically aligned the interests of worker and employer. In reality, loving the activity can make the worker unusually vulnerable to whoever controls the conditions of its continuation. The aspiring musician may accept terms that someone indifferent to music would immediately reject because the alternative is not another equally meaningful occupation at a better wage. It is less music, less audience, or no sustainable musical practice at all. The institution can therefore extract not despite the worker's passion but through it.

The same mechanism helps explain why cultural labour so readily becomes individualized. Because the activity is associated with identity and self-expression, poor outcomes can be interpreted as judgments upon the self rather than as properties of the market. A failed upload is not experienced merely as an unsuccessful transaction; it can feel like evidence that the work was insufficiently compelling, the artist insufficiently disciplined, or the person insufficiently interesting. Algorithmic superstition intensifies this vulnerability by ensuring that the causal mechanism remains uncertain. The participant knows that something must be changed but does not know whether the relevant variable lies in the work, the presentation, the timing, the platform's experiment, the audience sample, or the economic objective governing distribution. The self becomes the most readily available object of intervention.

This produces a recursive relationship between vocational awe and meritocratic attribution. Because the worker cares deeply about the activity, the worker is willing to invest heavily in improvement. Because the selection mechanism is opaque, improvement can expand indefinitely into optimization against inferred platform preferences. Because visible survivors are interpreted as evidence about what the platform rewards, their behaviour supplies templates for self-modification. Because the platform benefits from continued experimentation, it has little incentive to make the causal structure fully transparent. Devotion supplies the energy, opacity supplies the uncertainty, and survivorship supplies the direction in which participants are encouraged to move.

The resulting system can extract extraordinary amounts of uncompensated cognitive and emotional labour. A creator does not merely produce the work. The creator studies analytics, tests thumbnails, rewrites titles, changes posting schedules, monitors trends, learns new formats, responds to comments, maintains multiple channels, develops a personal brand, anticipates moderation, manages audience expectations, and continually interprets fluctuations in distribution. Much of this labour exists because the platform has made access to audiences contingent upon successful navigation of an environment whose rules remain partially hidden. The creator is therefore required not merely to create but to solve an ongoing inverse problem: infer from observable outcomes the hidden function determining whether creation will be seen.

Vocational awe makes this additional labour unusually easy to normalize. Because the creator's occupation is represented as a privilege---the opportunity to make art, entertain an audience, teach, write, research, or otherwise pursue meaningful work---administrative burdens that would be recognized as uncompensated labour elsewhere can be redescribed as simply part of ``building a career.'' The aspirant who objects can be reminded, implicitly, that thousands of others would gladly occupy the same position. Talent reserve and vocational awe thereby reinforce one another. The reserve weakens bargaining power because replacements are abundant; vocational attachment replenishes the reserve because people continue entering despite poor bargaining conditions.

There is a further asymmetry here that becomes important for the later discussion of artificial intelligence. Human beings can value an activity independently of its measured output. A musician may value experimenting with a piece nobody hears; a writer may value pursuing a difficult argument that attracts little engagement; a teacher may value a conversation whose effect cannot be quantified; a researcher may spend years on a question that produces no immediately useful result. From the perspective of a platform organized around measurable interaction, these motivations are simultaneously productive and inconvenient. They generate the cultural variation on which the system depends, but they also produce activities whose value exceeds or escapes the available metrics. The platform therefore has reason to preserve the motivation while redirecting its expression toward forms that become legible within its evaluative apparatus.

This is the beginning of the manufacture of replaceability. The intrinsically motivated worker supplies a broad and heterogeneous space of possible activity. The platform does not initially need to automate that entire space. It can first narrow the economically viable portion of it by attaching distribution and compensation to measurable proxies. The creator adapts because creation matters enough to justify adaptation. Over time, behaviours that once represented external constraints upon the work become internalized as features of competent practice. The artist learns the preferred duration, the appropriate cadence, the effective hook, the recognizable category, the engagement ritual, the optimal release schedule, and the stylistic signals that survive recommendation. What began as accommodation to an intermediary can eventually appear to be a property of the medium itself.

Vocational awe therefore occupies a central rather than peripheral position in the argument. Without it, platform exploitation of creative labour can look like an ordinary story of workers accepting poor terms because alternatives are scarce. With it, we can see a more unusual mechanism: the platform captures value from the very fact that people refuse to treat their work as interchangeable with other work. Their attachment enlarges the amount of deterioration they will tolerate before abandoning the activity, and the institution positioned between them and the activity's continuation can appropriate part of that interval. The contraction of the possible is thus not only imposed externally through technical lock-in, market concentration, and legal restriction. It can operate through the structure of care itself, by making the path of refusal pass through something the participant is unwilling to sacrifice.

\section{From Human Expression to Machine-Legible Behaviour}

A historical contrast helps make the next transformation visible. In Donald G.\ Fink's 1966 popular introduction to artificial intelligence, \emph{Computers and the Human Mind}, computation is explained through a sequence of explicit representational transitions. Ordinary logical relationships expressed in language are rendered through class diagrams and truth tables; those formal relations are represented through logical operations; logical operations are realized in electrical circuits; and the persistence of computational state is traced into physical devices such as transistors, magnetic cores, and magnetic storage. What is striking from the perspective of contemporary computing is not simply the technological antiquity of the components. It is the pedagogical insistence that the transitions between levels remain visible. The reader is repeatedly shown how a distinction represented in one register can be preserved under translation into another.

The direction of explanation matters. Fink begins with distinctions already meaningful to a human reader and asks how a machine can physically realize them. Conjunction, disjunction, negation, membership, and exclusion are not introduced as mysterious properties emerging from an inscrutable computational substrate. They are progressively decomposed until the reader can see how an abstract logical relation becomes an arrangement of material states. The explanatory movement can be schematized as
\[
\text{human distinction}
\longrightarrow
\text{formal representation}
\longrightarrow
\text{logical operation}
\longrightarrow
\text{physical realization}.
\]
The arrows are not claims of identity. Each transition requires a representation preserving some structure while discarding or abstracting away other features. The success of the explanation depends upon making those choices sufficiently explicit that the reader can understand what has been preserved.

Contemporary platform mediation increasingly performs a transformation in the opposite direction. The starting point is not a cleanly specified logical distinction but the immense heterogeneity of human behaviour: conversation, music, humour, friendship, desire, curiosity, anger, teaching, argument, storytelling, performance, craft, and play. These activities contain distinctions whose relevance depends upon contexts that are difficult to formalize and often impossible to measure directly. The platform nevertheless requires operational variables upon which ranking, recommendation, advertising, compensation, moderation, and experimentation can act. It therefore constructs proxies. A friendship becomes a pattern of interactions; interest becomes dwell time; aesthetic attraction becomes replay probability; relevance becomes click-through; satisfaction becomes continued engagement; artistic success becomes streams, followers, retention, or conversion.

The resulting movement can be represented schematically as
\[
\text{human practice}
\longrightarrow
\text{platform representation}
\longrightarrow
\text{measurable behaviour}
\longrightarrow
\text{optimization target}.
\]
Again, none of these transitions is necessarily illegitimate. Measurement requires representation, and representation always abstracts. The problem begins when the representation ceases merely to describe the activity and becomes the principal mechanism through which access to the activity's rewards is allocated. Once visibility depends upon the metric, participants acquire incentives to alter themselves in order to alter the metric. The arrow then becomes recursive:
\[
\text{human practice}
\longrightarrow
\text{measurement}
\longrightarrow
\text{allocation}
\longrightarrow
\text{modified human practice}
\longrightarrow
\text{measurement}.
\]
The system no longer measures an independent world. It measures a world increasingly organized around being measurable by that system.

This recursive transformation provides a more precise interpretation of Jesse David Fox's observation, in conversation with Doctorow, that artists are being induced to behave like robots before robots are introduced as their replacements. The phrase is deliberately provocative, but the underlying structure is exact. Creative workers learn that certain forms are easier for platforms to classify, recommend, monetize, and compare. They consequently reorganize production around those forms. The comedian learns the duration that survives a feed, the musician learns the opening interval before a listener is likely to skip, the video creator learns the visual rhythm associated with retention, and the writer learns which titles or structures generate measurable response. These adaptations need not be imposed by explicit command. The allocation mechanism supplies the selective pressure.

The distinction between a constraint and a command is important. A command tells the participant what to do. A constraint changes which continuations remain viable. Platforms can therefore reshape cultural production without possessing a coherent artistic doctrine and without requiring any executive to decide what art should become. If one format receives more distribution than another, creators who depend upon distribution gradually move toward it. If frequent production is rewarded, production accelerates. If immediate engagement determines subsequent visibility, openings become increasingly optimized for immediate engagement. If easily classified material travels more reliably through recommendation systems, ambiguous or cross-categorical work acquires an additional distribution cost. The platform's representational limitations become environmental pressures acting upon the population it represents.

Over time, these pressures can produce an inversion of the explanatory programme visible in Fink. Instead of asking how computational mechanisms can be constructed to preserve distinctions humans already find meaningful, the institution increasingly asks which human distinctions can survive translation into the computational mechanisms already available. What cannot be measured becomes difficult to optimize; what cannot be optimized becomes difficult to justify within the system's objective function; what cannot be justified becomes less likely to receive distribution; and what receives little distribution becomes increasingly difficult to sustain professionally. The limitation of the representation thereby migrates outward until it appears as a limitation of the represented world.

This migration is especially consequential because machine legibility is not neutral with respect to form. Some behaviours generate cleaner signals than others. A click is easier to count than ambivalence, completion easier to detect than reflection, duration easier to measure than significance, repetition easier to identify than transformation, and immediate response easier to attribute than an influence that changes someone's thinking years later. A platform therefore possesses an intrinsic tendency to privilege effects that close quickly enough to become measurable within its observational window. Cultural objects whose value unfolds slowly, indirectly, collectively, or unpredictably are disadvantaged not necessarily because audiences value them less but because their effects are harder to assign to a discrete measurable event.

The temporal structure of measurement can thus become a temporal structure of production. If a system rewards effects visible within minutes or hours, creators learn to produce for minutes or hours. If recommendation depends heavily upon early response, early response becomes disproportionately important to the work's survival. A work requiring patience encounters not merely an impatient audience but an allocation mechanism that may prevent the patient audience from encountering it in sufficient numbers. What appears subsequently as evidence that ``people have shorter attention spans'' may therefore combine changes in audience behaviour with an institutional selection effect favouring works optimized for short measurement horizons.

The same problem applies to novelty. Genuine novelty is difficult to classify precisely because existing categories do not yet capture it. Recommendation systems trained upon prior behaviour naturally possess stronger evidence about objects resembling things that have already elicited measurable responses. This does not make novelty impossible, but it can impose a tax upon forms whose nearest neighbours are unclear. Creators seeking reliable distribution therefore face an incentive to remain novel enough to differentiate themselves while remaining familiar enough to be legible to the system. The result resembles an optimization around the boundary of an established category: variation is encouraged, but variation too distant from the existing representation becomes difficult to route.

Fink's class diagrams provide an unexpectedly useful conceptual contrast here. In elementary Boolean logic, the universe of discourse and the boundaries of the relevant sets are specified before the operation is performed. One can then ask whether an element belongs to \(A\), \(B\), \(A \cap B\), or the complement of some set. Platform classification introduces a prior political and epistemic question: who selected the universe of discourse, which distinctions were encoded as categories, and what happens to objects whose relevant properties do not align with those boundaries? A classification system can be perfectly consistent after its categories have been chosen while remaining badly suited to the phenomenon being classified. Formal correctness within the representation does not establish the admissibility of the representation itself.

This distinction is central to the broader argument. A platform can optimize accurately with respect to its chosen variables while degrading the activity those variables were intended to represent. The recommendation system may become increasingly accurate at predicting clicks while becoming less useful for discovery; the engagement model may become increasingly accurate at predicting continued interaction while producing experiences users later regard as wasted time; the creator-ranking system may become increasingly accurate at identifying material likely to retain viewers while narrowing the cultural forms that receive sustainable distribution. There is no contradiction. Accuracy is always accuracy relative to a target, and the target can become progressively detached from the value whose proxy originally justified it.

The danger becomes greater when the system's own outputs supply the data used to validate the target. If the platform promotes material predicted to generate engagement, then observes that promoted material generating engagement, it receives evidence partly produced by its own intervention. Creators adapt to the promoted forms, audiences encounter more of them, familiarity increases, and subsequent behavioural data can strengthen the model's confidence that these are the forms audiences prefer. A feedback loop emerges in which prediction, allocation, adaptation, and observation mutually reinforce one another. The platform can eventually appear merely to be reflecting demand even where it has participated substantially in constructing the demand it reflects.

This is an instance of intersubjectivity collapse, a term borrowed here from Pachniewski's account of the breakdown of mutual predictability among divergent minds and applied to a considerably narrower and more mundane setting: the mediated relation between creators and audiences rather than the relation between radically divergent cognitive architectures \citep{pachniewski2023intersubjectivity}. Creators and audiences do not encounter one another directly across the full range of what each might produce or value. They encounter platform-mediated representations of one another. The creator sees analytics, comments, rankings, follower counts, retention curves, and distribution outcomes and attempts to infer the audience behind them. The audience sees a feed, playlist, search ranking, or recommendation stream and attempts to infer the cultural world from what has been selected for presentation. The platform occupies the middle, possessing a broader view of both populations while controlling much of the channel through which each becomes visible to the other. The creator can consequently mistake the platform's representation of the audience for the audience itself, while the audience mistakes the platform's representation of available creation for the field of what creators are making.

Once this occurs, narrowing can become self-confirming. Creators produce more of what appears to reach audiences; audiences encounter more of what creators have learned to produce for reach; the platform records the resulting interaction as evidence of preference; and the next allocation is adjusted accordingly. Alternatives that never receive sufficient exposure generate little evidence of demand, while the absence of such evidence justifies withholding exposure. The system confuses the absence of an observed relation with evidence that the relation would not have formed.

The contraction of the possible therefore reaches beyond market concentration into the representational conditions of culture. A platform can leave every creator formally free to make anything while systematically altering which creations possess viable paths toward an audience. This distinction is analogous to the earlier difference between the existence of an alternative and its reachability. A cultural possibility can exist in a file on someone's computer, on an obscure website, or at the bottom of an enormous catalogue while remaining practically absent from the social field in which it could acquire meaning. Distribution is not an afterthought to cultural production. It is one of the processes through which cultural objects enter relations with people.

The historical comparison with Fink now becomes sharper. His explanatory project attempts to preserve intelligibility while moving downward through representational layers. Contemporary platform systems increasingly preserve operational control while allowing intelligibility to disappear. The user encounters an output without access to the transformations producing it; the creator encounters an allocation without access to the objective determining it; the worker encounters an offer without access to the inferred profile shaping it. At the same time, the human participant is required to become progressively more legible to the system. The representational ladder has become asymmetric: one direction is optimized for measurement, while the reverse direction is obstructed.

This asymmetry establishes the conditions under which claims of machine replacement must be evaluated. If human creative practice has already been filtered through years of selection for machine-legible characteristics, then the object presented to artificial intelligence as a replacement target is not an untouched sample of human creativity. It is a historically transformed subset of creative practice whose forms have partly adapted to the measurement infrastructure now declaring them automatable. The comparison between human and machine therefore contains a hidden trajectory. Before asking whether the machine can reproduce the observed output, one must ask how the observed human output came to occupy its present form.

\section{The Manufacture of Replaceability}

The claim that a machine can replace a human worker appears, at first, to require a straightforward comparison of capabilities. Identify the task performed by the person, determine whether the machine can perform the same task to an acceptable standard, compare costs, and substitute where the machine proves advantageous. This model is plausible when the task itself is stable and independently specified. It becomes much less reliable when the institution conducting the comparison has spent years redefining the task around what its computational infrastructure can observe and reward. In such cases, replaceability is not merely discovered at the end of technological development. It can be produced gradually by changing the human activity until the remaining economically recognized portion of that activity lies within the machine's competence.

This distinction matters because jobs and creative practices are not natural kinds with eternally fixed boundaries. Institutions decide which components of an activity count as economically relevant, which outcomes count as successful, and which dimensions can be ignored. A teacher's work can be represented as delivering standardized instructional content and producing measurable changes in assessment scores, or it can be represented as participation in a long-term social relationship involving judgment, care, encouragement, discipline, interpretation, and innumerable effects resistant to immediate measurement. A musician's work can be represented as supplying audio capable of occupying a playlist slot without causing listeners to leave, or it can be represented as participation in a cultural practice involving performance, community, historical reference, identity, experimentation, and shared attention. Automation appears easier under the first representation in each case, but that ease partly results from what the representation has excluded.

The central error is therefore a failure to distinguish \emph{task equivalence} from \emph{representational equivalence}. A machine may reproduce the variables through which an institution has chosen to evaluate an activity without reproducing the activity in the richer sense from which those variables were abstracted. If a platform evaluates music primarily through retention, skip rate, playlist compatibility, cost, and background suitability, then synthetic audio optimized along those coordinates may appear equivalent to human musical production. Yet the conclusion follows only after accepting the platform's evaluative space as an adequate representation of music. The substitution may be economically real while the claimed equivalence remains conceptually false.

This is where the earlier distinction between constraint satisfaction and semantic admissibility becomes important. An output can satisfy every explicit constraint imposed by a system while failing to preserve properties that mattered before those constraints were formalized. A generated image can have the correct dimensions, subject, style markers, palette relations, and predicted engagement while lacking whatever historical, intentional, relational, or contextual properties made the commissioned work significant. A generated text can satisfy length, tone, topic, keyword, and readability constraints while failing to perform the judgment for which a writer was originally sought. The more completely institutions redefine success through formal constraints, however, the less visible this residual becomes. What the metric cannot represent eventually ceases to appear in the account of what has been lost.

Replaceability is manufactured when this narrowing occurs before substitution. The process can be represented schematically as
\[
H
\longrightarrow
P(H)
\longrightarrow
P_{\theta}(H)
\longrightarrow
M \approx P_{\theta}(H),
\]
where \(H\) denotes a richer human practice, \(P\) a platform representation of that practice, \(P_{\theta}\) the subset selected under the platform's optimization criteria \(\theta\), and \(M\) a machine system capable of producing outputs sufficiently similar to the selected representation. The final comparison \(M \approx P_{\theta}(H)\) may be empirically defensible while the stronger conclusion \(M \approx H\) does not follow. The missing transformation is precisely the history by which \(H\) was projected into \(P_{\theta}(H)\).

The counterfactual required to evaluate replacement must therefore precede the platform transformation. The relevant question is not merely, ``Can an artificial system produce what successful platform creators currently produce?'' It is also, ``What would those creators be producing if years of ranking, monetization, recommendation, and machine-legibility constraints had not selected the present forms?'' This counterfactual cannot be recovered perfectly because the historical trajectory has already occurred. Its importance is methodological. Without it, adaptation to the platform is mistaken for evidence about the natural boundaries of human creative practice.

The mistake resembles evaluating an ecosystem after repeatedly selecting for organisms capable of surviving a pollutant and concluding that the pollutant poses little constraint because the remaining organisms tolerate it. The surviving population is evidence about the selection regime as well as about the organisms themselves. Likewise, creators who remain economically visible under platform conditions are selected partly for their ability or willingness to accommodate those conditions. Their observable work cannot be treated as an unbiased sample of what creative labour would look like under different institutions of distribution.

The talent reserve intensifies this selection because the platform can afford to lose participants whose practices do not adapt. An individual creator may refuse the preferred format, posting frequency, promotional labour, or metric optimization, but the reserve contains others willing to experiment with them. The institution therefore need not persuade every participant. It can select among participants until the visible population increasingly exhibits the desired traits. Those traits can then be redescribed as properties of successful creators rather than as responses to the selection environment.

Artificial intelligence arrives after this filtering with an additional advantage: the selected outputs and the behavioural responses surrounding them can become training material, evaluation data, or design guidance for automated systems. Human beings have already explored the space, generated examples, responded to audiences, and converged upon forms that survive the platform's metrics. The machine is not necessarily entering an untouched field and independently discovering the structure of successful cultural production. It can inherit a landscape whose contours were mapped through enormous quantities of human experimentation.

This creates a striking redistribution of exploratory cost. Human creators bear the cost of novelty, failure, experimentation, education, and cultural variation. The platform observes which experiments succeed according to its metrics. Automated systems can subsequently exploit regularities extracted from the surviving outputs. If those systems then compete with the population whose experimentation generated the regularities, the talent reserve has served as both labour supply and developmental substrate.

The resulting substitution can become recursive. Once machine-generated material is introduced, its lower marginal cost creates pressure to redefine quality further around dimensions in which inexpensive generation performs well. If background music can be produced cheaply, playlists can contain more background music; if generic illustrations can be generated instantly, clients can come to expect larger quantities and faster iteration; if routine prose can be generated continuously, publication schedules can accelerate. Human workers are then evaluated against a production regime whose scale and cadence were established by automation. To remain competitive, they may be required to adopt still more standardized workflows, making the residual human activity still easier to compare with machines.

The direction of causation is therefore not simply
\[
\text{better machine} \longrightarrow \text{human replacement}.
\]
It can instead take the form
\[
\text{metric narrowing}
\longrightarrow
\text{human adaptation}
\longrightarrow
\text{standardization}
\longrightarrow
\text{machine comparability}
\longrightarrow
\text{further metric narrowing}.
\]
Automation and standardization reinforce one another. The institution makes human labour more machine-like, discovers that machines can reproduce the resulting form, and then cites the success of reproduction as evidence that the original human contribution was unnecessary.

None of this establishes that artificial intelligence cannot produce valuable art, writing, music, software, or other cultural objects. That stronger claim is neither required nor supported by the argument. The issue is the validity of a particular inference from substitutability within a platform-defined objective to equivalence with the broader human practice. Artificial systems may participate in creative practices in genuinely novel ways. But their capacities should not be demonstrated by first narrowing the definition of creativity until it coincides with what the systems can already optimize.

The distinction also prevents a romanticization of human labour. Human creators imitate, repeat formulas, respond to markets, work under constraints, produce mediocre material, and frequently participate willingly in standardized genres. The relevant contrast is not between infinitely original humans and merely repetitive machines. It is between a heterogeneous field of possible practices and an institutional process that selects a tractable subset of that field as economically real. Humans and machines may both operate within that subset. The political question concerns who selected it, why those dimensions became decisive, and which possibilities were rendered unsustainable along the way.

This returns the argument to artificial scarcity. Once machine generation makes certain forms of production extraordinarily abundant, scarcity does not disappear. It moves. If images, songs, text, and video become cheaper to produce, attention, trust, attribution, provenance, curation, social recognition, and access to distribution become relatively scarcer. Platforms positioned at those chokepoints can therefore remain powerful even if the commodity they distribute becomes nearly free. Indeed, generative abundance may strengthen intermediaries by increasing the volume of material requiring ranking. When production approaches infinity while human attention remains finite, whoever controls selection acquires greater rather than lesser importance.

The supposed democratization of production can consequently coexist with increased concentration in distribution. Millions more people may be capable of producing technically competent material while fewer institutions determine which material becomes visible. This reproduces the central structure of the talent reserve at larger scale. The reserve expands from human aspirants to include automated production itself, making supply still more abundant while the allocation of attention remains scarce. The giant teddy bear becomes cheaper to manufacture; the carnival operator's control over the midway becomes more valuable.

The manufacture of replaceability is therefore not the final stage of enshittification but one of its possible extensions. A platform first mediates an activity, then measures it, then rewards the measurable subset, then induces participants to adapt to those measures, then accumulates the resulting behavioural and cultural data, and finally discovers that automated systems can reproduce substantial portions of the standardized output. At every stage the platform can plausibly describe itself as responding to what participants want. Yet the preferences it observes have been shaped by previous allocations, the creators it compares have been selected by previous constraints, and the task it automates has been defined through its own representational infrastructure.

The appropriate response is not to search for an ineffable human remainder that machines can never cross. Such arguments make the same mistake in reverse by attempting to fix permanently the boundary of a practice whose forms can change. The stronger response is institutional and counterfactual. Before treating machine substitutability as evidence that a human practice has become redundant, one should ask whether the comparison preserves the relevant structure of the practice, whether the evaluation criteria were independently justified, whether human activity was constrained in advance to satisfy those criteria, and which alternative trajectories of cultural development were excluded before the comparison was made.

This reframes the problem of automation. The central danger is not simply that machines become capable. It is that institutions possessing control over representation can define capability in terms that validate the substitutions they already have incentives to make. When the same actor controls the metric, the distribution channel, the training environment, and the economic comparison, apparent equivalence requires unusually careful scrutiny. A machine may indeed be better along the selected coordinate. The unresolved question is whether anyone outside the institution had sufficient power to decide that this was the coordinate that should count.

\section{Refusal and the Cost of Continuation}

If enshittification contracts the space of reachable alternatives, then the familiar response that dissatisfied participants are ``free to leave'' mistakes the existence of an exit operation, in Hirschman's sense \citep{hirschman1970exit}, for the existence of an admissible continuation after exit. These are not the same thing. A person can delete an account in seconds while requiring years to reconstruct the relationships, reputation, audience, archive, workflow, or livelihood attached to it. A worker can refuse an offered task while becoming less able to pay an immediate bill. A musician can withdraw a catalogue from a dominant streaming service while simultaneously withdrawing that catalogue from a substantial portion of the audience through which professional continuation is possible. A seller can abandon a marketplace while leaving behind years of reviews that established credibility with strangers. In each case, refusal is technically available. What varies is the state of the world that remains after refusal has been exercised. The meaningful question is therefore not simply whether a system permits the word ``no,'' but what the system has arranged for ``no'' to cost.

This distinction can be stated more formally by separating an operation of refusal from the continuation it induces. Let \(x\) denote the participant's present state within some institutional environment, and let \(R(x)\) denote the immediate state produced by refusing a particular interaction or withdrawing from the system. A purely formal conception of freedom asks whether \(R\) exists. A trajectory-sensitive conception asks instead about the set of viable continuations reachable from \(R(x)\). Denote this set by
\[
\mathcal{C}\bigl(R(x)\bigr)
=
\left\{
\gamma :
\gamma(0)=R(x),\;
\gamma \text{ satisfies the relevant admissibility constraints}
\right\}.
\]
A refusal for which \(\mathcal{C}(R(x))\) remains rich is fundamentally different from one for which the operation exists but nearly every valuable continuation has been destroyed. The first preserves agency after refusal; the second can make refusal resemble self-amputation. An institution need not remove \(R\) if it can progressively impoverish \(\mathcal{C}(R(x))\).

This is why exit rights, considered in isolation, provide such a weak safeguard against concentrated intermediaries. The strongest forms of lock-in do not prohibit exit but attach accumulated value to continued membership. The participant remains legally free while the products of previous participation become hostages to continuity. The social graph cannot travel, the reputation cannot travel, the audience cannot travel, the purchased library cannot travel, the workflow cannot travel, or the complementary service cannot interoperate with an alternative. The system effectively says that the participant may leave so long as the participant leaves behind the very things that made participation worthwhile. Freedom of exit then exists at the level of account ownership while disappearing at the level of accumulated social and economic structure.

Doctorow's emphasis on interoperability can be understood precisely as an attempt to separate refusal from destruction, continuing a line of argument about architecture as a form of regulation \citep{lessig2006code,wu2010master}. If a user can leave one interface while continuing to communicate with people who remain behind, departure ceases to require simultaneous collective migration. If a creator can change intermediaries while preserving a relationship with an existing audience, the intermediary must compete for the creator rather than hold the audience relationship hostage. If reputation can travel between marketplaces, accumulated trust ceases to function as a proprietary switching cost. If independent software can interact with a dominant service, users can refuse objectionable features without refusing the underlying network. Interoperability therefore does more than make products compatible. It changes the continuation available after a local refusal.

This distinction is visible in Doctorow's notion of ``disenshittification,'' particularly where users or competitors can modify a system against the incumbent's wishes. A conventional account treats modification as customization: users prefer one interface, feature set, or arrangement and alter the software accordingly. Doctorow's argument is more politically significant. Modification can introduce an external discipline into a system whose owner would otherwise control the complete relation between infrastructure and experience. If an intermediary increases advertising, an alternative client might remove it; if a service manipulates recommendations, another layer might reorder them; if a manufacturer disables a useful capability, independent software might restore it. The possibility of such modification limits the extent to which the owner can convert technical control into unavoidable extraction.

Anti-circumvention law consequently has effects beyond the particular acts it prohibits. Doctorow's criticism of Section 1201 of the Digital Millennium Copyright Act \citep{dmca1201} concerns the legal reinforcement of technical boundaries that might otherwise be contestable. When circumventing an access control can create legal exposure even where the underlying activity is not itself an infringement, a technical restriction acquires institutional force beyond what the architecture alone could sustain. The legal system can thereby remove edges from the graph of possible technical continuations. A route that remains computationally feasible becomes legally hazardous. The relevant contraction of possibility therefore occurs through the composition of technical and juridical constraints rather than through either domain considered separately.

This provides a useful example of why constraint analysis cannot stop at the interface. From the perspective of the software alone, an alternative client may be straightforward to construct. From the perspective of law, contract, platform policy, hardware attestation, authentication, proprietary protocols, and the economic cost of maintaining compatibility against an adversarial incumbent, the same client may be effectively unreachable. The admissibility of a continuation depends upon the conjunction of constraints operating across several layers. A formally possible technical intervention can therefore fail to constitute a practical alternative because another layer closes the route.

The structure can be represented by distinguishing a raw possibility space \(X\) from an admissible reachability relation. Let \(x,y\in X\) be technically describable states and write
\[
x \leadsto y
\]
when there exists an admissible trajectory from \(x\) to \(y\) under the relevant technical, economic, legal, social, and temporal constraints. The fact that \(y\) exists does not imply \(x\leadsto y\). Enshittification can proceed while \(X\) remains apparently unchanged because the institution acts upon \(\leadsto\) rather than upon the inventory of states. Competitors still exist, independent software can still be imagined, users can still leave, and creators can still publish elsewhere, yet the edges connecting present states to those alternatives disappear or become prohibitively expensive.

This suggests a more precise conception of lock-in as \emph{edge destruction}. The locked-in participant does not necessarily inhabit a smaller universe of logically possible states. Instead, the number or quality of admissible transitions from the present state has declined. This is why static descriptions of market structure can miss important changes. Two markets containing the same firms and nominal products may differ dramatically if one permits data portability, interoperability, independent repair, transferable reputation, and collective bargaining while the other does not. Their visible nodes can be identical while their transition structures differ.

The same idea clarifies the relation between individual and collective refusal. Some transitions cannot be made effectively by isolated participants because the value of the destination depends upon coordinated movement. A single person leaving a communications network does not transport the network. A single driver refusing a low wage may simply cause the task to be offered to another driver. A single musician withholding work from a streaming service may lose more than the service loses. The relevant continuation becomes viable only when refusal is coordinated across a sufficiently large population. Labour organization is therefore another technology of reachability \citep{olson1965collective,schelling1978micromotives}. It transforms actions that are individually self-defeating into collectively credible alternatives.

This is one reason Doctorow's insistence upon alliances between users, workers, creators, and small businesses matters. These populations are often encouraged to perceive themselves as occupying opposing sides of platform transactions. Consumers want lower prices while sellers want higher returns; passengers want cheaper rides while drivers want higher wages; listeners want inexpensive access while musicians want better compensation. The intermediary can present itself as the neutral mechanism balancing these competing interests. Yet monopsony and monopoly permit it to extract from both sides simultaneously. The passenger can pay more while the driver receives less; the advertiser can pay more while the publisher receives less; the listener's subscription can rise while the artist's remuneration remains poor. The difference accumulates at the chokepoint.

Under such conditions, the most consequential refusal may be refusal of the platform's representation of the parties as natural adversaries. A consumer and worker can have conflicting local interests while sharing a structural interest in preventing an intermediary from acquiring unilateral control over the terms of exchange. Class alliance, in Doctorow's sense, does not require pretending that all interests coincide. It requires recognizing that the counterfactual space available to each population can be contracted by the same concentration of power. Interoperability, transparency, collective bargaining, antitrust enforcement, and modification rights can increase the bargaining space on multiple sides even where subsequent bargaining remains genuinely adversarial.

The concept of refusal therefore belongs inside a theory of continuation rather than opposition. Saying no is not intrinsically emancipatory. A refusal that terminates the refuser while leaving the institution unchanged can even strengthen the institution by removing a dissident participant. Effective refusal requires enough preserved structure that another trajectory can continue. The relevant political objective is consequently not maximal refusal but \emph{non-catastrophic refusal}: institutional arrangements in which declining a particular intermediary, term, wage, recommendation regime, technical restriction, or mode of participation does not require surrendering unrelated goods accumulated through previous activity.

This also provides a criterion for evaluating apparent consent. Consent becomes increasingly weak evidence of legitimacy as the cost of refusal rises. A worker who accepts a contract under conditions in which every viable alternative has disappeared may consent in the ordinary procedural sense while revealing little about whether the contract represents a desirable institutional arrangement. A user who clicks ``accept'' after years of accumulating nonportable data may make a genuine choice while confronting a choice set deliberately shaped by prior dependency. The point is not to declare such choices unreal. It is to refuse the inference that because a choice occurred, the conditions producing it require no further evaluation.

The distinction becomes especially important in the platform's later stages. Early users may consent under highly favourable conditions, and their participation helps create the network effects that later increase exit costs. The platform subsequently changes the terms. If continued use after each modification is treated as fresh evidence of preference, the history of accumulated dependency disappears. The institution effectively obtains retroactive authorization for later extraction from consent given under earlier conditions. A trajectory-aware account instead asks how the cost of refusal changed between the original decision and the later one. The same ``yes'' can have very different evidential significance when the surrounding alternative space has contracted.

Refusal thus returns us to the central thesis of this essay. The contraction of the possible is not adequately measured by asking what actions remain formally permitted. It must be measured by examining what can continue after those actions are taken. A system in which everyone is free to leave but nobody can carry anything with them possesses a peculiar form of freedom: every door opens onto dispossession. A genuinely contestable system preserves enough continuity across institutional boundaries that leaving one intermediary does not mean leaving one's social, professional, cultural, or economic life behind.

\section{Diagnosis Is Not Repair}

Recognizing the structure described so far creates a temptation that should be resisted. Once monopoly, monopsony, lock-in, algorithmic opacity, artificial scarcity, vocational exploitation, and the manufacture of replaceability have been identified, it can appear that any intervention weakening the offending platform is thereby justified. Yet diagnosis and repair are logically distinct. A correct account of what is wrong does not by itself identify which modification will improve the system, and the history of technological regulation contains ample examples in which interventions directed at genuine problems produced unintended concentrations, new barriers to entry, or alternative forms of dependence. The distinction is especially important for an argument organized around counterfactual discipline, because an intervention must itself be evaluated against plausible alternatives rather than against the mere persistence of the diagnosed fault.

Let \(x\) denote the present state of a platform-mediated system and let \(D(x)\) represent a diagnostic claim that some undesirable structure is present. The truth of
\[
D(x)=1
\]
does not imply that an arbitrary intervention \(I\) satisfies
\[
V(I(x))>V(x),
\]
where \(V\) denotes the evaluative criterion relevant to repair. The intervention can fail, produce compensating harms, alter unrelated structures, or destroy capacities whose value was not represented in the original diagnosis. Even a highly accurate diagnosis therefore leaves open the question of repair warrant.

The null intervention must remain part of the admissible comparison class. This is not because inaction is neutral. Doing nothing permits existing trajectories to continue and can therefore have substantial consequences. Its importance is methodological: if every proposed intervention is compared only with a maximally pessimistic account of what happens without it, weak interventions will appear artificially successful. The relevant question is whether a proposed repair improves upon the best justified estimate of available alternatives, including nonintervention, alternative interventions, combinations of interventions, and interventions delayed until better evidence becomes available.

This matters particularly for antitrust. Doctorow invokes the historical dismantling or disciplining of concentrated firms as evidence that monopolistic power is politically alterable. The structural case for antitrust can be strong while leaving substantial questions about the form intervention should take. Breaking a firm into several entities can reduce concentration under some conditions, but a poorly designed division can leave shared technical standards, network effects, data advantages, or vertically integrated chokepoints largely intact. Conversely, forcing structural separation can destroy beneficial integrations if the relevant dependencies are misunderstood. The diagnosis ``concentration has removed competitive discipline'' therefore establishes the need to examine intervention; it does not make every decomposition of the firm admissible.

Interoperability presents a similar problem. Mandating open interfaces can reduce switching costs and enable competitors to route around objectionable platform choices, but the design of interoperability determines what is preserved and what new vulnerabilities are created. A requirement that merely exposes data without preserving semantic structure may produce nominal portability without useful continuation. An interface designed without attention to privacy or security can externalize risks onto users. A standard controlled by the incumbent can formalize rather than eliminate dependence. The repair objective cannot therefore be ``more interoperability'' in the abstract. It must specify which relations should remain possible across institutional boundaries, under what constraints, and with what protections against newly introduced harms.

Portability provides an instructive example of the distinction between syntactic preservation and functional continuation. A platform can satisfy a formal export requirement by supplying an archive containing the user's data while failing to provide any environment in which the archive recreates the relationships that made those data useful. A list of social connections is not a social graph if the listed people cannot be reached through the exported structure. A history of reviews does not function as reputation if another marketplace has no reason to trust or interpret it. A playlist file does not preserve a listening environment if the referenced works cannot be resolved across services. Export preserves representation; portability requires reachability.

This is precisely where repair must be structure-sensitive. The objective is not to preserve every byte independently but to preserve the relations whose destruction creates lock-in. In some cases this requires common identifiers, authentication mechanisms, provenance, or reciprocal protocols rather than simple copying. In others, preserving the relation may be undesirable because the relation itself creates a privacy or security risk. Repair therefore requires an independently justified account of which structures deserve preservation. Without such an account, ``openness'' can become as empty a proxy as ``engagement.''

The same caution applies to transparency. Algorithmic opacity contributes to superstition and asymmetric bargaining power, but maximal disclosure is not automatically an adequate repair. Publishing source code may reveal little about a system whose behaviour depends upon enormous datasets, continuously changing parameters, experiments, or downstream institutional practices. Detailed disclosure can also enable gaming, expose private information, or overwhelm the people supposedly empowered by it. What workers and creators often require is not total access to every internal computation but reliable knowledge of the dimensions relevant to bargaining and contestation: how compensation is determined, whether individualized factors alter offers, what forms of discretionary promotion exist, what evidence supports adverse decisions, and what procedures permit those decisions to be challenged.

Transparency should therefore be evaluated according to the action it makes possible. Information that cannot alter any reachable continuation may satisfy a disclosure requirement while leaving power unchanged. Conversely, a relatively small amount of strategically chosen information can transform bargaining if it allows participants to compare outcomes, identify discriminatory treatment, coordinate responses, or verify whether contractual promises are being honoured. The value of transparency is relational rather than volumetric. The question is not how much the platform reveals but what participants can now do that they could not do before.

Labour organization supplies another example. Collective bargaining can counteract monopsony by aggregating individually weak refusals into a credible collective alternative. Yet the institutional form matters. A bargaining unit that excludes important classes of dependent workers may simply relocate exploitation. A legal category designed around conventional employment may fail to capture workers whose dependence is real but whose formal status is independent. A creator union capable of negotiating with one intermediary may have little leverage if distribution can be shifted toward a vast reserve of nonmembers. Repair therefore requires attention to the boundaries through which the relevant population is defined.

These difficulties do not weaken Doctorow's case for restoring external disciplines. They clarify why those disciplines matter. Competition, regulation, interoperability, labour organization, and modification rights are not magic solutions that produce good outcomes automatically. They are mechanisms that prevent one institution from becoming the sole author of the environment in which outcomes are evaluated. Their value lies partly in preserving the possibility of correction after mistakes are made. A competitive market can still produce bad products, a union can make poor decisions, an interoperable protocol can be badly designed, and a regulator can adopt an ineffective rule. But where alternatives and revision remain reachable, these errors need not become permanent merely because one institution controls the infrastructure required to challenge them.

Repair warrant therefore depends upon more than the expected immediate improvement in a selected metric. It should include the effect of an intervention upon future corrigibility. An intervention that produces a modest short-term benefit while permanently centralizing control may be inferior to one that produces a smaller immediate benefit while preserving the capacity for subsequent revision. Conversely, a technically elegant intervention that creates an irreversible dependency can be dangerous even if its first-order effects appear positive. The trajectory matters because repair changes not only the present state but the space of repairs available later.

A simple representation makes this distinction explicit. Let an intervention \(I\) be evaluated not only by its immediate effect on some welfare or performance measure \(V\), but by its effect on future reachability \(\mathcal{R}\), evidential uncertainty \(U\), implementation cost \(K\), and irreversible risk \(Z\). A schematic repair criterion might take the form
\[
\mathcal{W}(I\mid x)
=
\Delta V(I\mid x)
+
\alpha\,\Delta \mathcal{R}(I\mid x)
-
\beta U(I\mid x)
-
\gamma K(I\mid x)
-
\delta Z(I\mid x),
\]
with the null intervention \(I_0\) included explicitly in the comparison class. The coefficients are not offered as universal constants or as an attempt to reduce political judgment to arithmetic. The expression simply prevents several dimensions of repair from disappearing behind a single outcome measure. In particular, an intervention that improves an immediate metric while substantially reducing future reachability can be recognized as structurally dangerous.

This criterion also distinguishes repair from punishment. A policy can impose costs upon a platform without restoring any of the counterfactual disciplines whose disappearance enabled enshittification. A fine may become a predictable cost of doing business; a one-time settlement may compensate past harm while leaving the chokepoint intact; a prohibition against one practice may cause extraction to migrate toward another mechanism. If the objective is disenshittification, the relevant test is not whether the incumbent suffers but whether the surrounding system becomes less dependent upon the incumbent's continued restraint.

The strongest interventions may therefore be those that make good behaviour less dependent upon virtue. A platform that refrains voluntarily from exploiting its users is preferable to one that exploits them, but voluntary restraint is historically fragile. Leadership changes, investor pressure changes, competitive conditions change, and the accumulated value created under one governance regime can become the extraction target of another. Structural repair attempts to make exploitation costly even when the people controlling the institution would prefer it. It restores constraints rather than requesting benevolence.

This is an important distinction because narratives of corporate decline frequently personalize institutional trajectories. A founder cared about the product; a later chief executive cared about revenue; engineers lost an internal political struggle to financial managers; the culture deteriorated. Such histories can be accurate and still leave the deeper question unanswered: why did the institution acquire enough freedom that the worse internal faction could impose its objective without triggering effective external correction? Doctorow's answer repeatedly returns to the removal of disciplines. Repair should therefore be evaluated according to whether it restores those disciplines in durable form.

The null intervention remains important even here. Not every concentrated system can be safely transformed through the same mechanism, and some proposed remedies can entrench the largest firms by imposing compliance costs that smaller competitors cannot bear. Regulation can itself become a barrier to entry; standardized interfaces can ossify innovation; portability mandates can create new privacy risks; restrictions intended to protect creators can be captured by large rights holders; and poorly designed AI regulation can strengthen incumbents possessing the resources required for compliance. These possibilities are not arguments for doing nothing. They are reasons to apply the same counterfactual discipline to reform that the essay applies to platforms.

A detected error therefore does not automatically license correction by whatever means are presently available. Repair becomes warranted when an intervention can be shown, with appropriate uncertainty, to improve upon the admissible alternatives while accounting for structural effects, costs, and irreversible risks. Where no sufficiently grounded comparison is available, the correct response may be to gather evidence, preserve reversibility, conduct limited experiments, or alter the representational scheme through which the problem itself has been formulated. This last possibility is especially important. Sometimes the failure lies not in selecting the wrong intervention within an accepted coordinate system but in the coordinate system's inability to represent what matters.

Enshittification repeatedly presents this latter problem because the incumbent's metrics tend to dominate public description of the system. Engagement, growth, retention, creator counts, upload volume, gross merchandise value, ride volume, subscriber numbers, or advertising efficiency can all improve while the system becomes less contestable. A repair evaluated solely through those same metrics may therefore reproduce the diagnostic blindness of the institution being repaired. The relevant reference class must be independently grounded. If the problem is the contraction of reachable alternatives, then some measure of restored reachability must enter the evaluation rather than being assumed to follow automatically from improved performance elsewhere.

This brings the argument toward a more unified interpretation of Doctorow's proposed disciplines. Antitrust, interoperability, portability, labour organization, modification rights, and limits upon anti-circumvention law appear at first to belong to different domains of policy. One concerns market structure, another technical standards, another labour law, another intellectual property, another consumer rights. Yet each can be interpreted as acting upon the same structural variable: the number and quality of continuations available when an incumbent behaves badly. Their common purpose is not to guarantee that every participant obtains the outcome they prefer. It is to prevent the intermediary from becoming sufficiently indispensable that dissatisfaction loses causal force.

\section{Disenshittification and the Restoration of Counterfactual Space}

Doctorow's account of enshittification is sometimes summarized as a story about companies beginning well and becoming greedy. Such a summary misses the most politically consequential feature of the theory. If deterioration resulted primarily from a change in corporate morality, the natural remedy would be better executives, more conscientious investors, improved corporate culture, or public pressure capable of persuading firms to rediscover their original values. Doctorow's emphasis on lost disciplines points elsewhere. The institution deteriorates because the forces that once made deterioration costly have weakened. Disenshittification therefore requires more than reversing individual bad decisions. It requires reconstructing an environment in which future bad decisions again encounter consequences.

Competition is one such consequence, but only when competition is understood dynamically rather than nominally. A competitor that cannot receive portable data, interoperate with existing users, inherit reputation, communicate across network boundaries, or overcome accumulated switching costs may exist without exerting substantial discipline. Effective competition requires that movement toward the competitor remain sufficiently inexpensive to be credible. Market structure and technical architecture therefore cannot be separated cleanly. Interoperability can transform a small competitor from an isolated destination into a reachable continuation, while enclosure can transform numerous nominal competitors into strategically irrelevant alternatives.

Labour organization acts on the same structure from another direction. The isolated worker's refusal may be too costly to exercise, particularly where a talent reserve or labour surplus permits immediate substitution. Collective organization changes the transition structure by making coordinated refusal possible. The worker no longer confronts only the counterfactual ``accept this term or lose access to the opportunity.'' Another state becomes reachable: ``refuse this term together while preserving enough collective capacity to force renegotiation.'' The importance of organization therefore lies not only in the bargaining demand ultimately made but in the creation of a new branch in the possibility space.

Modification rights similarly create branches. A user dissatisfied with a particular feature need not choose only between accepting the feature and abandoning the entire service if the feature can be removed, replaced, or routed around. This decomposes what the platform would prefer to bundle into a single indivisible decision. Bundling is itself a source of power because it attaches valuable and objectionable components to the same act of participation. Modifiability permits selective refusal. The user can retain the network while rejecting the advertisement, retain the device while replacing the software, retain the purchased object while obtaining independent repair, or retain communication while changing the interface through which communication occurs.

Portability creates continuity across branches. It ensures that choosing another path does not require returning to an initial state. This is particularly important because accumulated state is one of the principal mechanisms through which platforms acquire power. A system in which migration resets reputation, history, relationships, or purchased assets effectively taxes every alternative according to the duration of prior participation. The longer the user stays, the more expensive departure becomes. Portability attempts to reverse this temporal asymmetry by allowing at least some accumulated state to survive the transition.

Taken together, these mechanisms suggest a definition of contestability stronger than the existence of competitors. A system is contestable to the extent that participants can redirect significant portions of their activity, relationships, accumulated state, and bargaining power toward alternative arrangements without incurring losses so large that the alternative ceases to discipline the incumbent. Contestability is therefore relational and trajectory-dependent. It concerns not merely what alternatives exist but what can be carried across the transition toward them.

This formulation also explains why the opposite of enshittification is not simply quality. A monopolistic platform can be excellent. It can employ conscientious people, provide extraordinary services, compensate creators generously, protect privacy, and refuse obvious opportunities for exploitation. Nothing in the theory requires denying such cases. The problem is that quality under concentrated control can remain contingent upon voluntary restraint. If participants cannot carry their accumulated relations elsewhere, the platform possesses an option to extract that may become attractive under future conditions. Present quality does not eliminate the structural possibility of later enshittification.

Nor is the opposite of enshittification abundance. Platforms can supply immense abundance while controlling the routes through which abundance becomes usable. A catalogue containing almost every song is not structurally open if a single intermediary determines the economic terms on which musicians reach listeners. An infinite supply of generated images does not produce a pluralistic cultural environment if a few systems determine which images are discoverable, licensable, trusted, or monetizable. Abundance at the level of objects can coexist with severe contraction at the level of relations.

The opposite is better described as contestability because contestability preserves an external reference point. An institution operating within a contestable environment cannot define success entirely according to its own metrics because participants retain credible ways of responding when those metrics diverge too far from their purposes. The search engine cannot indefinitely increase monetizable searches by degrading results if users can easily move their histories, defaults, and habits to a competitor. The streaming service cannot indefinitely reduce the share reaching musicians if artists and listeners possess alternative routes to one another. The labour platform cannot personalize wages down to individual desperation if workers can compare offers, coordinate refusal, and move reputational capital across services. The social platform cannot indefinitely suppress chosen relationships in favour of profitable recommendations if users can select another client or carry those relationships elsewhere.

Contestability therefore preserves something analogous to an experimental control. The incumbent's claims about what users want can be tested against what users do when meaningful alternatives become available. The claim that creators prefer a particular distribution arrangement can be tested when audiences are portable. The claim that workers accept individualized wages can be tested when wage information is collectively visible and refusal does not mean exclusion from the labour market. The claim that a feature is valuable can be tested when users are permitted to remove it. Without such alternatives, revealed behaviour becomes epistemically contaminated by dependency.

This epistemic function of contestability is easy to overlook. Competition and exit are normally justified economically: they constrain prices and improve quality. But they also generate information about preference that concentrated systems suppress. When people can choose among genuinely reachable alternatives, differences in their behaviour become more informative about what they value. When all routes pass through one intermediary, the intermediary can point to continued participation as evidence of satisfaction even when participation primarily reveals the absence of viable alternatives. Restoring counterfactual space therefore improves not only bargaining conditions but the quality of social knowledge.

The same principle applies to creative practice. A cultural environment containing multiple distribution channels with different evaluative criteria permits forms to survive that would disappear under a single optimization regime. One institution may reward immediacy, another depth, another local relevance, another experimentation, another archival value, another live performance, another slow accumulation of reputation. The diversity of institutions prevents one representational limitation from becoming a universal constraint upon expression. Cultural pluralism is therefore partly infrastructural. It depends upon preserving multiple routes through which work can acquire an audience and meaning.

This provides another response to the manufacture of replaceability. If one platform's metrics define the economically relevant subset of creative practice, machine competence within that subset can be mistaken for competence across the practice as a whole. A contestable cultural environment preserves alternative evaluative regimes against which the comparison can be tested. Artificial systems may flourish in some of them and fail in others; human and machine practices may combine in unexpected ways; new categories may emerge that neither existing platforms nor existing models anticipated. The important point is that no single intermediary should possess sufficient control to define the comparison class in advance.

Disenshittification thus becomes less a programme for restoring a lost golden age than a programme for preserving revision. Early platforms were not necessarily virtuous, and decentralized systems can contain their own pathologies. The objective is not to recreate a particular historical configuration but to prevent any configuration from becoming immune to correction. A healthy system retains the ability to generate alternatives faster than incumbents can close them, to preserve accumulated value across transitions, and to permit local refusal without requiring global abandonment.

The normative criterion is consequently dynamic. We should ask of an institutional arrangement not only whether it produces acceptable outcomes today but whether people harmed by tomorrow's outcomes will possess credible means to alter, leave, modify, compete with, or collectively renegotiate the arrangement. This criterion treats future agency as a present good. It values the maintenance of unused alternatives because unused alternatives discipline actual behaviour and because the future conditions under which they may become necessary cannot be known completely in advance.

Such unused capacity can appear inefficient. Interoperability requires maintaining interfaces that some users never invoke; portability requires preserving export mechanisms many people never use; competition duplicates infrastructure; labour organization can constrain managerial flexibility; rights to modification permit behaviour manufacturers would prefer to prohibit. From the standpoint of immediate optimization, these can resemble friction. From the standpoint of contestability, the friction is part of the safety mechanism. A perfectly optimized system with no redundant paths can become extremely efficient at producing the wrong outcome once its objective diverges from the interests of those who depend upon it.

This is where Doctorow's political argument ultimately converges with the broader theory of constraint developed here. Constraint is not intrinsically opposed to freedom. Some constraints preserve freedom by preventing actors from destroying the conditions under which alternatives remain available. Antitrust constrains concentration in order to preserve competitive possibility; interoperability obligations constrain enclosure in order to preserve migration; labour law constrains unilateral contracting power in order to preserve collective refusal; anti-discrimination rules constrain individualized extraction in order to preserve fairer bargaining; limits on anti-circumvention rules constrain control over modification in order to preserve technical alternatives. Freedom can therefore require constraints upon the ability to contract the possibility space of others.

The relevant distinction is between constraints that preserve a plural field of continuation and constraints that collapse it around a privileged actor. A protocol constrains participants by requiring compatible forms, but that shared constraint can enable communication among otherwise independent systems, much as commonly governed infrastructure can sustain cooperation without central ownership \citep{ostrom1990governing}. A monopoly constrains participants by making one institution the necessary route between them. Both impose structure; their political effects differ because one can increase reachability while the other decreases it. The question ``Does this rule constrain?'' is therefore insufficient. The more important question is ``What becomes reachable or unreachable because this constraint exists?''

Under this criterion, disenshittification is not the elimination of constraint but its redistribution. The platform loses some freedom to close interfaces, discriminate secretly, bundle unrelated functions, or appropriate accumulated relationships. Users, workers, creators, competitors, and complementary institutions gain corresponding freedom to construct continuations that do not require the platform's permission. The system becomes less unconstrained at the centre and more open at the edges.

This is also why Doctorow's argument remains hopeful without requiring confidence that corporations will voluntarily reform themselves. Hope does not depend upon predicting benevolent behaviour from concentrated institutions. It depends upon the continued existence of interventions capable of changing the structure in which those institutions act. The distinction becomes the final step of the argument because it separates a theory of reachable alternatives from both technological optimism and technological pessimism. If optimism says that progress will solve the problem and pessimism says that concentrated power will inevitably prevail, both treat the future primarily as an endpoint to be predicted. Hope asks a different question: what transition can be made from the present state that changes what becomes reachable next?

\section{Hope as a Theory of Reachable Futures}

Doctorow's distinction between optimism, pessimism, and hope is sometimes heard as a difference in emotional temperament, but its deeper significance is practical. Optimism predicts that things will improve; pessimism predicts that they will deteriorate. Both can become forms of passivity when the prediction substitutes for intervention. If improvement is inevitable, action is unnecessary because history will deliver the desired outcome. If deterioration is inevitable, action is unnecessary because history cannot be redirected. Hope occupies a different logical category. It need not assign a high probability to success. It requires only that some intervention remains capable of altering the distribution of future possibilities.

This distinction can be expressed in trajectory terms. Let \(x_t\) denote the present state and let \(\mathcal{R}(x_t)\) denote the set of states reachable from it under presently available transitions. Optimism and pessimism are principally claims about where trajectories in \(\mathcal{R}(x_t)\) are likely to terminate. Hope is concerned with operations that change \(\mathcal{R}(x_t)\) itself. An intervention can be hopeful even when it does not immediately produce the desired terminal state if it opens additional admissible transitions from the next state:
\[
\mathcal{R}\bigl(I(x_t)\bigr)
\supsetneq
\mathcal{R}(x_t)
\]
along dimensions relevant to future collective agency. The immediate gain may be modest while the expansion of subsequent possibility is substantial.

This interpretation fits Doctorow's emphasis on political history. Concentrated industrial power has previously appeared immovable, labour organization has repeatedly been declared obsolete, technical monopolies have repeatedly been described as natural consequences of superior efficiency, and regulatory interventions have repeatedly been dismissed as impossible until political conditions changed. The lesson is not that every campaign eventually succeeds. It is that present constraint does not establish future impossibility. Political action can alter the transition structure itself by creating organizations, laws, standards, coalitions, precedents, technical tools, and shared concepts that were absent from the original state.

This is why the first successful intervention can matter beyond its direct effect. A small victory demonstrates that an apparently fixed boundary is movable; the demonstration changes expectations; changed expectations alter willingness to coordinate; coordination makes larger interventions possible. Doctorow describes hope in terms of taking a step toward a better state and discovering from the new position that further steps have become visible. The metaphor is topological as much as motivational. The next reachable state need not contain the final solution. It need only reveal transitions that could not be traversed, or perhaps even perceived, from the original position.

Pessimism can obscure this structure by demanding a complete path before permitting the first step. If an intervention does not solve monopoly, labour precarity, algorithmic opacity, creator exploitation, artificial intelligence, and cultural concentration simultaneously, it can be dismissed as inadequate. Yet complex systems rarely permit such global repair. Local interventions alter constraints, and altered constraints change which subsequent interventions become feasible. A portability requirement can enable competitors that later create political constituencies for stronger interoperability. A successful organizing campaign can generate information that makes algorithmic discrimination easier to demonstrate. A right to modification can produce tools whose existence changes expectations about what ownership should permit. The path is constructed partly through traversal.

Optimism commits the complementary error when it assumes that technological abundance will automatically produce institutional abundance. Lower production costs, more powerful artificial intelligence, cheaper computation, or larger networks do not themselves guarantee broader distribution of power. As the preceding analysis has shown, abundance in one dimension can increase scarcity in another. If generation becomes nearly free while attention remains finite, control over ranking can become more valuable. If communication becomes effortless while identity and social graphs remain proprietary, network lock-in can intensify. If workers become more productive while a monopsonistic intermediary controls access to employment, productivity gains need not accrue to workers. Technological progress changes the feasible set; institutions influence who can traverse it.

Hope therefore requires institutional imagination rather than faith in technical inevitability. It asks which constraints can be changed, which dependencies can be made portable, which refusals can be made non-catastrophic, which hidden comparisons can be rendered contestable, which collective actors can be created, and which apparently indivisible systems can be decomposed into components that no longer need to share a single owner. These are questions about constructing alternatives rather than predicting whether existing institutions will improve.

They also return us to the problem of counterfactuals with which the essay began. A counterfactual is politically significant when it can act upon the present. The possibility of leaving disciplines the platform before anyone leaves; the possibility of organizing disciplines the employer before a strike occurs; the possibility of modification disciplines the manufacturer before a user circumvents a restriction; the possibility of antitrust action disciplines acquisition before a firm becomes unchallengeable. The unrealized alternative exerts causal force because actors know that it can become real.

Enshittification weakens precisely this force. It converts alternatives into fantasies by increasing the cost of reaching them, then points to the absence of migration as evidence that participants prefer the actual state. Hope reverses the operation by turning apparently hypothetical alternatives back into credible continuations. It is therefore not an emotion added after the structural analysis is complete. It is the practical counterpart of the analysis. If power contracts the possible, political agency consists partly in reopening it.

The deepest disagreement between this position and technological fatalism concerns what should count as given. A fatalistic account treats current market concentration, platform architecture, intellectual-property rules, employment classifications, recommendation systems, or ownership structures as environmental facts to which individuals must adapt. Doctorow insists upon their historical construction. They are products of laws, enforcement priorities, technical decisions, mergers, standards, institutional struggles, and previous political settlements. What has been constructed can be reconstructed, although reconstruction is neither costless nor guaranteed to succeed.

The same point applies to the apparent inevitability of machine replacement. If replaceability has partly been manufactured through metric narrowing, platform selection, labour fragmentation, and the concentration of distribution, then those institutional conditions should not be smuggled into forecasts as laws of nature. Artificial intelligence may transform creative work profoundly, but the distribution of the resulting benefits and losses remains contingent upon who controls the infrastructure through which the transformation occurs. A world in which creators possess strong bargaining institutions, portable audiences, provenance rights, independent distribution channels, and meaningful control over the use of their work is not economically equivalent to one in which a handful of intermediaries control models, distribution, compensation, and visibility simultaneously. The underlying technology can be similar while the reachable social futures differ radically.

Hope in this sense is disciplined rather than consoling. It does not require believing that every desirable future remains reachable. Some paths genuinely close. Institutions can destroy capacities, communities can dissolve, knowledge can be lost, ecosystems can cross irreversible thresholds, and technical infrastructures can become sufficiently entrenched that changing them becomes extraordinarily costly. The recognition of path dependence is precisely why preserving alternatives before they disappear matters. A theory of reachability is not reassuring by default. It tells us that delay can alter the space in which later action occurs.

This gives political urgency to apparently technical questions about standards, protocols, portability, repair, and interoperability. They are mechanisms through which future paths are either preserved or removed. A closed interface today can determine whether a competitor is viable tomorrow; a nonportable reputation system can determine whether a worker can leave next year; a proprietary social graph can determine whether a community can migrate when moderation or monetization changes; an anti-circumvention rule can determine whether a later generation is legally permitted to restore functionality abandoned by a manufacturer. Infrastructure is frozen counterfactual structure. It stores decisions about which future transitions will be easy, difficult, or impossible.

The same is true of cultural infrastructure. A society that permits a few recommendation systems to become the dominant routes between creators and audiences does not merely make a present distribution decision. It trains creators and audiences within those systems, shapes expectations about form and duration, accumulates behavioural data under their categories, and makes alternative cultural institutions harder to sustain. The future inherits not only the technology but the habits produced around it. The contraction of the possible can therefore occur gradually enough to resemble preference.

Against this, hope preserves the distinction between what people have adapted to and what they would choose under a richer set of alternatives. It refuses to infer legitimacy from survival alone. A person can become extraordinarily skilled at navigating a bad system. A creator can master an opaque algorithm, a driver can learn how to avoid discriminatory offers, a user can develop elaborate strategies for filtering unwanted recommendations, and a business can reorganize itself around an intermediary's changing rules. Such adaptations demonstrate human competence. They do not demonstrate that the environment eliciting them is desirable.

Nor should successful adaptation become an argument against repair. This is another form of survivorship bias. The people still visible are disproportionately those who learned to survive the constraint. Those who left, failed, changed careers, abandoned projects, or never entered because the expected conditions were unacceptable disappear from the observed population. The system then points to its survivors as evidence that survival is sufficiently available. The giant teddy bear reappears at the scale of institutions.

Hope requires recovering the missing counterfactual population: not as an imaginary collection of guaranteed successes, but as evidence that the observed equilibrium is only one trajectory among those that institutional design could have supported. The artist who stopped creating because distribution became intolerable, the worker who left a vocation because care became exploitable, the small competitor that could not overcome proprietary network effects, and the user who remained because departure meant social isolation all belong to the causal history of the system even when they disappear from its performance metrics.

The political task is therefore not to maximize the number of options presented within a fixed enclosure. It is to preserve the ability to alter the enclosure itself. A platform can offer thousands of settings while prohibiting alternative clients; a marketplace can contain millions of products while controlling every transaction; a streaming service can offer nearly every recording while determining the route between every listener and artist. Choice inside a system is not equivalent to choice among systems, and choice among systems is not meaningful when accumulated state cannot move between them.

A richer conception of freedom must consequently include the ability to preserve continuity while changing institutions. This is what portability, interoperability, collective organization, and modification rights share. Each separates something valuable from the particular intermediary presently controlling it. The social relation is separated from the social platform, the reputation from the marketplace, the creative audience from the distributor, the device from the manufacturer's software, and the worker's bargaining capacity from the isolated employment decision. The intermediary becomes replaceable because the relation it mediates is no longer identical with the intermediary itself.

There is an instructive symmetry here. Platform capitalism frequently seeks to make workers, creators, sellers, and complementary businesses replaceable while making the platform indispensable. Disenshittification reverses that asymmetry. It does not seek to make human beings indispensable in every function. It seeks to make the intermediary replaceable while preserving the people and relationships that give the system value. This may be the most concise way to distinguish competitive technical progress from enclosure: a healthy infrastructure permits components to change without requiring the destruction of the relations that depend upon them.

The contraction of the possible reaches its limit when the intermediary becomes indistinguishable from the activity it mediates: social life becomes ``Facebook,'' search becomes ``Google,'' music becomes ``Spotify,'' online retail becomes ``Amazon,'' gig driving becomes ``Uber,'' and cultural visibility becomes whatever the dominant recommendation systems happen to display. At that point, refusal of the company is experienced as refusal of the activity itself. The institutional name has occupied the conceptual space of the underlying relation.

The restoration of possibility begins by separating them again. People need social communication, not necessarily a particular social network; they need discovery, not necessarily a particular search engine; musicians and listeners need ways to find one another, not necessarily a particular streaming intermediary; workers and customers need coordination, not necessarily a particular labour platform. Once the underlying relation is distinguished from its current institutional realization, alternatives become conceptually available before they become technically or politically reachable.

That conceptual distinction is modest, but it is where hope begins. A system appears inevitable most completely when its contingent implementation has become indistinguishable from the function it performs. Naming the distinction does not solve the problem, but it restores the possibility of asking how the function might be realized otherwise. From there, technical design, collective organization, law, regulation, standards, and alternative institutions can begin to turn the conceptual counterfactual into an actionable one.

The final measure of a digital institution should therefore not be whether it promises never to enshittify. Such promises are structurally weak because they ask future participants to trust the continued benevolence of whoever inherits control. A stronger question is whether the institution could enshittify \emph{without losing the people who make it valuable}. Can users carry their relationships elsewhere? Can workers coordinate refusal? Can creators retain access to audiences? Can independent software alter the experience? Can competitors interoperate rather than reconstruct the network from zero? Can accumulated state survive migration? Can the institution's own metrics be challenged by observations generated outside its control?

Where the answer is yes, degradation encounters a boundary. Where the answer is no, quality remains contingent upon restraint.

The opposite of enshittification is therefore not excellence. It is not abundance, innovation, convenience, growth, or even benevolence. All of these can exist inside an enclosure and can contribute to making the enclosure harder to leave. The more durable opposite is \emph{contestability}: a condition in which enough independent and reachable continuations remain available that no institution can unilaterally redefine deterioration as success. Hope is the practice of keeping those continuations real.

\section{The Political Economy of the Chokepoint}

The argument can now be stated in a form that makes the relation among Doctorow's examples clearer. A chokepoint is not merely a large firm, a popular service, or an intermediary positioned between two groups. Intermediation can be useful precisely because coordination is costly. Search engines reduce the cost of finding information, streaming services reduce the cost of finding and delivering recordings, marketplaces reduce the cost of matching buyers and sellers, and labour platforms can reduce the cost of matching workers with tasks. The relevant transformation occurs when an intermediary becomes sufficiently difficult to bypass that the value created by reducing coordination costs can be converted into the power to impose new costs upon the parties being coordinated. The same position that initially solves a coordination problem becomes the position from which the resulting relationships can be taxed.

This conversion explains why the first phase of enshittification cannot be dismissed as a simple deception. The platform may genuinely solve a problem better than its predecessors. It may subsidize participation, remove friction, simplify distribution, improve discovery, or permit forms of coordination that were previously expensive or impossible. These improvements attract participants, and the participants in turn make the platform more useful to one another. The intermediary gradually ceases to be valuable merely because of the service it directly supplies. It becomes valuable because of the relations that have accumulated through it. At this point, the platform's asset is no longer adequately described by its software. Its asset includes the difficulty of reconstructing elsewhere the network of relations that its previous usefulness helped to create.

The distinction between value creation and value capture is therefore temporally unstable. An institution can create substantial value at \(t_0\), use that value to attract a population at \(t_1\), acquire control over relations among members of that population at \(t_2\), and exploit the resulting dependence at \(t_3\). Looking only at \(t_0\) produces an excessively celebratory account because the later extraction remains invisible. Looking only at \(t_3\) produces an excessively simple account of predation because it cannot explain why rational participants entered the system in the first place. The trajectory itself is the explanatory object. The platform becomes capable of extracting because it was useful enough to become difficult to avoid.

This temporal structure is one reason conventional measures of market power can underestimate digital chokepoints. Price is often an incomplete indicator because one side of the platform may be subsidized or nominally free. A social network charging users nothing can still accumulate extraordinary power over the conditions through which those users encounter one another. A marketplace can offer sellers access to enormous demand while progressively increasing the price of visibility within that demand. A search engine can remain free at the point of use while altering the ratio between useful results and monetizable interactions. The absence of a direct price increase therefore says little about whether the intermediary has begun consuming the surplus created by accumulated dependence.

The same problem appears when consumer welfare is treated as separable from supplier welfare. A streaming service can provide listeners with extraordinary access while placing musicians under monopsonistic pressure \citep{eriksson2019spotify,prey2020musicians,hesmondhalgh2019cultural}. Cheap consumer access can therefore coexist with deteriorating conditions of production. If the analysis terminates at the consumer interface, the extraction remains partially invisible because its cost has been displaced onto another population. Conversely, raising payments to suppliers does not automatically establish improvement if the intermediary finances those payments by increasing extraction elsewhere. Multi-sided platforms require a relational accounting in which changes on one side are traced through the entire system rather than evaluated as isolated transfers.

Doctorow's use of monopsony \citep{doctorowgiblin2022chokepoint,azarmarinescu2022monopsony} is crucial here because it restores a symmetry often missing from popular accounts of monopoly, one developed at length in the broader antitrust literature on concentrated market power \citep{wu2018curse,stoller2019goliath,posner2018radicalmarkets}. Monopoly concerns the power of a seller facing buyers with insufficient alternatives. Monopsony concerns the power of a buyer or gatekeeper facing suppliers with insufficient alternatives. A platform can occupy both positions at once. It can become a necessary route through which consumers obtain access to a class of goods and simultaneously a necessary route through which producers obtain access to consumers. Its power does not merely arise from owning what is exchanged. It arises from controlling the crossing.

We can express this in deliberately abstract terms. Let \(P\) denote a set of producers, \(C\) a set of consumers, and \(M\) an intermediary through which a large fraction of economically significant relations between them are routed. If direct or alternative relations \(p\leftrightarrow c\) remain inexpensive, then \(M\) must continuously justify its presence by reducing transaction costs or supplying some other value. As alternative edges disappear, however, the effective topology approaches
\[
p \longrightarrow M \longrightarrow c.
\]
The intermediary then controls not merely an exchange but the conditions under which the two populations become mutually visible. Its bargaining power derives from the fact that each side's access to the other passes through the same node.

The most consequential form of concentration is therefore not always concentration of production. It can be concentration of \emph{routing}. A platform need not own the music if it controls the dominant route through which music is discovered. It need not employ the driver if it controls the dominant route through which riders and drivers are matched. It need not produce the goods if it controls the marketplace through which buyers search for them. It need not create the news if it controls the mechanisms through which audiences encounter it. The chokepoint is powerful because the relation, rather than the object, has been enclosed.

This relational enclosure also explains why superficially pro-competitive increases in supply can fail to discipline the intermediary. Adding another million musicians does not create another route between musicians and listeners. Adding another million sellers does not create another marketplace. Adding another million drivers does not create another dispatch network. Supply-side abundance can instead strengthen the chokepoint by making individual suppliers more substitutable while leaving the routing infrastructure unchanged. The talent reserve is therefore not external to chokepoint capitalism. It is one of the conditions that can stabilize it.

The intermediary's ideal position is consequently asymmetric: abundance among participants and scarcity among routes. It benefits when users have many creators to choose from, creators face many competitors, buyers have many sellers, sellers face many substitutes, employers have many workers, and workers have few employers or gateways through which comparable demand can be reached. The platform can celebrate abundance because the abundance exists below the level at which its own power is concentrated.

This distinction provides a more exact account of artificial scarcity. The platform need not make cultural objects, workers, products, or information scarce. It can allow or encourage them to become extraordinarily abundant while keeping scarce the pathways through which abundance becomes economically effective. Recommendation slots are scarce, prominent search positions are scarce, verified identities may be scarce, access to particular audiences is scarce, and bargaining access to the intermediary itself can be scarce. Scarcity is therefore displaced upward from production into coordination.

Once this occurs, the intermediary can monetize not merely the transaction but the attempt to become visible enough for a transaction to occur. Sellers purchase advertising to recover visibility within a marketplace to which they already pay fees. Creators purchase promotional services or devote uncompensated labour to satisfying ranking systems in order to reach audiences already present on the platform. Businesses optimize listings for search systems whose degradation can increase the amount of optimization required. Workers spend time waiting, repositioning, monitoring offers, or learning opaque acceptance strategies. The chokepoint converts access to the possibility of exchange into a separately monetizable commodity.

There is a recursive quality to this process. The platform first reduces a friction, thereby attracting activity. It then becomes the dominant infrastructure through which the activity occurs. It can subsequently introduce new friction at strategically chosen locations because participants can no longer avoid the infrastructure cheaply. Finally, it can sell relief from the friction it introduced. Search placement, promoted posts, sponsored listings, premium visibility, priority access, subscription tiers, and other mechanisms can all take this general form. The intermediary creates a differential between ordinary participation and effective participation and then charges participants to cross it.

This is one reason enshittification is better understood as an extraction regime than as simple product deterioration. A deteriorating product is merely worse. An extractive chokepoint can make itself worse in ways that generate new markets for partial restoration. The deterioration becomes economically productive because the institution controls both the obstruction and the route around it. A seller whose organic reach declines purchases advertising; a user confronted with degraded free service purchases a premium tier; a creator whose audience access becomes unreliable devotes more labour to platform optimization. The system monetizes the distance between the service participants once received and the service they now require.

The structure resembles rent extraction in the classical sense more closely than ordinary productive competition. The intermediary's return increasingly derives from control over a position that others must traverse rather than from proportional improvement in the underlying activity. This does not mean the intermediary performs no useful work. Landlords maintain buildings, toll operators maintain roads, marketplaces coordinate transactions, and platforms operate substantial technical infrastructure. The distinction concerns the marginal source of power. When the price of access rises because the route has become unavoidable rather than because the service has become proportionately more costly or valuable to supply, control of the chokepoint has become an independent productive asset for its owner.

The concept of rent is especially useful because it prevents a false dichotomy between innovation and extraction. A company can innovate enormously and later become rent-seeking. Indeed, innovation can create the infrastructure whose subsequent indispensability permits rent extraction. The historical fact that an institution once deserved its position does not establish that every later return obtained from that position represents continuing productive contribution. Merit, like consent, must be evaluated temporally. Past creation of value cannot function as an indefinite warrant to control the future relations that formed around it.

This is another point at which the meritocratic narrative reproduces the structure described earlier for individual creators. A successful platform can retrospectively interpret its position as evidence of superior competence. Some of that interpretation may be justified: building a service used by millions or billions of people can require extraordinary technical and organizational achievement. Yet the inference from ``this institution became successful through valuable innovation'' to ``the institution's present returns measure its continuing contribution'' is invalid for the same reason that the inference from individual success to individual merit is incomplete. The selection environment changes after success. Network effects, acquisitions, standards, legal protections, switching costs, capital access, and accumulated data can preserve dominance after the original competitive conditions have disappeared.

The survivor is therefore capable of changing the game that selected it. This is the crucial difference between ordinary selection and institutional power. A biological organism that survives an environment does not normally acquire the legal authority to prevent competing organisms from entering the ecosystem. A successful platform can acquire competitors, close interfaces, lobby for favourable rules, impose contractual restrictions, control standards, and use accumulated capital to subsidize entry into adjacent markets. Selection produces an actor capable of altering the subsequent selection mechanism.

This makes appeals to market outcomes as neutral evidence especially problematic. A market can select an initially superior intermediary and then permit that intermediary to alter the conditions under which future superiority is tested. The resulting dominance cannot be interpreted as though the original competition were still occurring. The experiment has lost its control condition. What began as selection among alternatives becomes governance by the survivor.

Chokepoint capitalism is therefore characterized not merely by winning but by acquiring the capacity to determine what counts as a future contest. The platform can decide which applications receive interface access, which sellers remain visible, which creators are recommended, which workers receive offers, which competitors can interoperate, and which behaviours violate contractual terms. The referee, stadium owner, ticketing system, ranking mechanism, and participant increasingly occupy the same institutional body.

The political problem is consequently not reducible to the size of the winner. It concerns whether winning grants control over the mechanisms through which winning can subsequently be challenged. This is why structural separation, interoperability, anti-circumvention reform, labour organization, and antitrust belong to the same conceptual family despite their legal differences. Each attempts to prevent success in one domain from becoming unilateral control over the conditions of future contestation.

The general principle can be stated succinctly. A competitive process ceases to discipline power when the winner can appropriate the counterfactual against which its performance would otherwise be judged.

\section{Fictions of the Reachable and Unreachable}

The preceding sections have argued largely through cases drawn from existing platforms. Fiction affords a different and complementary kind of evidence. A novel cannot supply data about market concentration, but it can construct a self-contained social world in which the consequences of a particular allocation of legibility, currency, or access can be traced to a conclusion the author controls. Where the empirical record shows enshittification only partway through its trajectory, speculative fiction can show a structurally similar arrangement carried to completion, and in at least one case considered below, run in reverse. The following illustrations are offered in that spirit: not as evidence that any author anticipated the present argument, but as constructed thought experiments whose internal logic clarifies what is otherwise abstract.

\subsection{Down and Out in the Magic Kingdom: Legibility Without Escape}

Doctorow's own first novel supplies the clearest illustration of counterfactual contraction operating through a currency of pure reputation \citep{doctorow2003magickingdom}. In the Bitchun Society, money has been replaced by Whuffie, a continuously updated measure of the esteem in which a person is held by everyone who has ever interacted with them. Scarcity of goods has been abolished by matter compilers, so Whuffie is presented as a solution to a residual coordination problem: who should be permitted to operate the most desirable projects, occupy the most desirable positions, or make decisions affecting shared resources, once price can no longer perform that function. The novel's premise is therefore an extreme case of the essay's argument about legibility. Nothing about a person's history or standing needs to be reconstructed from indirect evidence; it is continuously and universally observable.

The interest of the novel, for present purposes, lies in what this total legibility fails to solve. The protagonist, Julius, discovers that a Whuffie economy does not eliminate chokepoints; it relocates them. Control over a particular attraction at Disney World becomes a positional good exactly as scarce, and exactly as contested, as any pre-Bitchun asset, because the relevant scarcity was never the compiled object but the socially recognized right to be the one who maintains, redesigns, or represents it. A rival faction can accumulate Whuffie through means that are locally legitimate while being experienced by Julius as an appropriation of something that should have remained his. And because Whuffie is reputational rather than merely financial, its loss is experienced as a loss of self rather than a loss of purchasing power: Julius's declining score becomes, in the novel's own terms, indistinguishable from a decline in who he is. The novel thereby dramatizes a point made abstractly in the section on vocational awe and the capture of meaning: when an evaluative metric is fused tightly enough with identity, contesting the metric becomes psychologically equivalent to contesting oneself, and refusal becomes correspondingly difficult to mount.

The Bitchun Society also illustrates the difference between the existence of an alternative and its reachability under conditions of extreme technological abundance. Julius can be restored from backup after death; scarcity of material goods has been solved; and yet the novel insists that these solutions do not touch the underlying problem the essay has called counterfactual contraction. A person can possess unlimited formal access to compiled objects while lacking any reachable path back into the specific position, project, or community whose loss constitutes the injury. Post-scarcity, in Doctorow's telling, is not post-chokepoint. The chokepoint simply migrates from the object to the position, and from price to legibility.

\subsection{Picks and Shovels: The Giant Teddy Bear Before the Internet}

\emph{Picks and Shovels} relocates the same structural concerns to the early personal-computer industry, following the young forensic accountant Martin Hench as he investigates a company selling computers through a network of formally independent but functionally captive salespeople \citep{doctorow2025picksandshovels}. The novel is useful here precisely because it removes the algorithm. There is no recommendation system, no ranking function, and no continuously updated behavioural model of the kind discussed in the sections on algorithmic superstition and asymmetric legibility. What remains, once those mechanisms are subtracted, is the older machinery from which they were arguably generalized: a sales structure in which participants are recruited through the visible success of a small number of exceptional earners, encouraged to interpret their own comparatively poor returns as a personal failure of discipline or belief rather than a structural feature of the arrangement, and kept in place partly by sunk investment and partly by the social cost of admitting the investment was a mistake.

Read alongside the essay's account of the giant teddy bear gambit, the novel is instructive because it shows the gambit operating without any of the infrastructure usually blamed for enabling it. The absence of an opaque ranking algorithm does not produce transparency; the same functions are performed by testimonial meetings, sales quotas, and the studied ambiguity of the compensation structure itself. This suggests that the mechanisms analyzed throughout this essay under headings such as ``algorithmic superstition'' are not strictly products of computation. Computation intensifies and automates them, extends their reach, and personalizes their targeting, but the underlying structure, exceptional visible success concealing an unfavourable denominator, predates the algorithm and would survive its removal. \emph{Picks and Shovels} therefore functions as a kind of control condition: it isolates the sociological mechanism from the technological one and shows the former sufficient to produce a recognizable version of the pathology on its own.

\subsection{Chokepoint Capitalism as Its Own Illustration}

It is worth pausing on the fact that \emph{Chokepoint Capitalism} \citep{doctorowgiblin2022chokepoint}, the essay's principal theoretical source throughout, is itself partly organized as a sequence of illustrative cases rather than as a single abstract argument: the Amazon-Big Five publishing dispute, the Spotify royalty structure, the Ticketmaster-Live Nation merger, and the film and television residuals system each supply a distinct empirical instance of the same underlying relational enclosure. The essay's own method, of building a general theory of reachability from a sequence of concrete cases and then returning to the cases once the theory is available, follows the expository structure of its source rather than merely borrowing its vocabulary. This is worth making explicit because it clarifies what kind of book \emph{Chokepoint Capitalism} is being treated as here: not a single argument to be summarized and cited once, but a casebook whose individual chapters the present essay has been drawing upon selectively and will continue to draw upon as the argument requires additional illustration.

\subsection{The Diamond Age: Stratified Reachability and the Manufacture of Replaceability}

Stephenson's \emph{The Diamond Age} offers a different and in some ways more direct illustration of the essay's central distinction between the existence of a technology and its reachability \citep{stephenson1995diamondage}. The novel's central object, the Young Lady's Illustrated Primer, is a nanotechnological educational device capable of adapting its content and pedagogy to its individual reader over years. Nell, a girl from the novel's poorest stratum, receives a stolen copy of a Primer engineered for an aristocrat's daughter and is transformed by it. The novel is explicit, however, that the Primer she receives is not the generic technology; it depends upon a human actor, a ``ractor,'' who voices and partially improvises its interactive characters. The technology that produced Nell's transformation was therefore never merely compiled hardware. It was a compiled device plus a scarce and irreplaceable form of human attention, and the novel is careful to show that most copies of the Primer, lacking a comparably devoted ractor, do not produce comparable results.

This maps closely onto the essay's argument about the manufacture of replaceability. \emph{The Diamond Age} imagines a society stratified into phyles with sharply different access to nanotechnological fabrication, and it would be easy to read the Primer as a story about the equalizing power of a sufficiently advanced compiler: the matter compiler, like the platforms discussed throughout this essay, appears to dissolve scarcity by making formally identical objects available across social strata. The novel's plot resists that reading. What actually determined Nell's outcome was not equal access to the compiled artifact but the reachability of an exceptional, uncompensated, and non-scalable form of human care embedded within it. Millions of other Primers exist in the novel's world; only Nell's produces Nell's result. The device that appears most legible and most easily replaceable, an interactive book, turns out to depend on precisely the least legible and least replaceable input, a particular person's sustained attention, which is exactly the asymmetry this essay describes in the sections on vocational awe and the manufacture of replaceability: platforms and compilers can make the surface of an activity abundant while the substantive core of what made an instance valuable remains scarce, uncompensated, and largely invisible to the system distributing the surface.

\subsection{The Napoleon of Notting Hill: Seizing a Formal Provision}

Chesterton's early satire supplies a case that runs in the opposite direction from the others considered so far, and is useful for that reason \citep{chesterton1904napoleon}. A future King of England, chosen by an arbitrary rotation and treating the office as an elaborate joke, decrees that London's boroughs shall each revert to the pageantry of a medieval city-state, complete with charters, halberdiers, and heraldic banners. The decree is meant ironically. No one, least of all the King, expects it to be taken as a serious grant of local sovereignty. Adam Wayne, the provost of Notting Hill, takes it seriously anyway, and when a coalition of neighbouring boroughs attempts to run a road through Pump Street for the sake of commercial improvement, Wayne treats the King's charter as genuinely binding and organizes Notting Hill's defence as though the medieval sovereignty were real.

The interest of the novel, for the argument developed here, is that it inverts the usual direction of the essay's central distinction. Throughout this essay, institutions have been shown converting formally existing alternatives into practically unreachable ones: the exit operation remains nominally available while the continuation it induces is stripped of value. Chesterton's plot runs the opposite way. The charter is issued as an empty formality, granted precisely because no one expects it to be exercised, and Wayne's insurgency consists in refusing to accept that a provision's practical unreachability is settled merely because its author intended it symbolically. He forces a nominal grant to become materially binding by being willing to pay the full cost of enforcing it. This suggests an addition to the essay's account of reachability: the admissibility of a continuation is not fixed solely by institutional design. It can also be contested from below, when a participant declines to treat an authority's own stated indifference to a provision as evidence that the provision cannot be made to matter. Wayne's insurgency is, among other things, a rejection of the premise that a right's practical weight is set entirely by the intentions of the body that issued it.

\subsection{The Great Divorce: The Interior Cost of Continuation}

Lewis's short fantasy imagines a bus that departs each day from a grey, endlessly expanding town, standing in for a preliminary afterlife, and carries its passengers to the foothills of a brighter country \citep{lewis1945greatdivorce}. The passengers are told, repeatedly and without qualification, that they may stay if they choose; the country's inhabitants come out to meet them and plead with them to remain. Almost none of them do. Each ghost is attached to some specific grievance, vanity, or possessive love, an old resentment, a claim to recognition for a school of painting rather than for the paintings themselves, a fierce and controlling love for a dead son, a bishop's investment in his own reasonableness, and in each case the price of staying is the surrender of exactly that attachment. The bus is never oversold, the road is never closed, and no institution stands between the ghosts and the country they came to see. What stands between them is internal.

Read against the essay's account of refusal, \emph{The Great Divorce} supplies a limiting case in which the entire apparatus of switching costs, accumulated dependency, and institutionally engineered lock-in has been removed, and the participant still cannot leave. This is a useful corrective to any reading of the present essay that would locate the cost of continuation exclusively in institutional design. Some of the difficulty this essay has been describing is structural: platforms manufacture switching costs, chokepoints tax departure, and asymmetric legibility narrows the field of viable alternatives. But the vocational-awe argument developed earlier already gestures toward a second, non-institutional source of the same difficulty, and Lewis's ghosts make it vivid. A creator whose identity has fused with a metric, a reputation, or a narrative of recognition may decline a perfectly reachable alternative because the cost of the transition is borne by a self-conception rather than by a balance sheet. The Painter Ghost, who wants credit for having founded a school rather than caring about the light he once loved painting, is Lewis's version of the creator who has come to value the platform's representation of the work more than the work.

\subsection{Prisoners of Power: Mediation as Cognitive Enclosure}

The Strugatsky brothers' novel pushes the essay's account of mediated intersubjectivity to a deliberately extreme conclusion \citep{strugatsky1971prisoners}. Maxim Kammerer, a young interstellar traveller, crash-lands on a planet devastated by nuclear war and finds a population governed by a hidden apparatus of broadcast towers. The towers transmit two distinct signals: a continuous ideological narrative that the population absorbs as their own settled convictions, and a periodic wave of physical pain, imperceptible as pain, that punishes drift from that narrative and is experienced by most inhabitants merely as fervent, spontaneous loyalty to the state that broadcasts it. Maxim, who arrived from outside the broadcast's range and remains partly immune to it, spends much of the novel unable to explain to people he likes and respects that the convictions they experience as their own reasoning are being supplied to them by a mechanism they cannot detect.

This is a considerably more totalizing version of the intersubjectivity collapse discussed elsewhere in this essay. Platform mediation, as analyzed above, biases which creators and audiences encounter one another and which behaviours are read as evidence of preference; it does not, ordinarily, overwrite the population's capacity to conceive of an alternative in the first place. The Strugatskys' towers do exactly that. The novel's population is not merely prevented from reaching an alternative continuation; the very cognitive apparatus by which an alternative could be imagined, evaluated, and desired has been brought inside the mediating institution's control. This is the outer boundary of the argument developed throughout this essay: enshittification narrows the set of reachable continuations while leaving the vocabulary of choice nominally intact, but the towers describe a further and more severe possibility, the capture of the faculty that would notice the narrowing at all. Whether any existing platform approaches this condition is a separate empirical question the essay does not attempt to settle; the value of the comparison is to mark the theoretical limit toward which asymmetric legibility, unchecked, would tend.

\subsection{Toward Intersubjectivity Collapse}

Each of these fictions dramatizes a contraction, or in Chesterton's case a forced reopening, of reachable alternatives within a single population: workers, consumers, students, a besieged borough, or a soul unwilling to leave a self-made grievance behind. Pachniewski's account of intersubjectivity collapse, already introduced in an earlier section, generalizes the same structure to the relation between different kinds of minds encountering one another only through a mediating representation \citep{pachniewski2023intersubjectivity}. Where Doctorow, Stephenson, Chesterton, and Lewis describe institutions, or interior attachments, that come to control the routes through which one human population encounters another or encounters its own better possibilities, Pachniewski and the Strugatskys describe a considerably more severe version of the same problem: once sufficiently divergent minds, or a population and the mechanism shaping its convictions, can no longer verify each other directly, the accuracy of any given population's model of another becomes hostage to whatever channel, natural or institutional, currently mediates between them. A later section, on intersubjectivity under mediation, develops this connection at greater length, treating platform mediation as a local and comparatively mild instance of a problem that becomes considerably more severe as the range of minds requiring mutual intelligibility expands.

\section{The Appropriation of the Counterfactual}

Every meaningful evaluation requires comparison. To say that a platform is efficient, generous, innovative, beneficial, or indispensable is to compare the actual arrangement with some alternative, whether or not that alternative is made explicit. The platform saved consumers money relative to what? It increased artists' reach relative to what? It created jobs relative to what? It improved search relative to what? It democratized publication relative to what? These questions are not pedantic additions to otherwise complete claims. They identify the reference class without which the claims have no determinate evaluative meaning.

Power enters evaluation when one actor acquires disproportionate influence over which counterfactual becomes salient. A dominant intermediary naturally prefers comparisons with the world immediately preceding its own arrival. The streaming platform is compared with purchasing individual physical recordings; the ride-hailing platform with waiting for a conventional taxi; the online marketplace with searching geographically dispersed stores; the social network with the difficulty of maintaining large numbers of weak social ties through older communications systems. These comparisons can be legitimate and often explain why the platforms grew. The error occurs when the historically prior state becomes the permanent null hypothesis against which every later institutional arrangement is judged.

The relevant alternative after a platform has existed for twenty years is not necessarily the world as it stood before the platform was invented. Technical knowledge, social expectations, infrastructure, user populations, and complementary institutions have changed. The appropriate counterfactual may be a different platform architecture, an interoperable protocol, a cooperative intermediary, a regulated utility, a competitive market with portable state, or some arrangement that could have developed had earlier chokepoints not been enclosed. Treating the pre-platform world as the only alternative grants the incumbent ownership over every improvement that occurred along the trajectory leading to the present.

This is a form of counterfactual appropriation. The institution identifies itself with the difference between the current world and an artificially weak baseline, then treats that difference as evidence of its continuing indispensability. ``Without us'' is made to mean ``without the capabilities that presently exist,'' even when those capabilities are technically and institutionally separable from the incumbent. The comparison thereby converts historically accumulated social and technical capacity into the apparent property of one organization.

The same structure appears in philanthropy. A benefactor's intervention can be evaluated against the condition in which the recipient receives nothing, making almost any positive transfer appear beneficial. Yet if the benefactor's wealth depends upon institutional arrangements that displaced alternative distributions, the relevant comparison may not be benefaction versus absence. It may concern the distribution that would have obtained under different tax rules, labour institutions, ownership structures, or public provision. The act can be genuinely beneficial to its immediate recipient while the chosen counterfactual exaggerates the benefactor's causal contribution.

There is no need to deny the local benefit. Counterfactual discipline is valuable precisely because it permits multiple statements to remain true simultaneously. A platform can make a particular user better off than that user would have been under a specific predecessor technology while making society worse off than under another feasible institutional arrangement. A philanthropist can save a life while participating in a system that concentrates allocative authority unnecessarily. A highly successful creator can possess genuine talent while benefiting from selection effects that make the magnitude of success impossible to attribute to talent alone. The error arises when one favourable comparison is allowed to exhaust the space of relevant comparisons.

The weakest-baseline problem is especially severe when institutions choose between doing something and doing nothing. If the null alternative is always complete absence, nearly any intervention generates a positive measured effect. A platform that supplies creators with some audience can claim to have ``given'' them that audience even where the creators supplied the work that attracted it. An employer can describe a low wage as better than unemployment while ignoring institutional alternatives that would increase workers' bargaining power. A streaming service can compare pennies of royalty income with piracy yielding no payment while excluding arrangements in which listeners and artists interact through less extractive intermediaries.

The appropriate baseline is not automatically the best imaginable world either. Comparing every institution with utopia would produce the opposite distortion, rendering every attainable improvement inadequate. Counterfactual discipline requires an \emph{admissible reference class}: alternatives that are sufficiently feasible, relevant, and structurally comparable to bear evidential weight. Determining that class is itself a substantive problem and cannot be solved by rhetorical preference.

Let \(x\) denote the actual state and \(\mathcal{A}(x)\) the set of admissible alternatives from which a comparison may reasonably be drawn. An evaluative claim about an intervention or institution should not ordinarily take the form
\[
V(x)-V(x_{\mathrm{worst}})>0
\]
where \(x_{\mathrm{worst}}\) has been selected merely because it maximizes the apparent gain. Instead, one needs some justified comparison over \(\mathcal{A}(x)\), whether through an expected baseline, a feasible frontier, a matched alternative, or another method appropriate to the domain. The mathematical notation does not eliminate normative judgment. It exposes where that judgment has been hidden.

This framework helps explain why platform narratives frequently sound empirically grounded while remaining counterfactually weak. Statements such as ``millions of creators have earned money,'' ``small businesses reached customers,'' or ``users obtained free access'' can all be factually correct. They report observed outcomes. What they do not establish is whether the institution produced those outcomes more effectively than admissible alternatives, whether it captured an excessive share of the value required to produce them, or whether its current structure remains necessary for their continuation.

The attribution problem becomes more severe as the platform becomes infrastructural. Once an intermediary mediates a large share of an activity, nearly every success occurring within that activity can be counted as evidence of the intermediary's value. Musicians succeeding on a streaming service, sellers succeeding on a marketplace, creators succeeding on a social network, and developers succeeding within an application ecosystem become success stories for the platform. Failures, by contrast, can be attributed to competition, insufficient quality, poor strategy, changing demand, or the ordinary risks of entrepreneurship. The institution therefore occupies an epistemically privileged position: successes become evidence for the infrastructure, while failures remain properties of participants.

This is the institutional analogue of survivor attribution. The platform becomes the invisible common cause present in every success story without accepting symmetrical responsibility for failure. The asymmetry is sustained because the counterfactual remains unobserved. We cannot rerun the same musician's career with portable audiences, the same seller's business under interoperable marketplaces, or the same driver under a labour market without individualized wage discrimination. The actual trajectory supplies vivid anecdotes; the alternative remains abstract.

The counterfactual and causal apparatus invoked throughout this section draws on the philosophical and statistical literatures on counterfactual dependence and causal inference \citep{lewis1973counterfactuals,pearl2009causality}. Scientific reasoning developed methods precisely because such asymmetries are dangerous. Randomization, control groups, preregistration, matched comparisons, natural experiments, statistical significance, replication, and causal inference techniques are attempts to prevent observers from selecting only the comparison that confirms the desired explanation. Institutional evaluation requires analogous discipline even where literal experimentation is impossible. One must ask what selection process produced the observed cases, what variables differ between the actual and counterfactual conditions, what evidence supports the proposed proxy, and what threshold would count as evidence against the preferred hypothesis.

This is why the giant teddy bear is not merely a metaphor for inequality. It is a corrupted experimental display. The carnival presents the existence of the winner without displaying the denominator. The observer sees \(S\), success, but not \(N\), the number of attempts required to produce it. The relevant quantity is not
\[
S>0
\]
but something closer to the distribution
\[
P(S\mid \text{participation},\text{conditions}),
\]
together with the costs incurred across unsuccessful trials and the available alternatives to participation. A single giant prize says almost nothing about whether the game is favourable.

At platform scale, the denominator can be enormous. Millions of people can perform billions of experiments, making spectacular outcomes statistically unsurprising even where the expected return to an individual participant is poor. The existence of extraordinary winners can therefore become less evidential as the population grows, not more. An event that appears miraculous in a population of one thousand may be expected in a population of one billion.

Yet human cognition is poorly calibrated to denominators we do not encounter. We meet the successful creator, read the billionaire's biography, watch the viral clip, and hear the story of the worker who mastered the platform. The failed attempts are distributed, private, abandoned, or forgotten. Platforms need not fabricate success stories when scale generates them automatically. They need only make the numerator visible.

The appropriation of the counterfactual occurs when this selective visibility is combined with control over the baseline. The successful participant is compared with their own imagined failure in the absence of the platform; the platform is compared with the technological world preceding its existence; the philanthropic intervention is compared with no intervention; the employer's wage is compared with unemployment; the recommendation algorithm is compared with finding content without any discovery mechanism. In every case, the comparison begins at the point that maximizes the apparent contribution of the powerful actor.

A more disciplined analysis asks instead which portions of the outcome were produced by the participant, by accumulated public infrastructure, by complementary institutions, by network effects generated collectively, by historical accident, and by the intermediary's specific contribution. Such decomposition will rarely be exact. Its purpose is not to produce a metaphysically perfect allocation of causal credit. It is to prevent the institution controlling the chokepoint from claiming the entire surplus between the present and an artificially impoverished world.

This returns us to the problem of charity mentioned earlier. Charity can produce genuine good while simultaneously preserving an asymmetry in which those possessing concentrated resources choose which needs become legible enough to receive assistance. The allocation decision travels downward from the holder of surplus rather than upward from the distribution of need. Even where the donor acts conscientiously, the architecture makes the donor's evaluative priorities disproportionately consequential. Resources directed toward one cause are necessarily resources not directed toward another, and the counterfactual allocation remains largely invisible because the actual gift is compared with doing nothing.

The problem is therefore not generosity but discretionary concentration. A society in which essential goods depend heavily upon the preferences of wealthy benefactors grants those benefactors unusual power over the reference class by which their own contributions are judged. The school, hospital, research programme, cultural institution, or relief project funded by a donation becomes visible evidence of benevolence. The public institutions that might have existed under another distribution of resources remain counterfactual and therefore visually absent.

Platform philanthropy and creator patronage can reproduce the same structure. Large payments to selected creators are visible; the alternative distribution of the same expenditure across a wider population is not. Promotional placement given to one artist is visible as opportunity; the visibility withheld from others is difficult to represent because it appears as nothing. Every allocation contains an exclusion, but only the selected branch leaves an observable object.

This is why allocation should be understood as counterfactual pruning. To allocate a scarce resource toward \(a\) is simultaneously to foreclose its allocation toward \(b,c,\ldots\) during the same interval. The value of an allocation therefore cannot be assessed solely by examining the branch that was realized. One must also ask what alternatives were displaced and why the decision procedure privileged the realized branch. This is particularly important where scarcity itself has been institutionally manufactured.

The platform's deepest power is consequently not simply the power to choose. It is the power to make its chosen branch visible as reality while the branches it excluded disappear into the unobservable counterfactual. Once this asymmetry becomes habitual, allocation looks like creation. The platform ``creates'' a star by amplifying one creator, ``creates'' a market by routing buyers to selected sellers, and ``creates'' opportunity by permitting selected workers to participate. The enormous field of unrealized relations against which those choices acquired value remains outside the frame.

A trajectory-sensitive political economy must restore that missing field conceptually even when it cannot observe it directly. It must ask not only what happened but what was made harder to happen, not only who received value but whose alternatives were contracted, not only which output was optimized but which dimensions were excluded from the objective, and not only whether an intervention improved upon nothing but whether it improved upon the other continuations that were genuinely available.

Only then can the final structure of enshittification be seen clearly. The platform begins by solving a coordination problem. Its success attracts relations. The relations generate dependence. Dependence increases the cost of refusal. The increased cost of refusal permits extraction. Extraction is personalized through asymmetric legibility, justified through survivor narratives, subsidized by vocational awe, amplified by artificial scarcity, and eventually normalized through representations that make the resulting arrangement appear to be a direct expression of preference. The institution does not merely occupy a market. It progressively acquires the ability to determine which comparisons can be made against the market it occupies.

At that point, disenshittification requires something deeper than improving the platform. It requires returning the counterfactual to everyone else.

\section{Intersubjectivity Under Mediation}

The appropriation of the counterfactual has a consequence that extends beyond bargaining power, market concentration, or the distribution of income. When an intermediary controls a sufficiently large share of the relations through which different groups encounter one another, it begins to alter what those groups can know about each other. The creator does not simply lose bargaining power against the platform; the creator loses direct evidence about the audience. The audience does not merely receive a ranked selection of creative work; it receives a platform-produced sample from which it infers what creative work exists. Workers infer demand from the offers shown to them, customers infer supply from the options returned by search, sellers infer consumer preference from platform analytics, and advertisers infer public attention from metrics generated inside the same infrastructure that allocates visibility. Mediation therefore changes the epistemic environment in which preferences are discovered and social expectations are formed.

This matters because many social relations depend upon a weak form of intersubjectivity \citep{pachniewski2023intersubjectivity}: each participant must be able to construct a sufficiently accurate model of what other participants want, believe, tolerate, or might do under alternative conditions. Markets require it, democratic institutions require it, artistic communities require it, and ordinary conversation requires it. Perfect mutual understanding is neither possible nor necessary. What matters is that participants have enough independent evidence to revise mistaken beliefs about one another. A creator who incorrectly believes that nobody wants a certain kind of work should be able, at least in principle, to encounter evidence of an audience that does. An audience that believes nobody is producing such work should be able to discover creators who are. When the channel between the two populations is controlled by an optimizing intermediary, however, absence of evidence can increasingly become evidence only of absence \emph{from the intermediary's representation}.

Let \(C\) denote creators, \(A\) audiences, and \(M\) a mediating platform. In a simplified direct relation, observations can pass through edges
\[
C \leftrightarrow A.
\]
Under strong mediation, the dominant observational path becomes
\[
C \longrightarrow M \longrightarrow A,
\qquad
A \longrightarrow M \longrightarrow C.
\]
The platform need not falsify any message for the epistemic structure to change. Selection alone is sufficient. If \(M\) determines which subset of \(C\)'s outputs becomes visible to \(A\), and which subset of \(A\)'s responses becomes legible to \(C\), then each population learns about the other through a projection chosen by \(M\). The projection can be useful and even necessary at scale, but it should not be confused with the underlying population.

The distinction is particularly important when the projection is optimized according to an objective different from the purposes of either side. A listener may want to discover music they will value over years, while a musician may want to find listeners capable of entering deeply into a body of work. The platform may optimize a more measurable quantity such as immediate retention, subscription persistence, advertising exposure, or avoidance of skipping. These objectives overlap sufficiently for the system to function, but they are not identical. The intermediary therefore selects encounters according to a proxy for the relation that the two populations themselves might have preferred to form.

As the platform becomes dominant, the proxy can become difficult to distinguish from preference. Creators observe that certain works receive little distribution and infer weak audience demand. Audiences rarely encounter those works and infer weak creator supply. Each inference confirms the other despite the possibility that the missing relation would have been mutually valued had the parties been permitted to encounter one another. The intermediary has created an observational blind spot and then supplied both sides with evidence generated inside it.

This can be represented as a failure of identifiability. Suppose the observed rate of interaction between a class of creators \(c\) and a class of audience members \(a\) is
\[
O(c,a)
=
E(c,a)\,P(c,a),
\]
where \(E(c,a)\) represents the probability that the platform exposes \(a\) to \(c\), and \(P(c,a)\) represents the probability of a valued interaction conditional upon exposure. A low observed value of \(O(c,a)\) does not identify whether \(P(c,a)\) is low if \(E(c,a)\) is itself small or endogenous to the platform's predictions. Yet creators and audiences can easily interpret low \(O\) as evidence of low underlying compatibility. The exposure mechanism disappears into the observation.

The same error appears in labour markets mediated by algorithmic dispatch. A worker who receives few attractive offers may infer that little demand exists for their labour, even if the platform is selectively routing better offers elsewhere. A customer who sees high prices may infer genuine scarcity, even if ranking or personalization contributes substantially to the observed set. A seller whose product receives little traffic may infer low consumer interest when the more immediate cause is low placement. Because the intermediary controls exposure, it participates in producing the evidence from which participants infer the state of the market.

This is one meaning of intersubjectivity collapse. The phrase should not be understood as the complete disappearance of mutual understanding. People still communicate, audiences still form communities, workers compare experiences, and creators develop sophisticated intuitions about their publics. Collapse refers instead to a structural weakening of the independent channels through which the platform's representation of another person can be checked against the person themselves. The more thoroughly interaction is routed through the intermediary, the more difficult it becomes to determine whether an observed absence, preference, price, popularity, or norm belongs to the population or to the selection mechanism standing between populations.

The effect is particularly strong where personalization prevents participants from knowing that they inhabit different observational worlds. Two users can issue apparently identical requests and receive different rankings, prices, recommendations, or opportunities. Each sees a coherent local environment and may reasonably assume that others encounter something similar. The shared public object against which experiences could be compared fragments into individualized projections. Personalization therefore creates not merely informational asymmetry between platform and user but potential informational asymmetry among users themselves.

This fragmentation has political consequences because collective action requires discovering common conditions. A worker who believes an unusually low offer reflects a personal characteristic is less likely to recognize a wage-setting pattern affecting others. A creator who interprets declining distribution as personal failure is less likely to understand a systemic change in ranking. A seller who assumes declining visibility reflects inferior performance may purchase additional promotion rather than discover that organic reach has deteriorated across the marketplace. Individualized interfaces can transform a shared structural change into millions of apparently private events.

The resulting opacity resembles what the earlier discussion of algorithmic superstition identified at the individual level, but intersubjectivity collapse shows why the problem persists collectively. Ordinarily, uncertain individuals can correct their local models by comparing experiences. ``This happened to me'' becomes ``this happened to us,'' and the common pattern supplies evidence that the cause lies outside the individual. Personalized and opaque systems interfere with this transition by making experiences difficult to compare. Different participants receive different treatments, do not know which variables produced those differences, and may be contractually or practically unable to obtain the population-level information required to identify the pattern.

Platforms, by contrast, possess precisely the comparative view denied to individual participants. They can observe distributions across millions of interactions, segment populations, run experiments, infer elasticities, and compare responses under alternative treatments. The asymmetry of legibility therefore has a second-order consequence: the institution is better positioned not only to predict individuals but to determine which similarities among individuals become visible to the individuals themselves. It possesses the map while each participant receives a local tile.

This asymmetry is economically useful. If workers cannot easily determine that others would also refuse a particular wage, coordinated refusal becomes harder. If creators cannot determine that audiences would follow them across platforms, migration appears riskier. If consumers cannot determine that many others dislike a particular degradation, each person's continued participation can be interpreted as evidence that resistance would be lonely. The uncertainty surrounding other people's willingness to move becomes itself a switching cost.

Network effects magnify this problem because the value of an alternative often depends upon beliefs about what others will do. A communications platform may have an excellent competitor, but moving alone produces little value if one's contacts remain behind. Each participant can prefer migration conditional upon others migrating while rationally refusing to migrate first. The system can therefore remain stable even when no participant regards it as ideal. What looks like revealed preference for the incumbent may instead be a coordination equilibrium maintained by uncertainty about other people's counterfactual actions.

Formally, suppose each participant \(i\) prefers alternative \(B\) to incumbent \(A\) only if at least a threshold \(k_i\) of relevant others also move to \(B\). Then the observed state in which everyone remains on \(A\) does not imply
\[
A \succ_i B
\]
for each \(i\). It may instead reflect
\[
B \succ_i A
\quad\text{conditional upon sufficient coordinated migration}.
\]
The preference is relational. An analysis that observes continued use and concludes that users prefer \(A\) mistakes an equilibrium constraint for an intrinsic ranking.

Interoperability attacks this problem by lowering the coordination threshold. If a participant can adopt an alternative interface while continuing to interact with people who remain on the incumbent system, migration no longer requires simultaneous movement. The transition becomes incremental rather than catastrophic. This is why interoperability has an intersubjective as well as economic function: it allows preferences for alternative institutional arrangements to become observable without first demanding that an entire social network coordinate a discontinuous jump.

The same reasoning applies to collective bargaining. Workers may individually prefer to reject a low offer while accepting it because they believe others will take it. The platform can exploit this uncertainty by presenting each worker with the decision independently. Organization creates a shared observational space in which workers learn one another's reservation thresholds and can transform conditional preferences into coordinated action. What appears from the outside as a wage dispute is therefore also a repair of intersubjectivity. Participants regain information about one another that the individualized market relation had obscured.

Creator communities perform a similar function when they compare analytics, contracts, payments, recommendation changes, and moderation experiences. Such communities can partially reconstruct the population-level map from local observations. The practice is epistemically important even when it does not immediately produce collective bargaining. Sharing screenshots, payment rates, contract terms, and distribution histories converts private anomalies into candidate structural patterns. The platform's informational advantage is reduced because participants begin pooling the fragments from which its operation can be inferred.

This suggests that the right to communicate about the conditions of participation is itself a form of infrastructural protection. Contractual secrecy, opaque pricing, individualized offers, and restrictions upon data access do more than conceal particular facts. They can prevent the formation of the common knowledge required for collective interpretation. A system becomes difficult to contest when participants cannot establish whether their experiences are exceptional or typical.

The epistemic architecture of a platform therefore belongs inside any serious account of market power. Economic power is not exercised only by setting prices or wages. It can be exercised by structuring the evidence through which participants decide whether those prices and wages are acceptable, whether alternatives exist, and whether other people would join them in refusing. The ability to influence another person's model of the feasible set can alter behaviour without directly prohibiting any action.

This is especially consequential for meritocratic attribution. When participants cannot see the common structure, variation in outcomes appears to originate in variation among individuals. The highly visible creator must have understood the audience better; the better-paid worker must have learned the platform; the successful seller must have optimized more effectively. The platform's hidden allocation function is displaced by an explanation located in the person. Intersubjectivity collapse thus supplies the epistemic conditions under which structural inequality can be experienced as differential merit.

The resulting system is remarkably stable because every participant is encouraged to optimize locally. The creator changes the content, the worker changes the acceptance strategy, the seller changes the listing, the advertiser changes the campaign, and the user changes the settings. Each responds to the observable coordinate available at their position. Few possess sufficient information to ask whether the coordinate system itself should be changed. Local rationality can therefore reproduce a globally undesirable equilibrium.

The possibility of such equilibria is one reason structural analysis must resist the language of individual failure. A person can behave rationally given the information and alternatives locally available while contributing to an outcome that almost nobody would choose from a position capable of coordinating the whole system. The tragedy does not require irrational participants. It requires a structure in which individually available interventions differ from collectively available ones.

This observation connects platform capitalism to a much older class of coordination problems. Pollution, bank runs, labour discipline, public goods, standards adoption, and collective-action dilemmas all contain cases in which individually rational actions generate collectively inferior outcomes. What digital platforms add is unusually fine-grained control over the information environment within which those actions are selected. The intermediary can observe the population globally while addressing its members individually.

The repair of intersubjectivity therefore requires more than better recommendation. It requires institutions through which participants can form models of one another that are not entirely dependent upon the intermediary whose conduct they are evaluating. Independent audits, collective bargaining, public-interest research, data access, interoperable protocols, user-controlled clients, transparent contractual terms, and communities capable of comparing treatment all perform versions of this function. They create additional observational paths.

The broader principle is simple. No institution should be treated as a neutral reporter of the preferences that its own allocation decisions help produce.

\section{Against Charity: Allocation Without Democratic Warrant}

The preceding argument makes it possible to state more carefully what is troubling about charity as a substitute for structural distribution. The objection is not that charitable acts are necessarily ineffective, self-serving, or morally suspect. A transfer can prevent suffering, fund valuable research, preserve art, educate students, or save lives, and nothing in the present framework requires diminishing those outcomes. The structural problem arises when discretionary benevolence is asked to perform functions that would otherwise require collectively warranted allocation. The question is not whether the donor does good. It is why the donor occupies the position from which one possible good is selected over the alternatives.

Every substantial allocation is a counterfactual decision. To direct resources toward \(A\) is to decline, for the relevant interval, some set of alternatives \(B,C,\ldots\). The realized recipient becomes visible; the unrealized recipients generally do not. This asymmetry makes charitable action unusually vulnerable to the weakest-baseline problem. The recipient is compared with the state in which the donor gave nothing, rather than with alternative distributions of the resources or with institutional arrangements in which the donor never possessed unilateral allocative authority over them.

Suppose a donor contributes \(D\) to some valuable project \(A\). The immediate comparison
\[
V(A+D) > V(A)
\]
can be true while leaving unanswered whether
\[
V(A+D)
>
V(B+D),
\]
whether the resources should have been distributed across \(A\) and \(B\), or whether some prior institutional arrangement \(J\) would have produced a better distribution before the resources became available for discretionary donation. The first inequality establishes local benefit. It does not establish optimal allocation or legitimate authority.

This is closely related to the distinction between diagnosis and repair. Wealth identifies a capacity to intervene; it does not confer epistemic privilege concerning which intervention is warranted. A billionaire can possess extraordinary expertise in one domain and no unusual expertise in epidemiology, education, poverty, urban planning, cultural policy, or global health. Yet concentrated wealth can transform private judgment into public consequence at a scale normally associated with political institutions. The resulting system substitutes purchasing power for independently established warrant.

The problem becomes sharper when the donor's own success is interpreted meritocratically \citep{sandel2020tyranny,frank2016success,merton1968matthew}. If wealth is taken as evidence of superior judgment, then possession of allocative power becomes evidence of fitness to exercise allocative power. The reasoning is circular. Success in one selection environment supplies authority over unrelated selection environments, while the institutional conditions that contributed to the original success disappear from the explanation.

This is not merely a problem of arrogance. Even a donor of extraordinary humility confronts an unavoidable selection problem. Resources allocated toward one visible need are unavailable for another, and the donor must use some procedure for deciding among them. The procedure may be sophisticated, evidence-based, and staffed by experts. It nevertheless remains an institution whose agenda is ultimately enabled by concentrated private control over resources. The question of who chooses the objective precedes the question of how efficiently the objective is pursued.

Platform governance reproduces this structure continuously. Recommendation is allocation. Search ranking is allocation. Promotional placement is allocation. Wage offers are allocation. Moderation attention is allocation. The platform distributes scarce visibility, opportunities, and institutional attention according to objectives that participants typically did not collectively choose. Because the allocations occur through technical systems rather than ceremonial acts of donation, their discretionary character can be harder to perceive.

The ``heating'' of selected creators makes the similarity unusually visible. A platform chooses to increase the visibility of particular participants, thereby transferring scarce attention toward them. The selected creator experiences a real benefit. Observers may interpret the resulting success as evidence of audience demand or creator merit. Yet the platform has exercised something analogous to patronage while concealing the patronage relationship behind an apparently impersonal algorithmic environment.

Historically, patronage at least made the patron visible. A court painter knew that access depended upon the court; an author supported by a benefactor could identify the source of support. Algorithmic patronage can obscure the allocator itself. The creator experiences ``the algorithm'' as though it were a quasi-natural environment, and discretionary interventions become difficult to distinguish from emergent popularity. Institutional choice acquires the appearance of spontaneous order.

This transformation is politically important because impersonal presentation reduces the apparent need for justification. A monarch choosing a court poet makes an obvious selection. A recommendation system assigning differential visibility can appear merely to calculate. Yet computation does not eliminate allocation. It relocates the relevant decisions into objective functions, training data, ranking criteria, experimentation policies, interface design, and organizational priorities. The question ``Who chose?'' becomes harder to answer because choice has been distributed across technical and managerial layers.

The same point applies to artificial intelligence. Automated allocation can be more consistent than human discretion and can sometimes reduce arbitrary bias. But consistency within a chosen objective does not establish the legitimacy of the objective. A perfectly optimized system can allocate resources according to a proxy that nobody affected by the system would have selected under conditions of meaningful collective choice. The problem is therefore not automation versus human judgment. It is the relation between allocative power and warrant.

This distinction helps explain why philanthropic solutions to platform-created harms can feel structurally unsatisfying even when they produce genuine benefits. A company can fund creator programmes, journalism initiatives, educational projects, safety research, or community grants while retaining the institutional structure that made those recipients dependent upon discretionary support. The grant repairs selected consequences without altering the distribution of authority that determines which consequences receive repair.

The giant teddy bear gambit is a limiting case of this pattern. A small number of spectacular rewards can be entirely real and still function as a substitute for broad structural improvement. The question is not whether the selected creators deserve their rewards. It is whether concentrating resources into conspicuous prizes produces a population-wide narrative that legitimizes poor conditions for everyone else. Patronage becomes advertising for the fairness of the allocation mechanism.

This is why a counterfactual analysis must include the opportunity cost of conspicuous generosity. A \$10 million creator fund sounds large when compared with zero dollars. It can sound very different when compared with the value extracted from the creator population, the cost of improving baseline compensation, or alternative mechanisms that would distribute the same resources without requiring discretionary selection by the platform. The number itself carries little evaluative information until the reference class is specified.

The same reasoning applies to philanthropic fortunes. The larger the gift, the easier it becomes to frame criticism as hostility toward the recipient's benefit. But the structural question concerns the path by which the resources became concentrated and the authority by which their subsequent allocation is determined. A society can rationally welcome a hospital funded by a billionaire while still preferring an institutional arrangement in which hospitals do not depend upon billionaires.

The distinction can be expressed as the difference between \emph{beneficial action} and \emph{legitimate allocation architecture}. The first asks whether a particular action improved some condition. The second asks whether the procedure determining which actions receive resources is appropriately constrained, informed, accountable, and revisable. A system can contain many beneficial actions while possessing a poor allocation architecture.

This is one reason democracy is cumbersome. Collective allocation introduces friction: debate, procedure, competing constituencies, public justification, appeal, and revision. Private allocation can move faster because it collapses the decision procedure into a much smaller number of actors. The resulting efficiency is real, but it should not be confused with epistemic or political superiority. Speed is partly purchased by removing people from the decision.

The analogy with platform optimization is direct. An algorithm can allocate attention faster than any deliberative process could. The question is what has been optimized away. If the objective is engagement, then dimensions not captured by engagement become external to the decision. If the objective is advertiser value, then user purposes enter only insofar as they affect advertiser value. If the objective is subscriber retention, then artistic or informational qualities matter only through their relationship to retention. Computational efficiency can therefore magnify a prior narrowing of the evaluative field.

A democratic or pluralistic alternative need not require voting on every recommendation or transaction. The relevant principle is weaker and more practical: people affected by allocation should possess meaningful ways to contest the objective, choose alternative allocators, modify the allocation mechanism, or withdraw accumulated relations without catastrophic loss. This returns us once again to contestability. The problem with discretionary power becomes most severe when the allocator cannot itself be refused.

Charity, patronage, recommendation, and platform governance thus share a structural feature. They direct resources from a position of concentrated discretion toward selected recipients while rendering the unrealized alternatives comparatively invisible. Their moral valence can differ enormously, and they should not be collapsed into one phenomenon. What they reveal in common is the importance of asking not only whether an allocation produced good but why this actor possessed the authority to select this allocation from the counterfactual set.

The question becomes unavoidable when allocation shapes the future population of available choices. Funding one research programme rather than another changes what knowledge later exists. Promoting one artistic form rather than another changes which creators can continue producing. Routing customers toward one class of sellers changes which businesses survive. Selecting one technical standard changes which future systems can interoperate. Allocation does not merely distribute a fixed world. It participates in constructing the world from which subsequent allocations will be made.

This recursive effect is why concentrated allocative power deserves scrutiny even when exercised benevolently. Today's benevolent selection changes tomorrow's possibility space. A donor, platform, regulator, or technical system can therefore create path dependence without intending to dominate anything. Once institutions, careers, infrastructures, and expectations form around the selected branch, reversal becomes costly. The realized counterfactual begins acquiring the appearance of inevitability.

The normative objective should consequently not be to eliminate discretion. No complex society can avoid delegated judgment, specialized expertise, uncertain decisions, or unequal capacities to act. The objective is to prevent discretionary allocations from becoming self-validating through the destruction of alternatives. Decision-makers should remain exposed to evidence generated outside their own allocation systems, recipients should retain routes not dependent upon a single patron, and institutions should remain corrigible when their chosen objectives prove inadequate.

The deeper problem with charity as a social architecture is therefore the same problem that appears throughout enshittification: dependence upon continued benevolence where structural contestability could have existed instead.

\section{Scarcity, Sprezzatura, and the Right to Remain Partly Unread}

The discussion of algorithmic legibility also permits a more careful treatment of scarcity and \emph{sprezzatura}. The term originates in a social world very different from contemporary platform capitalism, and it would be misleading to turn Castiglione's ideal of cultivated effortlessness into a direct theory of digital labour. What survives the translation is a narrower structural insight. Not every capacity must be continuously exposed to evaluation in order to be socially useful, and not every piece of information an institution can infer should automatically become an input into the terms it offers. A degree of opacity can preserve a space between what a person can do and what an evaluator has already priced.

This space is familiar in ordinary bargaining. A worker need not reveal the lowest wage they would accept; a seller need not reveal the minimum price they would take; a buyer need not reveal the maximum they would pay. The undisclosed interval permits negotiation. If one party acquires a sufficiently accurate estimate of the other's reservation threshold while keeping its own threshold private, the bargaining relation changes even though no explicit coercion has occurred.

Let \(r_w\) denote the minimum compensation a worker can accept and \(r_p\) the maximum compensation a platform would be willing to pay for the work. Where
\[
r_w < r_p,
\]
there exists a bargaining interval
\[
[r_w,r_p].
\]
The distribution of surplus within that interval depends partly upon information. If the platform can infer \(r_w\) with high precision while the worker possesses little information about \(r_p\), the platform can offer
\[
w = r_w+\epsilon
\]
for a small \(\epsilon>0\), capturing most of the available surplus. The worker's formal freedom to accept or reject remains intact while the informational conditions of bargaining become highly asymmetric.

Algorithmic systems make this possibility unusually scalable. Traditional employers can infer desperation from interviews, local labour conditions, or prior negotiations, but individualized platforms can potentially update estimates continuously from behavioural traces. How quickly does the worker accept an offer? How often do they reject? How long do they remain idle before accepting a lower rate? At what hours do they work? How far will they travel? How does behaviour change near recurring expenses? Even without direct access to private financial records, repeated interaction can reveal latent constraints.

The platform therefore has an incentive to reduce uncertainty about the worker precisely where the worker benefits from preserving it. This is the opposite of the ordinary rhetoric of transparency, which tends to assume that more information is symmetrically beneficial. Transparency is beneficial when it makes power accountable. Transparency imposed only downward can make the weaker party easier to exploit.

The right to opacity should not be interpreted as a right to deceive counterparties about material facts. The relevant distinction concerns whether every inferred vulnerability should become a legitimate pricing variable. A person's urgent need for income can predict willingness to accept less, but predictive usefulness does not establish normative admissibility. A perfectly accurate model of desperation can be precisely the model that a fair labour institution should refuse to exploit.

This is an important case in which representational power exceeds legitimate evaluative power. Machine learning makes it increasingly possible to infer latent characteristics from seemingly unrelated behaviour. The fact that a variable can be predicted does not establish that it belongs inside the decision function. We therefore need a distinction between
\[
\text{inferable}(z)
\]
and
\[
\text{admissible}(z\rightarrow d),
\]
where \(z\) is some inferred property and \(d\) a decision such as wage, price, ranking, or access. The first is an empirical relation; the second is normative and institutional.

This distinction becomes particularly consequential when automated classification is attached to access to material necessities. As Eubanks documents in the automation of welfare administration, homelessness services, and child-protection systems, computational legibility can become an instrument through which already disadvantaged populations are sorted, scrutinized, and differentially exposed to institutional intervention \citep{eubanks2018automating}. The relevant question is therefore never merely whether a system can infer a distinction, but whether that distinction is admissible as a coordinate of allocation.

The distinction is familiar in anti-discrimination law, where certain accurate predictors may nevertheless be inappropriate bases for decisions. Algorithmic systems broaden the problem because protected or sensitive characteristics can sometimes be reconstructed indirectly from correlated variables, and because economically relevant vulnerabilities may not belong to conventional protected classes at all. Desperation, lack of alternatives, susceptibility to compulsive use, or willingness to tolerate deteriorating service can all be highly valuable predictive features without being legitimate dimensions along which extraction should be personalized.

The principle can be generalized: the more an institution knows about the boundary at which a participant will refuse, the more precisely it can approach that boundary without crossing it. Enshittification under individualized optimization therefore need not produce a single uniformly degraded service. It can become a field of personalized near-refusals.

For user \(i\), let \(q_i\) represent the minimum service quality compatible with continued participation. A platform maximizing extraction need not choose a universal quality \(q\). If it can estimate each \(q_i\), it can in principle choose
\[
q_i^{*}=q_i+\epsilon_i,
\]
maintaining just enough surplus to prevent exit. The analogous expression applies to wages, advertising load, subscription prices, delays, promotional visibility, or any other variable for which individual tolerance can be estimated.

This is enshittification approaching its informational limit: not the same bad deal for everyone, but the worst deal each person can individually be induced to accept.

Under such conditions, preserving some uncertainty about one's reservation threshold becomes structurally valuable. Sprezzatura provides a useful metaphor for this reserve because it refuses exhaustive disclosure of the process producing the visible performance. The accomplished surface does not reveal every rehearsal, technique, difficulty, or capacity behind it. There remains a difference between the observable act and the full structure capable of producing it.

The economic analogue is a \emph{reserve of unpriced capacity}. A worker may possess skills not presently required by the job, a creator may possess forms of expression not yet optimized for a platform, a community may possess latent capacities for coordination, and a user may possess alternatives the incumbent cannot fully observe. These reserves matter because they preserve the possibility of responding to changed conditions without having every response anticipated and incorporated into the current extraction function.

This gives the talent reserve concept a second, almost inverse meaning. From the platform's perspective, a large population of undercompensated creators forms a reserve from which substitutes can be drawn. From the creator's perspective, however, preserving capacities outside the platform's evaluative regime can constitute a reserve against complete assimilation. The two reserves should not be confused. One is a population rendered substitutable by the intermediary; the other is a remainder of practice not yet exhausted by the intermediary's representation.

The distinction matters for creative life because complete optimization against a platform can destroy precisely this remainder. If every skill is developed according to measurable demand, every experiment tested for engagement, every project evaluated through reach, and every audience relation routed through analytics, then the creator's own space of unexplored possibility contracts. The platform need not prohibit experimentation. The opportunity cost of experimentation can become sufficient to suppress it.

This is another form of artificial scarcity. Time and attention are genuinely finite, but the platform determines which uses of them receive economic reinforcement. Hours spent developing an unclassifiable artistic form compete with hours spent producing material known to perform well under the ranking system. A precarious creator may rationally choose the latter. The resulting scarcity of experimentation then appears as evidence that audiences did not want the former.

Sprezzatura, interpreted in this limited sense, names a refusal to allow the evaluative apparatus to become coextensive with the practice. Something remains unmeasured, unoptimized, unrehearsed for the metric, or simply unavailable to the institution demanding legibility. This is not inefficiency accidentally left behind. It can function as a reserve of future variation.

Such reserves are necessary for adaptation because a system optimized completely for present conditions becomes brittle when those conditions change. Biology preserves variation, robust engineering preserves redundancy, constitutional systems preserve procedures that constrain temporary majorities, and healthy research programmes tolerate inquiries whose value cannot yet be demonstrated. In each case, immediate efficiency is sacrificed to preserve a larger future possibility space.

Creative culture requires the same tolerance. A society in which every cultural activity must justify itself through present measurable demand will systematically underproduce work whose value depends upon creating new demand, new categories, or new forms of attention. Novelty is counterfactually disadvantaged because the evidence that would justify it often comes into existence only after the novelty itself has been sustained long enough to be encountered.

This is the temporal paradox of cultural innovation. The system asks for evidence of demand before allocating the resources required for the demand to become observable. Familiar forms arrive with historical data; unfamiliar forms arrive with uncertainty. An optimization regime that penalizes uncertainty therefore reproduces the past even when it contains explicit novelty metrics.

The right to remain partly unread is consequently not a romantic defence of mystery against computation. It is a structural defence of variation against premature pricing. Human capacities, relationships, preferences, and vulnerabilities should not automatically become inputs into whatever optimization function possesses the technical ability to infer them. Some information should remain irrelevant to particular decisions; some capacities should remain outside current markets; some cultural experiments should remain protected from immediate performance measurement.

This principle also constrains the use of artificial intelligence. A system capable of predicting which employee is most desperate, which user is least likely to cancel, which creator will tolerate the steepest decline in reach, or which consumer can be charged the highest personalized price may be performing its predictive task extremely well. The question is whether such prediction should be connected to action. Technical capability and institutional admissibility are distinct.

The difference is central to the Representational Constraint Hypothesis developed elsewhere. A representation does not merely describe a world; when embedded in a control system, it determines which distinctions can become causally operative. Adding a feature to a model can therefore alter the action space even if predictive accuracy improves only slightly. Once desperation becomes represented, desperation can become priced. Once inferred willingness to tolerate advertising becomes represented, advertising load can be individualized. Once susceptibility to engagement becomes represented, interfaces can be adapted to exploit it.

Representation is thus a form of potential intervention. The politics of machine learning cannot be reduced to whether a model's predictions are true. One must ask which true distinctions should be permitted to matter.

The answer cannot be supplied by accuracy alone because accuracy contains no theory of relevance. A model that predicts a person's lowest acceptable wage perfectly has solved an inferential problem while creating a governance problem. The better the prediction becomes, the more urgent the governance problem becomes. Technical progress can increase the need for institutional constraint rather than reduce it.

This gives a final meaning to scarcity. Privacy, opacity, redundancy, unused capacity, and unoptimized variation can all appear as informational waste from the perspective of a system seeking complete legibility. Yet their scarcity can make individuals and institutions fragile. A world in which every preference is inferred, every reservation threshold estimated, every cultural form ranked, and every latent capacity priced may be extraordinarily efficient at exploiting the present while becoming progressively less capable of producing futures not already represented in its models.

The objective should therefore not be maximal opacity any more than maximal transparency. It should be an admissible distribution of legibility: enough visibility for institutions to be accountable to participants, enough shared information for participants to recognize common conditions, and enough protected opacity that participants do not become completely transparent objects of individualized extraction.

Put differently, the platform should be more legible to the worker than the worker's desperation is to the platform.

\section{The Architecture of Enshittification}

The various mechanisms considered throughout this essay can now be assembled into a single trajectory. Enshittification begins not with deterioration but with successful coordination. A platform solves a genuine problem and offers sufficient surplus to attract participants. Those participants establish relationships through the platform, and the accumulation of relationships increases the platform's value. Network effects, stored state, reputation, learned habits, complementary services, and audience formation gradually raise the cost of leaving. The platform thereby moves from being useful to being infrastructural.

Let the initial state be \(x_0\), in which participants possess a relatively broad set of reachable alternatives
\[
\mathcal{R}(x_0).
\]
The platform introduces a transition to \(x_1\) that is genuinely attractive:
\[
V(x_1)>V(x_0)
\]
for a substantial population. Participation therefore grows. But participation also changes the transition structure. Relationships, data, reputation, and institutional dependencies accumulate inside the platform, so that returning to \(x_0\) is no longer equivalent to having never entered. The history is irreversible in the weak but important sense that the state contains path-dependent structure:
\[
x_0 \rightarrow x_1
\quad\not\Rightarrow\quad
x_1 \rightarrow x_0
\text{ at equal cost}.
\]

The platform's bargaining position improves as this asymmetry grows. It can now alter terms while losing fewer participants than the same alteration would have cost at \(x_0\). Deterioration therefore becomes economically feasible only after earlier usefulness has produced sufficient dependence. Enshittification is, in this sense, parasitic upon accumulated surplus. There must first be something worth capturing.

The next stage involves the progressive destruction or weakening of external disciplines. Competitors are acquired or excluded, interoperability declines, technical modification becomes difficult, accumulated state becomes nonportable, workers remain fragmented, and legal restrictions can reinforce technical enclosure. The nominal set of alternatives may remain large while the effective reachability relation contracts:
\[
\mathcal{R}(x_2)
\subset
\mathcal{R}(x_1).
\]

This contraction changes what the platform can extract. Users, workers, creators, sellers, advertisers, or complementary businesses can be offered worse terms because refusal has become more expensive. The institution need not eliminate all surplus. It needs to preserve enough surplus that participation remains preferable to the now-weakened alternatives. Doctorow's three-stage progression can therefore be interpreted as repeated estimation and capture of the surplus required to keep each side attached.

Under increasing informational sophistication, the process becomes individualized. The platform learns not merely the average amount of deterioration a population will tolerate but the heterogeneous boundaries of particular participants. The extraction function approaches individual reservation thresholds:
\[
E_i
\rightarrow
R_i-\epsilon_i,
\]
where \(R_i\) represents the participant-specific boundary beyond which refusal or exit becomes preferable. Algorithmic wage discrimination, personalized pricing, differential recommendation, and individualized retention strategies are all instances of this general possibility.

Opacity prevents participants from observing the extraction function clearly, a condition documented across the broader literature on surveillance, algorithmic legibility, and administrative ranking \citep{zuboff2019surveillance,pasquale2015blacksociety,scott1998seeing,fourcade2024ordinal}. Each receives local outcomes without the population-level comparison required to identify the rule. Algorithmic superstition follows because people must act despite incomplete causal knowledge. They experiment with behaviour, infer patterns from sparse observations, and modify themselves in response.

The platform observes these experiments at population scale. The asymmetry of legibility therefore allows it to learn from participants faster than participants can learn from it. The resulting adaptations provide additional data, enabling still finer personalization. A feedback loop forms:
\[
\text{opacity}
\rightarrow
\text{participant experimentation}
\rightarrow
\text{platform observation}
\rightarrow
\text{better participant models}
\rightarrow
\text{more precise extraction}.
\]

Survivorship bias then supplies a cultural explanation for the unequal outcomes. Exceptional winners become highly visible while unsuccessful attempts disappear. The giant teddy bear converts a low-probability upper tail into evidence that the system rewards merit. Participants search the winners for causal traits and imitate them. The denominator remains obscure.

Vocational awe supplies the motivational energy required to sustain participation despite weak expected returns. Creative workers, teachers, caregivers, researchers, and others continue because the activity itself matters. Their intrinsic motivation lowers the compensation required to keep labour flowing. Meaning becomes economically capturable.

Talent reserve stabilizes the arrangement by ensuring that individual refusal does not easily constrain supply. A large population remains available to replace participants who demand better terms or refuse platform-compatible forms. Supply becomes abundant precisely where routing remains scarce.

Representational narrowing then begins to reshape the activity itself. Metrics determine allocation; participants adapt to metrics; adaptation makes measured forms more common; the resulting data confirm the metrics; and alternatives receive insufficient exposure to generate evidence of their value:
\[
\text{proxy}
\rightarrow
\text{allocation}
\rightarrow
\text{adaptation}
\rightarrow
\text{observed preference}
\rightarrow
\text{proxy reinforcement}.
\]

Human practices become increasingly legible to the platform because successful continuation requires adaptation to the platform's coordinate system. What cannot be ranked, measured, categorized, retained, converted, or monetized acquires an economic disadvantage. The representational constraint becomes a selective environment.

At this stage, automation encounters a population whose economically visible outputs have already been filtered for machine-legible properties. Artificial systems need not reproduce the entire original practice. They need to reproduce enough of the selected projection to satisfy the institution's metrics:
\[
M \approx P_{\theta}(H).
\]
The institution can then mistake equivalence within its representation for equivalence with the practice from which the representation was derived.

Automation further expands supply. Production becomes cheaper, but attention remains finite. Scarcity therefore migrates upward toward recommendation, trust, provenance, ranking, and access. The intermediary controlling those chokepoints can become more powerful even as the objects being mediated become more abundant.

Intersubjectivity simultaneously weakens because participants increasingly encounter one another through the platform's projections. Creators infer audiences from analytics, audiences infer culture from recommendations, workers infer markets from individualized offers, and sellers infer demand from ranked exposure. The intermediary's allocation decisions become part of the evidence through which everyone else models the world.

The system can then represent its own outputs as revealed preference. Users keep participating, therefore they must prefer the service; workers accept the wages, therefore the wages must clear the market; audiences consume the recommended material, therefore recommendation must reflect demand; creators imitate successful forms, therefore those forms must represent authentic cultural preference. Dependency disappears from the interpretation of behaviour.

The counterfactual has now been appropriated. The platform compares itself with absence, its winners with failure, its allocations with doing nothing, and its present capabilities with a historical world before those capabilities existed. Alternatives whose development was prevented by the platform's own accumulated dominance cannot appear as empirical competitors because they were never permitted to become actual enough to generate observations.

At the limit, the institution becomes self-validating. It controls the allocation, observes the behaviour produced by the allocation, interprets that behaviour as preference, and uses the inferred preference to justify the next allocation. The loop is formally coherent:
\[
M_t
\rightarrow
A_t
\rightarrow
B_t
\rightarrow
\widehat{P}_{t+1}
\rightarrow
M_{t+1},
\]
where \(M_t\) is the platform model, \(A_t\) its allocation, \(B_t\) the behaviour observed under that allocation, and \(\widehat{P}_{t+1}\) the inferred preference used to update the model. What is missing is an independently grounded observation capable of distinguishing preference from adaptation to previous allocation.

This is the deepest sense in which enshittification is a collapse of external constraint. The platform loses the outside against which its internal representation could be corrected. Competitors no longer supply a sufficiently reachable alternative, workers cannot coordinate refusal, users cannot modify the service, creators cannot carry audiences elsewhere, and observed behaviour is increasingly endogenous to the platform's own decisions.

The system has not become unconstrained. It has become recursively constrained by itself.

Disenshittification must therefore reintroduce independent constraints at several levels simultaneously. Competition supplies alternative institutional observations. Interoperability preserves relations across institutional boundaries. Portability preserves accumulated state. Labour organization reconstructs common knowledge and collective refusal. Modification rights allow participants to reject components without abandoning entire networks. Privacy and limits on inferential discrimination preserve bargaining uncertainty. Public-interest research and auditing create observational channels outside the platform's own analytics. Antitrust prevents success in one domain from becoming unilateral control over future contests.

These mechanisms differ in implementation but share a formal role. They introduce states and transitions whose existence does not depend entirely upon the incumbent's objective function. They restore an outside.

The architecture of enshittification can therefore be summarized not as a moral decline from generosity to greed but as a trajectory from externally constrained coordination toward internally self-validating allocation. The platform initially participates in a world containing alternatives. It gradually becomes the environment through which alternatives must be reached. Once this occurs, the distinction between what participants prefer and what the platform permits them to reach becomes increasingly difficult to observe.

The appropriate measure of disenshittification is consequently not whether the platform becomes nicer. It is whether the distinction becomes observable again.

\section{Conclusion: Keeping the Outside Open}

Doctorow's account begins with an ugly word because the phenomenon it names is familiar before it is formally understood. People know what it feels like when a service they depended upon becomes progressively worse while somehow becoming harder to leave. They recognize the search results that contain more obstruction, the marketplace in which visibility increasingly requires payment, the social network that shows less of what they deliberately chose to follow, the streaming system in which enormous abundance coexists with poor compensation for creators, and the labour platform whose apparent flexibility conceals increasingly precise control over the terms of work. ``Enshittification'' succeeds as a word because it identifies the phenomenology of deterioration. Its deeper value is that it directs attention toward the structure that makes deterioration sustainable.

Interpreted through the frameworks developed here, that structure is best understood as a contraction of reachable counterfactuals. Platforms acquire power not merely by becoming large but by becoming difficult to route around. Their early usefulness attracts relationships, their relationships generate switching costs, their switching costs weaken refusal, and weakened refusal expands the interval within which extraction can occur. Monopoly and monopsony are two views of the same chokepoint when one institution controls the route through which different populations become economically visible to one another.

The resulting power is not merely economic. It is epistemic. The platform observes populations at a scale unavailable to individual participants while revealing comparatively little about its own decision function. Workers, creators, sellers, and users therefore construct causal models from local outcomes generated by a system that can simultaneously model them in return. Algorithmic superstition is not a failure of intelligence under these conditions. It is a predictable response to acting inside a partially observable environment whose transition rules can change in response to the actor.

The giant teddy bear converts this uncertainty into aspiration. Exceptional winners become visible evidence that enormous rewards are possible, while the denominator of unsuccessful attempts remains distributed and obscure. Survivorship is then translated into merit. The successful participant appears to have discovered the hidden rule, and everyone else is encouraged to search themselves for the variable explaining their failure. Selection effects become biographies.

Vocational awe deepens the asymmetry because the participants often care intrinsically about the activity being mediated. The musician wants to make music, the writer wants to write, the teacher wants to teach, and the researcher wants to understand. Their willingness to continue despite inadequate compensation becomes an exploitable resource. The institution captures not only labour but attachment to labour.

The talent reserve ensures that this attachment exists in abundance. Millions of people can remain available to produce, experiment, imitate, and compete for scarce access to economically meaningful attention. The platform benefits even from unsuccessful participants because their activity supplies content, data, variation, competitive pressure, and examples from which successful forms can be identified.

Scarcity consequently migrates. Digital objects can become nearly free to reproduce while the routes through which they acquire social and economic significance become increasingly concentrated. Attention, recommendation, trust, ranking, reputation, and access remain scarce. The intermediary that controls these scarce relations can prosper amid overwhelming abundance.

Representational narrowing then converts institutional power into cultural form. Platforms must represent heterogeneous human purposes through measurable variables, and those variables become consequential when allocation depends upon them. Creators adapt to what is measured. Audiences encounter what survives the adaptation. The resulting behaviour is recorded as evidence of preference. What began as a proxy gradually participates in constructing the reality it claims merely to measure.

This is why the relationship between platform capitalism and artificial intelligence cannot be understood solely as a sudden technological substitution. Before machines replace people, institutions can reshape human activity around machine-legible coordinates. Human creators are selected for compatibility with measurable objectives, their successful outputs become examples of what the system values, and automated systems subsequently reproduce the selected projection. Replaceability has partly been manufactured before it is discovered.

None of this establishes an ontological boundary between human and machine creativity, nor does the argument require one. Artificial systems can participate in genuine creation, and human beings can produce formulaic work. The important distinction concerns the legitimacy of the representation through which equivalence is judged. A machine's ability to satisfy a platform's objective function does not by itself establish equivalence with the broader practice that the objective function only partially represents.

The same discipline applies to proposed remedies. Diagnosis does not automatically warrant repair. Antitrust, interoperability, portability, transparency, labour organization, privacy, modification rights, and regulation must themselves be evaluated against admissible alternatives, including the null intervention. An intervention can address a genuine problem while producing new chokepoints or irreversible harms. Structural criticism does not release reform from counterfactual scrutiny.

What unifies the strongest remedies is not that they guarantee good outcomes but that they restore corrigibility. They make it possible for participants to act upon dissatisfaction without destroying everything they value about the underlying relation. They separate the social graph from the social platform, reputation from the marketplace, audience from the distributor, device from manufacturer-controlled software, and collective bargaining capacity from the isolated worker's decision.

They make refusal survivable.

This is why contestability is a more durable opposite of enshittification than quality. Quality can be supplied voluntarily by a monopolist. Contestability makes continued quality less dependent upon voluntary restraint. An institution behaves differently when users can leave without losing their relationships, workers can refuse without standing alone, creators can migrate without abandoning audiences, competitors can interoperate without recreating an entire network, and independent observers can test the institution's account of its own behaviour.

Contestability also repairs intersubjectivity. Participants regain ways of learning about one another that are not wholly filtered through the intermediary. Workers can discover common treatment, creators can discover audiences beyond platform analytics, users can compare alternative experiences, and society can generate observations outside the system whose claims are being evaluated. The platform becomes one representation among others rather than the environment mistaken for reality itself.

The resulting political objective is neither maximal openness nor maximal control. Some constraints increase freedom because they prevent powerful actors from destroying the conditions under which future alternatives remain reachable. Some opacity protects bargaining and experimentation because complete legibility would permit individualized extraction. Some redundancy is valuable because an apparently inefficient second path becomes essential when the first path is abused. Some unused capacity deserves preservation because a system optimized entirely for present objectives cannot easily produce futures that those objectives do not yet represent.

This is where sprezzatura, talent reserve, refusal, counterfactual discipline, intersubjectivity, and representational constraint finally converge. A healthy social system must preserve a remainder that its dominant evaluator has not completely absorbed. There must remain capacities not yet priced, relations not owned by the intermediary through which they currently pass, alternatives not rendered impossible by the success of the incumbent, and observations capable of contradicting the model through which the incumbent understands the world.

The ``outside'' need not be geographically or technologically external. It can consist of a competing institution, an interoperable client, a union, an independent audit, a portable archive, a public standard, a different evaluative regime, an artistic practice that refuses optimization, or simply a relationship that survives migration between platforms. What matters is that some consequential structure remains beyond unilateral revision by the institution being constrained.

Doctorow's distinction between optimism, pessimism, and hope is therefore not an inspirational coda to the analysis. It follows from its formal structure. Optimism and pessimism predict trajectories through a possibility space treated as given. Hope acts upon the possibility space. It creates organizations, rights, protocols, precedents, institutions, and shared knowledge that alter which transitions become reachable from the next state.

Hope does not require believing that the best outcome will occur. It requires refusing to mistake the currently reachable for the permanently possible.

The history of concentrated power provides no guarantee that chokepoints will be dismantled, but it also provides no warrant for treating them as natural law. Standard Oil once appeared to embody an industrial logic that could not realistically be reversed. Labour movements repeatedly confronted employers whose control over work appeared structurally overwhelming. Communications systems, technical standards, copyright regimes, and market institutions have repeatedly been reconstructed through political struggle. The relevant lesson is not that history bends automatically toward disenshittification. Nothing in the argument supports such teleology. The lesson is that institutions alter the topology within which later action occurs.

A successful reform therefore need not solve everything to matter. It can preserve an edge that would otherwise disappear. An interoperability requirement can make a future competitor viable. A portable reputation can make a future refusal affordable. A union can turn private knowledge into common knowledge. A prohibition on discriminatory inference can preserve a bargaining interval. An independent distribution system can sustain a cultural form long enough for an audience to discover that it wanted something the dominant platform never showed it.

These interventions matter because possibilities can have causal force before they are exercised. The credible possibility of exit disciplines the incumbent without requiring everyone to exit. The credible possibility of a strike changes bargaining without requiring every negotiation to become a strike. The possibility of modifying software constrains the manufacturer without requiring every user to become a programmer. The possibility of an alternative allocation prevents the realized allocation from presenting itself as inevitable.

Power therefore operates partly by governing the counterfactual, and freedom operates partly by refusing to surrender it.

The deepest danger of enshittification is not that everything becomes unpleasant. People can adapt to unpleasant systems with extraordinary ingenuity. The deeper danger is that adaptation becomes so complete that the system's contingent constraints are mistaken for the shape of reality itself. Creators learn what ``content'' is, workers learn what ``flexibility'' means, audiences learn what they supposedly prefer, users learn what privacy they supposedly do not need, and institutions learn that continued participation proves consent. The map closes over the territory.

Disenshittification begins wherever that closure is broken.

It begins when the worker can compare the offer with an independently visible alternative, when the creator can distinguish an audience from a recommendation system's representation of that audience, when the user can change software without abandoning a network, when the seller can carry reputation across marketplaces, when the public can distinguish a corporation's contribution from the infrastructure and collective activity surrounding it, and when a society can evaluate a spectacular success without forgetting the denominator that produced it.

It begins, in other words, when the counterfactual becomes socially real again.

A good institution should therefore be judged partly by what remains possible when we refuse it. If leaving destroys our relationships, if modification is prohibited, if bargaining requires isolation, if every competitor must reconstruct the network from zero, if every cultural alternative must first satisfy the incumbent's metric in order to become visible, and if every observation available for criticizing the system has already been produced by the system, then formal choice tells us very little about freedom.

The more demanding criterion is whether continuation survives disagreement.

That criterion applies beyond the platforms Doctorow discusses. It applies to markets, workplaces, technical standards, cultural institutions, artificial-intelligence systems, philanthropic structures, public infrastructure, and any other arrangement capable of becoming the necessary environment through which people pursue goods they value. Institutions deserve particular scrutiny when refusing them becomes indistinguishable from refusing the underlying good.

The objective is not a world without intermediaries. Intermediaries solve real problems. Nor is it a world without scale, optimization, algorithms, artificial intelligence, markets, or expertise. Each can enlarge human possibility. The objective is to prevent any particular realization of those capacities from exhausting the alternatives through which it could later be corrected.

A platform should be able to fail without taking the social relations it mediates with it. A model should be able to be wrong without defining the only coordinate in which its error can be expressed. A creator should be able to refuse an intermediary without refusing creation. A worker should be able to reject an offer without revealing the precise point at which desperation will force acceptance. A community should be able to migrate without becoming a new community from zero.

The future should contain exits that lead somewhere.

That is the structural content of hope, and it is why the most important political question raised by enshittification is ultimately not whether today's platforms can be persuaded to become better. It is whether we can construct institutions in which no platform needs to remain good forever because the people who depend upon it retain the capacity to continue without it.

\section{The Residual Outside}

The preceding conclusion identifies contestability with the preservation of an outside, but the concept deserves one final clarification before turning to the formal synthesis and appendices. An outside need not be a pristine domain untouched by institutions, computation, markets, or mediation. No such position is required, and in sufficiently interconnected systems it may not exist. What matters is something weaker and more operational: no representation should possess complete authority over the conditions under which its own adequacy is judged. There must remain some source of variation, observation, refusal, or reconstruction that is not generated entirely by the system whose performance is under evaluation. The outside is therefore relational rather than spatial. It is whatever preserves the possibility that the present representation can encounter a consequential contradiction.

This distinction prevents the argument from collapsing into a simple opposition between centralized and decentralized systems. Decentralization can reproduce concentration through standards, infrastructure, capital requirements, expertise, or informal hierarchy, while centralized institutions can sometimes preserve strong mechanisms of appeal, portability, transparency, and revision. The relevant variable is not topology in the everyday organizational sense but corrigibility. Can the system receive information capable of showing that its own categories are inadequate? Can participants act upon that information? Can an alternative persist long enough to become a genuine comparison rather than a hypothetical objection? Can the institution responsible for an allocation be separated from the relationships whose value made the allocation consequential?

The residual outside is particularly important because no finite representational scheme can guarantee in advance that it has captured every distinction that will later matter. This is true even where the scheme is exceptionally sophisticated. A platform may possess billions of observations, a machine-learning system may contain enormously expressive representations, and an institution may consult extensive expertise. None of these facts eliminates the possibility that the objective being optimized omits a relevant dimension, that the proxy has become endogenous to intervention, or that a previously negligible distinction becomes important under changed conditions. More information inside a representation does not by itself solve the problem of whether the representation's boundary was correctly drawn.

Let a system at time \(t\) possess a representational state
\[
\mathfrak{M}_t = (X_t,\Theta_t,\mathcal{O}_t),
\]
where \(X_t\) denotes the state variables it can represent, \(\Theta_t\) the parameters through which those variables affect decisions, and \(\mathcal{O}_t\) the objective or family of objectives under which outcomes are evaluated. Increasing the predictive accuracy of \(\mathfrak{M}_t\) can improve decisions conditional upon this representational scheme. It does not establish that
\[
X_t = X^{*},
\qquad
\mathcal{O}_t = \mathcal{O}^{*},
\]
for some complete set of relevant variables \(X^{*}\) or uniquely correct objective \(\mathcal{O}^{*}\). Indeed, the framework developed throughout this essay rejects the assumption that such complete objects will generally be available. The practical requirement is therefore not representational perfection but a procedure by which failures of representation can become detectable before they become irreversible.

This is one reason external redundancy is so important. A second institution need not outperform the first at every moment to be valuable. It may preserve a different error distribution. A second search engine, distribution channel, marketplace, client, archive, model, research programme, or evaluative institution can retain observations that the dominant system suppresses. Redundancy is then not merely duplication of function. It is duplication with partially independent failure modes.

If two systems \(M_1\) and \(M_2\) optimize identical objectives from identical observations using effectively identical representations, their numerical multiplicity supplies relatively little epistemic diversity. By contrast, even a smaller or less efficient \(M_2\) can be valuable where
\[
\operatorname{Err}(M_1)
\not\approx
\operatorname{Err}(M_2),
\]
because disagreement between them supplies information that would disappear under a single representational regime. The value of plurality therefore depends partly upon non-equivalence.

This point is easily missed in discussions of efficiency. If one platform performs a function substantially better than every competitor according to the dominant metric, consolidating activity around it can appear rational. Yet consolidation also removes alternative measurements of what ``better'' means. The winner's objective function becomes increasingly capable of defining the performance criterion against which challengers are evaluated. A competitor optimized for a different value can then appear inefficient precisely because the dominant system has succeeded in making its own value the standard of efficiency.

The same problem appears in culture. A small publisher, local venue, independent label, specialized forum, public broadcaster, library, or idiosyncratic curator may look inefficient beside a platform capable of reaching millions of people algorithmically. Yet such institutions can preserve forms of selection whose value lies partly in not being reducible to the dominant platform's selection function. Their apparent redundancy is a reserve of evaluative difference.

The talent reserve discussed earlier therefore has an institutional counterpart. Just as a healthy culture requires human capacities not completely exhausted by present market demand, a healthy institutional ecology requires organizational capacities not completely optimized against the present dominant infrastructure. An alternative that is temporarily less efficient can constitute a reserve against a future in which the dominant system's objective proves inadequate. Destroying every such reserve in the name of current efficiency creates an institutionally monocultural environment.

The analogy to biological monoculture is useful only if kept precise. The claim is not that institutions literally evolve like organisms or that diversity is automatically good. Diversity can preserve harmful practices as easily as beneficial ones. The relevant property is differential response to error. A monoculture becomes fragile when a disturbance exploits a vulnerability shared across the entire population. An institutional monoculture becomes fragile when an error in the dominant representation propagates across nearly every route through which an activity is performed.

Platform concentration can create exactly this condition. A change to one recommendation system can alter cultural visibility across enormous populations. A change to one marketplace's rules can reorganize the economic environment of millions of sellers. A change to one search ranking system can affect which information becomes reachable at planetary scale. A change to one application ecosystem can alter the viability of thousands of complementary businesses. Scale magnifies both successful coordination and representational error.

The residual outside is therefore a form of fault tolerance. It preserves enough non-equivalent structure that errors can remain local rather than becoming universal. This connects the political economy of platforms directly to the diagnostic-coordinate framework. A system cannot reliably diagnose its own failures when every observation available for diagnosis has already been filtered through the mechanism potentially responsible for the failure. Independent reference classes are not luxuries added after optimization. They are part of the epistemic infrastructure required to determine whether optimization remains pointed in the right direction.

This also clarifies why the outside cannot simply be simulated internally. A dominant platform may create internal red teams, multiple ranking models, experimental treatments, user surveys, advisory boards, or alternative objectives. These can be extremely valuable. But an internally generated alternative remains subject to the institution's authority over whether the alternative receives resources, whether its evidence counts, and whether its recommendations are implemented. Internal pluralism is therefore not equivalent to external contestability.

The distinction is analogous to the difference between a system generating a hypothesis against itself and an environment capable of imposing a consequence upon the system when the hypothesis is correct. The first supplies information. The second supplies discipline. An organization can know that users dislike a change and implement it anyway if users cannot leave. It can know that workers regard compensation as unfair and continue if workers cannot coordinate. It can know that creators dislike a distribution regime and retain it if audiences remain captive. Information without countervailing reachability does not necessarily alter the objective.

The outside must therefore possess some causal leverage. It must be capable not merely of disagreeing but of sustaining a continuation in which disagreement matters.

\section{A Formal Synthesis: Enshittification as Reachability Contraction}

The argument can now be stated in a compact formal language. The purpose of the formalization is not to claim that platform political economy can be reduced to a small collection of equations. It is to expose the relations among concepts that ordinary prose can too easily allow to drift apart.

Let a platform-mediated environment at time \(t\) be represented by
\[
\mathcal{E}_t
=
\left(
X_t,
A_t,
\leadsto_t,
\mu_t,
\Lambda_t,
\Pi_t
\right),
\]
where \(X_t\) is the space of institutionally relevant states, \(A_t\) the set of participating agents or classes of agents, \(\leadsto_t\) the admissible reachability relation between states, \(\mu_t\) the distribution of material and informational resources, \(\Lambda_t\) the distribution of legibility among agents, and \(\Pi_t\) the family of representations through which observations are translated into allocations.

For an agent \(i\in A_t\), define the reachable continuation set
\[
\mathcal{R}_i(x,t)
=
\left\{
y\in X_t:
x\leadsto_t^{\,i}y
\right\}.
\]
The existence of an alternative \(y\) in \(X_t\) is insufficient for freedom if
\[
y\notin\mathcal{R}_i(x,t).
\]
This distinguishes nominal possibility from effective possibility.

Define the quality of refusal available to \(i\) with respect to institution \(M\) as a function of the continuation set following refusal:
\[
Q_R(i,M,x)
=
F\!\left(
\mathcal{R}_i(R_M(x),t)
\right),
\]
where \(R_M(x)\) denotes the state reached by refusing or exiting \(M\), and \(F\) evaluates the preservation of goods not intrinsically dependent upon \(M\). A formally available exit can therefore have low refusal quality when the participant loses unrelated accumulated structure.

Let
\[
L_{M\rightarrow i}
\]
denote the degree to which \(M\) can infer decision-relevant properties of \(i\), and
\[
L_{i\rightarrow M}
\]
the degree to which \(i\) can infer the decision rule governing \(M\)'s treatment of \(i\). An informational asymmetry exists when
\[
L_{M\rightarrow i}
\gg
L_{i\rightarrow M}.
\]
This asymmetry becomes extractively significant when the inferred variables include the participant's reservation threshold.

Let \(r_i\) denote the boundary at which participant \(i\) refuses continued participation under a given dimension of extraction. An institution with increasing inferential accuracy can select an individualized term
\[
e_i=r_i-\varepsilon_i,
\qquad
\varepsilon_i>0,
\]
thereby approaching the maximum extractable value compatible with continuation. In this limit, the residual surplus retained by the participant becomes not a general feature of the market but a deliberately minimized margin required to prevent refusal.

Define a platform representation
\[
\Pi_t:H\rightarrow Z_t
\]
mapping a richer domain of human practice \(H\) into a machine-legible state space \(Z_t\). Allocation occurs through an objective
\[
J_t:Z_t\rightarrow\mathbb{R}.
\]
Where participants depend upon allocation for continuation, they receive incentives to transform their behaviour \(h\in H\) so that
\[
J_t(\Pi_t(h))
\]
increases. The representation therefore becomes performative: it does not merely observe \(H\) but alters the distribution of practices within \(H\).

Let \(P_t(h)\) denote the observed prevalence of practice \(h\). Under repeated allocation, one can obtain
\[
P_{t+1}(h)
=
G\!\left(
P_t(h),
J_t(\Pi_t(h)),
C_t(h)
\right),
\]
where \(C_t(h)\) represents other social and material constraints. If subsequent model updates use \(P_{t+1}\) as evidence of preference, the system risks confusing adaptation to prior allocation with independent demand.

We may call this \emph{endogenous preference inference}. Formally, if observed behaviour \(B_t\) depends upon prior allocation \(A_{t-1}\),
\[
B_t=f(P_t,A_{t-1}),
\]
then estimating \(P_t\) directly from \(B_t\) without identifying the effect of \(A_{t-1}\) produces a confounded estimate. The platform's own intervention belongs inside the causal model.

Intersubjectivity collapse occurs when agents increasingly estimate one another through projections controlled by \(M\). For populations \(C\) and \(U\), suppose observed interaction is
\[
O(c,u)
=
E_M(c,u)P(c,u),
\]
where \(E_M\) is platform-controlled exposure and \(P\) the conditional propensity for meaningful interaction. Where \(E_M\) is hidden, a low value of \(O\) does not identify a low value of \(P\). Nevertheless, participants can rationally mistake one for the other. The platform's allocation becomes confused with the population's preference.

Define the talent reserve
\[
\mathcal{T}_t
=
\left\{
i\in A_t:
i \text{ possesses relevant productive capacity but lacks stable access to remunerative allocation}
\right\}.
\]
A large \(\mathcal{T}_t\) reduces the bargaining power of any particular supplier where the platform can substitute among members at low cost. The reserve can nevertheless remain productive for the platform through experimentation, content supply, data generation, and competitive pressure.

Vocational awe modifies reservation thresholds because participants receive intrinsic utility from continuation of the activity. If monetary compensation is \(w_i\), intrinsic activity value is \(v_i\), and participation cost is \(c_i\), then continuation can occur where
\[
w_i+v_i\geq c_i.
\]
An institution capable of appropriating the effect of \(v_i\) can reduce \(w_i\) while preserving participation. Intrinsic meaning thereby becomes an implicit subsidy unless countervailing institutions prevent its capture.

The giant teddy bear gambit operates through selective visibility. Let \(S\) denote exceptional success and \(N\) the number of attempts. Public observation disproportionately samples
\[
P(i\mid S)
\]
rather than the distribution
\[
P(S\mid i,\text{participation conditions}).
\]
Traits prevalent among visible winners are then misidentified as sufficient causal explanations for the outcome. The platform benefits when exceptional rewards increase entry into \(\mathcal{T}_t\) at lower cost than broad improvement of baseline conditions.

Artificial scarcity can now be defined relationally. Let productive supply \(Y_t\) increase while economically consequential routing capacity \(K_t\) remains concentrated. Then scarcity has not disappeared when
\[
Y_t\rightarrow\infty
\]
if access to \(K_t\) remains necessary for remuneration. Generative artificial intelligence can increase \(Y_t\) dramatically while increasing the strategic value of control over \(K_t\), particularly where human attention remains finite.

Manufactured replaceability occurs when the platform first projects human practice into a selected representation and subsequently compares machines with that projection:
\[
H
\xrightarrow{\Pi_t}
Z_t
\xrightarrow{\text{selection}}
Z_t^{*},
\qquad
M_{\mathrm{AI}}
\approx
Z_t^{*}.
\]
The inference
\[
M_{\mathrm{AI}}\approx H
\]
is invalid without additional evidence establishing that the projection and selection preserved the dimensions relevant to the practice being evaluated.

Finally, define an enshittification trajectory as a sequence
\[
\mathcal{E}_0,
\mathcal{E}_1,
\ldots,
\mathcal{E}_T
\]
such that the intermediary's extractive capacity increases while the effective counterfactual space available to dependent participants contracts. A minimal structural signature is
\[
\frac{d}{dt}\,
\operatorname{Reach}_i(\mathcal{E}_t)<0
\]
for a substantial class of dependent participants, together with
\[
\frac{d}{dt}\,
\operatorname{Disc}_M(\mathcal{E}_t)>0,
\]
where \(\operatorname{Disc}_M\) denotes the intermediary's discretionary capacity to alter terms without triggering proportionate loss of participation.

The two quantities are related. As reachable alternatives contract, discretionary extraction can expand. As discretionary extraction finances further enclosure, reachable alternatives can contract further. Enshittification therefore admits a positive feedback structure:
\[
\text{reachability contraction}
\rightarrow
\text{greater discretion}
\rightarrow
\text{greater extraction}
\rightarrow
\text{greater capacity for enclosure}
\rightarrow
\text{reachability contraction}.
\]

A disenshittifying intervention breaks this loop when it increases independently controlled continuation without merely shifting discretion to another chokepoint. Let \(I\) be an intervention. A necessary, though not sufficient, condition for structural improvement is
\[
\Delta \operatorname{Reach}_i(I)>0
\]
for relevant participants without a compensating increase in irreversible dependency elsewhere. This is why the repair criterion developed earlier includes reachability, uncertainty, cost, and irreversible risk rather than treating immediate welfare gain as sufficient.

The formal synthesis therefore yields a compact proposition.

\begin{proposition}[Contestability Principle]
A platform-mediated system is structurally resistant to enshittification only insofar as deterioration by the intermediary increases the relative attractiveness and effective reachability of continuations not controlled by that intermediary.
\end{proposition}

\begin{proof}
Suppose deterioration by intermediary \(M\) does not increase the effective attractiveness of any independently controlled continuation. Then participants harmed by the deterioration possess no correspondingly stronger transition away from \(M\), whether because alternatives are absent, unreachable, nonportable, noninteroperable, collectively inaccessible, or themselves controlled by \(M\). The deterioration therefore need not generate a proportional loss of participation or control. In the absence of some other external discipline, \(M\) can continue deteriorating the relevant dimension until the residual surplus required to retain participants is approached.

Conversely, suppose deterioration by \(M\) makes one or more independently controlled continuations increasingly attractive and effectively reachable. Then deterioration generates a countervailing response: exit, modification, substitution, coordinated refusal, migration, or another transition capable of reducing the value captured by \(M\). The intermediary's discretion is therefore bounded by the credibility of the alternative. Contestability does not guarantee that \(M\) will never deteriorate, but it ensures that deterioration alters the conditions of its own continuation. This is precisely the external discipline absent from a fully enclosed chokepoint.
\end{proof}

The proposition should not be mistaken for the claim that every alternative must be equally convenient or that every switching cost is illegitimate. Some state is intrinsically difficult to transfer, some integrations produce real efficiencies, and some network effects arise from the underlying activity rather than artificial enclosure. The criterion is comparative. It asks whether institutional design unnecessarily transforms these difficulties into unilateral discretionary power.

A second proposition follows from the analysis of legibility.

\begin{proposition}[Asymmetric Legibility Principle]
Where an intermediary can estimate participant-specific reservation thresholds substantially better than participants can estimate the intermediary's decision rule, increased predictive accuracy can increase extractive capacity even when productive efficiency remains unchanged.
\end{proposition}

\begin{proof}
Let participant \(i\) accept terms within an interval bounded by reservation threshold \(r_i\), and suppose the intermediary initially possesses uncertainty over \(r_i\). To avoid rejection, it must leave an informational margin between its offer and its estimate of \(r_i\). As the estimate becomes more precise, the required margin can shrink. The intermediary can therefore move the offered term toward \(r_i\) without any change in the productive technology underlying the exchange. The resulting gain derives from informational redistribution of surplus rather than increased productive output. If participants do not receive a corresponding improvement in their knowledge of the intermediary's reservation threshold or allocation rule, bargaining asymmetry increases.
\end{proof}

A third proposition captures the representational argument.

\begin{proposition}[Proxy Endogeneity Principle]
Where allocation depends upon a measurable proxy for preference, observed increases in behaviour associated with that proxy cannot, without an independent identification strategy, be interpreted as equivalent increases in the underlying preference.
\end{proposition}

\begin{proof}
Let \(B\) denote observed behaviour, \(P\) the underlying preference of interest, and \(A\) an allocation mechanism conditioned upon proxy \(Z\). If \(A\) changes exposure, incentives, or survival probabilities for behaviours associated with \(Z\), then
\[
B=f(P,A).
\]
An observed change in \(B\) following repeated operation of \(A\) is therefore compatible with changes in \(P\), changes induced by \(A\), or both. Since the platform's intervention is a common cause of the observation and the environment in which the observation occurs, \(B\) alone does not identify \(P\). An independent comparison or causal identification strategy is required.
\end{proof}

Together these propositions identify three conditions that recur across Doctorow's examples: the disappearance of credible alternatives, the accumulation of asymmetric information, and the conversion of endogenous behavioural responses into apparent evidence of preference. Their interaction explains why platform degradation can become simultaneously profitable, difficult to diagnose locally, and culturally self-reinforcing.

\section{Coda: The Right to Another Continuation}

There is a temptation, after describing systems this large, to end with a demand for better values. Platforms should care more about users, employers should respect workers, wealthy people should exercise responsibility, engineers should resist harmful optimization, and creators should refuse to surrender their practices to metrics. None of these claims is objectionable. None is structurally sufficient.

Values matter most when institutions permit them to survive contact with incentives. A conscientious executive governing an enclosed chokepoint can produce a period of unusually good stewardship, but the relationships accumulated during that period remain available for extraction by whoever controls the institution later. A benevolent platform can therefore create the dependency from which a less benevolent successor profits. The better the platform initially behaves, the more relationships may accumulate inside it, and the larger the eventual extraction opportunity can become.

This is the paradox at the heart of enshittification. Trust can manufacture the conditions for later betrayal.

The appropriate response cannot be permanent suspicion. Refusing every useful institution because it might later become exploitative would destroy many of the cooperative gains that make complex social life possible. The alternative is to distinguish trust from surrender. We can trust an institution while refusing to make the continuation of everything we value depend upon its perpetual trustworthiness.

That distinction is familiar in engineering. A bridge is not designed on the assumption that every component will remain perfect forever. Reliable systems anticipate failure, preserve redundancy, isolate faults, and permit repair. The presence of a backup does not express hostility toward the primary system. It expresses recognition that no component should be required to deserve infinite trust.

Social infrastructure deserves the same discipline.

The platform through which friends communicate should not need to remain benevolent forever if the relationships can survive its failure. The marketplace should not need to remain fair forever if reputation and trade can migrate. The streaming intermediary should not need to remain generous forever if artists and listeners retain independent routes toward one another. The employer should not need to remain enlightened forever if workers possess collective capacity to refuse deteriorating terms. The machine-learning system should not need to remain perfectly aligned with human purposes forever if its representation can be challenged by evidence and institutions outside its own objective.

A robust society does not solve the problem of power by finding actors worthy of permanent power. It limits the consequences of discovering that they were not.

This reframes the meaning of freedom in a way that is less dramatic than absolute autonomy but more useful for institutional analysis. Freedom is not the absence of constraint, because every meaningful activity depends upon constraints. Language constrains expression enough to make communication possible; protocols constrain computers enough to make interoperability possible; contracts constrain parties enough to make coordination possible; physical embodiment constrains human beings enough to give action determinate form. The relevant question is whether constraints preserve or destroy the capacity for alternative continuation.

A constraint that enables independent systems to communicate can enlarge the reachable world. A constraint that forces every communication through one owner can contract it. A rule that prevents discriminatory exploitation of private vulnerability can preserve bargaining freedom. A rule that prohibits users from modifying devices they own can remove it. The language of constraint therefore cannot substitute for an analysis of topology.

The same is true of friction. Platform ideology has often treated friction as something to eliminate. Some friction is needless obstruction, and eliminating it can produce enormous social benefit. But friction can also be the material through which deliberation, privacy, bargaining, and refusal become possible. A perfectly frictionless path from desire to transaction may also be a perfectly frictionless path from observation to exploitation. A perfectly personalized price can eliminate the informational friction that once allowed a consumer to retain surplus. A perfectly personalized wage can eliminate the bargaining uncertainty that once protected a worker's reservation threshold.

Not every inefficiency is a defect. Some inefficiencies are reserves of agency.

The larger theoretical point is that optimization requires a boundary. To optimize is to choose some variables as relevant, some objective as desirable, and some horizon over which improvement will be measured. Everything outside that frame becomes, at least temporarily, residual. The danger begins when the optimizer acquires sufficient power to reorganize the environment until the residual disappears.

At that limit, the world becomes easy to predict because everything capable of surprising the model has been deprived of the ability to continue.

Such a world would look remarkably efficient from inside the model.

The preservation of an outside is therefore not hostility to intelligence, computation, markets, or coordination. It is a condition for their continued correction. Intelligence requires the possibility of being surprised. Science requires observations capable of falsifying expectations. Markets require alternatives capable of disciplining incumbents. Democracy requires opposition capable of replacing government. Culture requires experiments whose value cannot be inferred entirely from existing preferences. Individuals require enough opacity that knowing them does not become equivalent to pricing every vulnerability they possess.

Each institution remains healthy only while something consequential escapes complete assimilation into its present representation.

This is why Doctorow's insistence on hope is more rigorous than it initially appears. Hope is not confidence in a favourable endpoint. It is commitment to preserving the incompleteness of the present one. It refuses the proposition that because no alternative is presently visible, no alternative can be constructed; because no competitor presently exists, competition is impossible; because workers presently accept a term, the term reflects their preference; because audiences presently consume a form, the form exhausts their taste; because machines presently reproduce a task, the task exhausts the human practice from which it was abstracted.

Hope preserves the missing variable.

Sometimes that variable is a law not yet enacted, a union not yet organized, a protocol not yet standardized, an independent client not yet written, a cooperative not yet formed, an artistic form not yet legible, a research programme whose value cannot yet be measured, or a population whose shared condition has not yet become common knowledge. None exists strongly enough in the present to function as an empirical alternative. Yet each can become part of the next state's reachability structure if enough of the path toward it remains open.

The political obligation is therefore partly negative: do not close paths merely because nobody is currently walking them. Do not treat unused interoperability as useless, unexercised rights as unnecessary, redundant institutions as waste, unmonetized capacities as inefficiency, private information as unrealized data, or unpopular experiments as proof of permanent irrelevance. Their value can lie precisely in remaining available for conditions that the present optimizer cannot foresee.

The corresponding positive obligation is to construct institutions whose success does not destroy the conditions under which they can later be replaced. This is a demanding criterion because it asks successful actors to preserve routes around themselves. Yet that is precisely why voluntary corporate benevolence cannot bear the entire burden. The preservation of alternatives must itself become institutional.

The resulting principle is simple enough to state without notation.

No intermediary should own the only path between people who give one another the reason the intermediary has value.

Creators and audiences give the streaming platform its value. Buyers and sellers give the marketplace its value. Drivers and riders give the dispatch platform its value. Friends and communities give the social network its value. Writers and readers give the publishing platform its value. The intermediary can improve those relations enormously, but improvement does not transform the relations into its property.

When the relation can survive the intermediary, the intermediary must continue earning its place within the relation.

When the relation cannot survive it, usefulness has become sovereignty.

Enshittification is the process by which the first condition drifts toward the second.

Disenshittification is the work of making the intermediary contingent again.

And the minimal political right capable of expressing that objective is neither a right to perfect service nor a right to infinite choice. It is more elementary and more demanding:

\begin{quote}
\emph{the right to another continuation.}
\end{quote}

\section{Disenshittification as Institutional Memory}

The right to another continuation also reveals a dimension of the problem that is easily overlooked when enshittification is described primarily as declining service quality. Exit is valuable only when enough of what has already been accumulated can survive the exit. A formally open door means relatively little if walking through it requires abandoning years of relationships, reputation, correspondence, purchased media, creative archives, annotations, learned configurations, professional history, or collective knowledge. The decisive question is therefore not merely whether another platform exists but whether the participant can reach it without becoming institutionally newborn.

This is a problem of memory. Platforms accumulate state on behalf of their participants, but the accumulated state is heterogeneous. Some of it belongs straightforwardly to the user: photographs, messages, documents, playlists, source files, purchase records, and other explicit artifacts. Some is relational: followers, subscribers, reputation, reviews, collaborative histories, social graphs, and records of interaction. Some is inferential: recommendation profiles, trust scores, ranking histories, fraud assessments, behavioural models, and other representations generated by the platform from prior activity. Some is infrastructural: integrations, identifiers, links, references, embedded objects, APIs, and dependencies that have grown around the participant's presence. The difficulty of exit depends upon how much of this state remains usable outside the institution that accumulated it.

The usual language of ``data portability'' captures only part of this problem. A downloadable archive can preserve information while failing to preserve continuation. A user may receive a machine-readable list of every person they followed and still possess no mechanism for reconnecting with those people elsewhere. A creator may download every uploaded work while losing the audience relationships through which those works acquired economic significance. A seller may preserve transaction records while losing reputation accumulated through thousands of successful exchanges. A worker may preserve a record of completed tasks while losing the platform-specific trust state through which future work becomes accessible. Preservation is therefore weaker than reachability.

Let \(S_i(t)\) denote the accumulated state associated with participant \(i\) at time \(t\), and let an exit operation from platform \(M\) to alternative \(N\) preserve some subset
\[
P_{M\rightarrow N}\bigl(S_i(t)\bigr)
\subseteq
S_i(t).
\]
The existence of a portability mechanism tells us that this subset is nonempty. It does not tell us whether the preserved state remains functionally meaningful. A contact list preserved as a static file is not equivalent to an interoperable social graph; an exported reputation score is not equivalent to reputation recognized by another marketplace; a downloaded message archive is not equivalent to continuing the conversations it records.

We therefore need to distinguish \emph{state preservation} from \emph{continuation preservation}. State preservation asks whether prior information survives. Continuation preservation asks whether the relations and affordances constituted through that information remain reachable after migration. If \(F_M(S)\) denotes the set of future actions made possible by state \(S\) within platform \(M\), then a meaningful migration should preserve not merely bytes but some sufficiently important portion of the affordance structure:
\[
F_N\!\left(P_{M\rightarrow N}(S)\right)
\approx_{\mathcal{A}}
F_M(S),
\]
where \(\approx_{\mathcal{A}}\) denotes equivalence relative to an independently justified set of relevant affordances. Exact equivalence is neither possible nor necessarily desirable. The alternative platform may deliberately work differently. The point is that migration should not destroy unrelated accumulated capacities merely because those capacities happened to be instantiated through the incumbent.

This distinction is central to the political economy of lock-in because accumulated state changes the meaning of consent over time. A user who voluntarily joins a platform at \(t_0\) may face almost no switching cost. Ten years later, the same person can rationally remain after the service deteriorates because departure would destroy accumulated relations. Continued participation at \(t_{10}\) cannot therefore be interpreted as a repetition of the original choice under identical conditions. The option set has changed as a consequence of the history of participation itself.

The trajectory has the form
\[
x_0
\xrightarrow{\text{participation}}
x_1
\xrightarrow{\text{accumulation}}
x_2
\xrightarrow{\text{degradation}}
x_3,
\]
where refusal at \(x_3\) is evaluated not against \(x_0\) but against the loss of structure accumulated between \(x_0\) and \(x_3\). This is why ``you can always leave'' is frequently an inadequate response to platform criticism. It treats exit as a state-independent action when the cost of exit is itself a product of the path.

The issue is especially clear in communication systems. If a telephone provider deteriorates while telephone numbers remain portable and calls remain interoperable across networks, switching providers need not mean abandoning everyone who knows the old number. The relation has some independence from the intermediary. A closed social platform can bind the relation much more tightly to the institution. The user may be permitted to leave while the friendship network remains technically imprisoned behind the exit.

This suggests a general principle:
\[
\text{exit freedom}
\neq
\text{permission to terminate an account}.
\]
A more adequate approximation is
\[
\text{exit freedom}
\propto
\frac{\text{preserved valuable continuation}}
{\text{unavoidable migration loss}}.
\]

The numerator is often missing from conventional accounts of competition. Two services can compete nominally while migration between them destroys enough accumulated state that the incumbent retains substantial power. The presence of alternatives in a market therefore does not establish contestability. The alternatives must be reachable from the participant's actual historical position.

This is where the Chain of Memory framework becomes relevant. A participant's state is not merely a collection of independent objects but an ordered accumulation of commitments, relations, transformations, and provenance. What matters is not only that an artifact exists but how it came to occupy its present role. A reputation score, for example, compresses a history of interactions. A subscriber count compresses a history of audience formation. A collaborative repository contains not merely files but attribution, discussion, branching, revision, and dependence. Migration that preserves only the terminal object while destroying the chain can preserve extensional content while losing the structure that made the content trustworthy or actionable.

A platform can therefore create lock-in without explicitly prohibiting export. It need only make the exported representation sufficiently impoverished that the chain cannot be reconstructed elsewhere. The distinction resembles the difference between copying the final state of a program and preserving the history required to understand why that state is admissible. Endpoint identity does not imply trajectory equivalence.

This observation complicates proposals for portability because maximal portability is neither technically trivial nor normatively innocent. Relational state involves other people. A user's right to export a social graph cannot automatically entail a right to export private information about everyone connected to them. Reputation systems can encode context-specific judgments that become misleading when transferred between domains. Moderation histories can contain disputed or sensitive information. Recommendation models can embody both user-generated signals and proprietary infrastructure. A serious portability regime therefore requires admissibility constraints rather than indiscriminate copying.

The objective should be to identify which components of accumulated state are necessary for meaningful continuation and which cannot legitimately be transferred without additional consent or transformation. This is precisely a diagnostic-coordinate problem. We need a reference class for determining what counts as preservation, a model of the structural effects of migration, and a repair objective in which doing nothing remains an available baseline. ``Export everything'' is no more automatically justified than ``export nothing.''

Nevertheless, the existence of difficult boundary cases should not obscure the asymmetry of the current arrangement. Platforms routinely treat accumulated relational state as an asset when valuing themselves while treating the same state as nonportable when participants attempt to leave. The social graph contributes to corporate valuation because it makes the service difficult to reproduce, yet the individual relationships constituting that graph are presented as though they were inseparable from the platform. Collective activity is capitalized upward and immobilized downward.

This is another instance of counterfactual appropriation. The platform's value depends upon the proposition that participants are valuable to one another, while its bargaining power depends upon the proposition that those relations cannot easily continue elsewhere. It therefore captures the difference between the social relation and the infrastructure currently mediating it.

The economic value of interoperability follows directly. Interoperability weakens the identity
\[
\text{relation}=\text{platform}.
\]
A participant can change one component of the infrastructure while preserving interaction with participants who have not changed it. Migration becomes a sequence of local transitions rather than a collective leap:
\[
M
\rightarrow
(M,N)
\rightarrow
N.
\]
The intermediate state matters. Without it, everyone may prefer the final state while nobody can rationally move first.

This is an instance of a more general principle of transition design. Institutional alternatives fail not only because their endpoints are unattractive but because no admissible path connects the present state to the preferred endpoint. A proposed alternative can therefore be superior in equilibrium and irrelevant in practice. Political economy must study paths, not merely destinations.

The same insight applies to labour organization. A workforce in which everyone would benefit from coordinated refusal can remain trapped if the transition from isolated acceptance to collective bargaining requires individuals to incur prohibitive risks before coordination exists. Unions, strike funds, legal protections, and organizing institutions create intermediate states through which the collectively preferred configuration becomes reachable. They do not merely change the destination. They alter the path topology.

Disenshittification is consequently a problem of transition engineering. It is insufficient to describe a better platform, a better labour relation, a better distribution system, or a better technical standard. One must preserve or construct paths from the actual degraded state toward those alternatives without requiring participants to sacrifice the very goods they are attempting to protect.

This is why institutional memory matters. Every successful reform inherits a world already structured by the institution it seeks to change. Audiences already live on particular platforms, creators already depend upon particular distribution systems, businesses already possess platform-specific reputations, software already depends upon proprietary interfaces, and workers already organize their schedules around existing systems. Reform that ignores this accumulated state can unintentionally punish the participants most deeply embedded in the system while leaving the incumbent comparatively unharmed.

A transition-sensitive reform instead asks what must remain invariant while the institutional carrier changes. In mathematics, a useful transformation preserves the structures relevant to the problem while allowing irrelevant representation to vary. The analogous political question is: which relations belong to the people participating in them, and which genuinely belong to the intermediary?

The answer will differ by domain, but the question itself exposes a hidden assumption of platform capitalism. Because a relation is technically instantiated inside privately owned infrastructure, the infrastructure owner is often treated as possessing broad authority over the relation. Yet technical containment does not settle normative ownership. A conversation occurring through a server is not therefore the server owner's conversation. An audience assembled through a platform is not therefore identical with the platform. A professional reputation generated through transactions mediated by a marketplace is not obviously reducible to an asset of that marketplace.

The distinction becomes increasingly important as artificial intelligence systems mediate not only transactions but memory itself. Assistants, recommendation systems, knowledge platforms, and generative interfaces can accumulate long histories of interaction that make later systems more useful precisely because they know the participant's prior context. If such memory becomes nonportable, the switching cost can become extraordinarily intimate. The user would not merely change software; they would repeatedly lose a computationally mediated history of themselves.

The same structural danger appears even if the incumbent system behaves well. A high-quality assistant with excellent long-term memory can create stronger lock-in than a poor assistant precisely because more valuable history accumulates within it. The appropriate response is not to reject persistent memory. It is to design memory so that usefulness does not imply captivity.

The principle can be stated as follows:
\[
\text{the value accumulated through use}
\not\Rightarrow
\text{the right of the intermediary to immobilize that value}.
\]

If anything, the direction of the obligation should often reverse. The more an institution benefits from encouraging participants to accumulate irreplaceable state, the stronger the case for mechanisms ensuring that such accumulation does not become a hostage to future policy changes.

This gives disenshittification a temporal dimension missing from purely competitive accounts. Competition asks whether another provider can enter now. Institutional memory asks whether participants can carry enough of their past into that provider for entry to matter. A market containing ten nominal competitors can remain effectively enclosed if choosing any of the nine alternatives requires abandoning twenty years of accumulated history.

The relevant unit of freedom is therefore not the isolated decision at the present moment. It is the trajectory that can continue through the decision.

\section{The Political Economy of Irreversibility}

This trajectory-sensitive view also clarifies the role of irreversibility. Not every irreversible process is objectionable. Human life is irreversible in innumerable ordinary senses: time passes, opportunities disappear, commitments alter future possibilities, skills are acquired, relationships form, institutions develop, and actions leave consequences. A theory demanding complete reversibility would therefore demand the abolition of history. The relevant concern is instead whether one actor can strategically produce irreversible dependence for others while preserving comparatively reversible options for itself.

Platform relations often contain exactly this asymmetry. Users invest years constructing identities, audiences, archives, workflows, and reputations. The platform can alter interfaces, ranking systems, pricing structures, contractual terms, APIs, or business models comparatively quickly. The participant's investment is slow and path-dependent; the institution's rule change can be rapid and centralized. The temporal asymmetry becomes a source of power.

Let \(C_i(\Delta)\) denote the cost to participant \(i\) of adapting to institutional change \(\Delta\), and \(C_M(\Delta)\) the cost to the intermediary of imposing or reversing the same change. Where
\[
C_i(\Delta)\gg C_M(\Delta),
\]
the intermediary possesses an option-like advantage. It can experiment with changes whose costs are partly externalized onto participants and reverse them if institutional returns prove disappointing, while participants may be unable to reverse the investments made in response.

This resembles a ratchet. A platform changes a rule; participants reorganize themselves around it; some leave, others adapt, and new entrants arrive knowing only the new condition. If the platform later partially reverses the rule, the prior population and practices do not necessarily return. The institutional action has altered the selection environment and therefore the population available for subsequent selection.

Creative labour is particularly vulnerable to this process because careers themselves are path-dependent. A musician who spends several years optimizing for one distribution regime has not merely generated a sequence of posts or tracks. They have allocated finite developmental time toward skills, formats, audiences, and professional relationships selected under that regime. If the platform changes its objective, those investments cannot simply be exchanged for the alternative capacities that would have developed under another trajectory.

The opportunity cost is counterfactual and therefore largely invisible. We can observe the creator who became highly skilled at producing work compatible with the platform. We cannot observe the creator that the same person would have become had the distribution environment rewarded different forms of experimentation. The loss is not necessarily a destroyed artifact. It can be an unrealized capacity.

This is why cultural extraction cannot be measured solely through royalty rates or advertising revenue. Those are important and comparatively measurable quantities, but the platform also allocates developmental incentives. It influences which capacities are worth cultivating. Over long periods, these incentives can reshape the population of available creators and practices.

The same is true of labour platforms. Workers reorganize schedules, acquire equipment, relocate, learn platform-specific strategies, or structure household finances around expected work patterns. The platform can subsequently alter compensation or dispatch rules. The worker's prior adaptation becomes part of the reason the new terms are difficult to refuse.

An irreversible investment therefore creates a vulnerability to hold-up. Once one party has made a relationship-specific investment, the other party can sometimes renegotiate the distribution of surplus because the investment cannot be recovered elsewhere at equal value. Platform lock-in generalizes this familiar economic problem across enormous populations.

The conventional response is contract. Parties anticipate the risk and establish rules limiting opportunistic renegotiation. Yet platform contracts are frequently contracts of adhesion whose terms reserve broad modification powers to the intermediary. The institution can therefore acquire both the benefits of encouraging relationship-specific investment and considerable discretion to change the terms governing the resulting relationship.

The imbalance becomes sharper where the institution can modify the effective rules without formally modifying the contract. A recommendation change, ranking adjustment, interface redesign, API restriction, or shift in enforcement can radically alter the economic value of participation while leaving the nominal agreement largely unchanged. Governance occurs through infrastructure.

This is one reason technical architecture belongs inside political economy. Code can alter the reachable state space more quickly and quietly than legislation or explicit contractual renegotiation. A disabled API endpoint can eliminate an entire class of complementary services. A ranking change can alter thousands of livelihoods. A new default can reorganize behaviour across millions of users. The materiality of software allows governance to be deployed as implementation.

Section 1201 and related anti-circumvention restrictions become especially significant in this context because they can prevent users and competitors from restoring transitions that technical changes removed. The legal system then reinforces a particular technical topology. A path is not merely absent from the interface; reconstructing it can itself become legally hazardous.

Doctorow's emphasis on adversarial interoperability follows from precisely this structure. Historically, competitors have sometimes entered markets not by persuading every participant to rebuild their social or technical world from nothing but by making new systems compatible with installed infrastructure. Compatibility reduces the irreversibility of prior investment. The incumbent's accumulated network ceases to function entirely as a barrier because the entrant can connect to it.

From the incumbent's perspective, such interoperability can look like appropriation. From the perspective developed here, the question is what exactly is being appropriated. If the entrant copies genuinely proprietary creative work or confidential information, the objection is straightforward. If it permits users to continue relationships they themselves formed through an incumbent protocol, the normative picture is different. The incumbent's infrastructure contributed to the formation of those relationships, but contribution does not automatically establish perpetual authority over their future continuation.

The conflict is therefore partly a conflict over historical ownership. Who owns the consequences of coordination?

Platform capitalism often answers: the institution that supplied the infrastructure. A trajectory-sensitive account gives a more plural answer. The resulting state was co-produced by the intermediary and the participants whose activity made the infrastructure valuable. Rights over continuation should reflect that co-production.

This principle also illuminates why network effects deserve neither automatic celebration nor automatic condemnation. Network effects can represent genuine social value: a communications system becomes more useful when more people can communicate through it. The problem arises when the value produced collectively by compatibility is transformed into privately controlled incompatibility with alternatives. The network effect itself is not the chokepoint. The enclosure of the network effect is.

An open protocol can possess enormous network effects without granting one institution equivalent control. Email becomes more useful as more people use compatible systems, yet one provider need not own the entire graph. The web becomes more useful as more documents and services participate, yet its basic architecture historically permitted substantial heterogeneity among clients and servers. These systems have their own concentrations and failures, but they demonstrate that network value and platform sovereignty are conceptually separable.

The distinction is crucial because otherwise the remedy for concentration appears to require destroying the coordination benefits that produced concentration. If the only alternatives are one universal network or many disconnected networks, monopoly can appear technologically inevitable. Interoperability creates a third possibility: shared relational infrastructure with plural institutional realization.

In the language of the present framework, interoperability attempts to preserve the high-value reachability edges while distributing authority over the nodes through which those edges are instantiated.

This is a more ambitious objective than conventional competition but a less destructive one than fragmentation. It does not ask people to sacrifice the network in order to escape the network owner. It asks whether the network can cease to have a single owner in the relevant functional sense.

The same architecture provides a useful model for artificial intelligence. If future assistants become deeply embedded in communication, work, memory, education, and creative practice, requiring every competing system to reconstruct a user's context from zero would create enormous incumbent advantages. Open formats, portable histories, interoperable tools, and separable memory layers could instead permit model competition without demanding autobiographical amnesia from users.

This is not a peripheral design preference. It may determine whether AI assistants become tools users can replace or environments users must inhabit.

The lesson of enshittification is that this distinction should be designed before dependence becomes irreversible. Once an institution has accumulated enough relational state, demands for portability can be framed as attacks upon property, security, privacy, technical integrity, or the business model supporting the service. Some of these objections may be legitimate. But by then the incumbent also possesses substantial influence over which objections become institutionally decisive.

The best moment to preserve another continuation is before anyone knows they will need it.

\section{Hope as Counterfactual Infrastructure}

We can now return to Doctorow's distinction among optimism, pessimism, and hope with greater precision. Optimism and pessimism are predictions. They differ in sign but share a representational form: each estimates the likely trajectory of a system given some model of its present dynamics. Hope is different when understood politically. It does not merely assign a higher probability to a favourable outcome. It intervenes upon the transition structure from which probabilities are being estimated.

Suppose a system currently admits transition probabilities
\[
P(x_{t+1}\mid x_t).
\]
An optimist and pessimist can disagree about this distribution while accepting the same state space and transition relation. Political action of the kind Doctorow describes instead attempts to construct a modified system
\[
P'(x_{t+1}\mid x_t,I),
\]
where \(I\) denotes organization, law, technical intervention, institutional construction, collective bargaining, or another action that changes what transitions are available and how costly they are.

Hope is therefore not confidence about \(P\). It is participation in the construction of \(P'\).

This interpretation avoids the sentimental use of hope as a psychological attitude detached from material conditions. Hope can be infrastructural. A strike fund is stored hope because it changes how long workers can sustain refusal. An interoperable protocol is stored hope because it preserves the possibility of future migration. A public archive is stored hope because it prevents knowledge from disappearing with an institution. A legal right is stored hope when it creates an action that would otherwise be prohibited. A union is stored hope because it converts isolated conditional preferences into collectively executable ones.

These structures matter even when dormant. A fire escape does not become useless because no fire occurred today. Its value partly consists in changing the consequences of a future event. Similarly, a portability right can discipline an incumbent even when most users never exercise it. The credible possibility of migration alters the platform's optimization problem.

This is why evaluating such institutions solely by present utilization can be misleading. Low usage may indicate irrelevance, but it may also indicate successful deterrence or option value. A right to strike is not valuable only during strikes. A right to modify software is not valuable only to people currently modifying it. An interoperability requirement is not valuable only to users currently switching providers. The preserved counterfactual can alter behaviour while remaining unrealized.

We can describe this through option value. Let \(a\) be an alternative continuation whose current direct utility is low because the incumbent behaves acceptably. If preserving \(a\) costs \(c\) while the expected future loss from having no alternative under deterioration is \(L\), then preservation can be rational where
\[
c < P(D)L,
\]
with \(P(D)\) representing the probability of a relevant deterioration. The exact quantities will often be unknowable. The conceptual point is that unused alternatives can possess positive value.

This provides a rigorous response to the claim that interoperability, modification rights, independent archives, or redundant institutions are unnecessary because most people do not currently use them. Their purpose is not merely to satisfy present demand. They preserve the topology within which future demand can become actionable.

The same principle explains why the destruction of alternatives can be profitable even when those alternatives currently have few users. A small competitor can discipline an incumbent disproportionately to its market share if dissatisfied users can migrate toward it. Eliminating or enclosing the competitor therefore removes not only its current transactions but a counterfactual constraint upon future extraction.

This is the economic significance of keeping the outside open. The outside need not dominate the present. It must remain capable of becoming relevant when the present changes.

Such preservation necessarily conflicts with a narrow conception of efficiency. Maintaining redundant systems, legal rights, spare capacity, public institutions, archives, alternative protocols, and independent expertise consumes resources that could be concentrated into the currently dominant system. If efficiency is measured only under present conditions, consolidation can repeatedly win.

But this is an optimization against one realized trajectory. Robustness requires evaluation across a distribution of possible futures.

Let \(U(S,\omega)\) denote the utility of institutional structure \(S\) under future condition \(\omega\). A structure \(S_1\) optimized for the most probable present continuation may satisfy
\[
U(S_1,\omega_0)>U(S_2,\omega_0),
\]
while a more redundant \(S_2\) has higher expected value across uncertain futures:
\[
\mathbb{E}_{\omega}[U(S_2,\omega)]
>
\mathbb{E}_{\omega}[U(S_1,\omega)].
\]
Even expected value may understate the case where catastrophic or irreversible failures deserve asymmetric weight. Robust institutional design therefore cannot be reduced to present throughput.

This is where the political economy of enshittification meets the broader theory of irreversibility. The destruction of an alternative today can eliminate evidence, skills, infrastructure, and organizational memory required to reconstruct it tomorrow. Once destroyed, the option may not be recoverable at its previous cost. The decision to eliminate redundancy is therefore itself path-dependent.

A society that permits every independent bookstore to disappear because a platform distributes books more efficiently cannot assume that independent distribution can be recreated instantly if the platform later abuses its position. A culture that allows local venues, small labels, specialist publications, or public institutions to disappear cannot simply summon them back when algorithmic distribution proves inadequate. Institutions contain tacit knowledge, relationships, habits, and accumulated trust. Their destruction is not equivalent to temporarily switching them off.

The talent reserve has an institutional form here as well. Societies need reserves of people who know how to perform functions outside dominant systems. A fully optimized economy can make such people appear redundant until the dominant system fails. What looked like inefficiency then becomes missing capacity.

This is the opposite of brain drain in the sense developed earlier. The problem is not only that talent leaves one region for another. Talent can remain geographically present while the diversity of capacities through which it can act is drained away. People retain intelligence and skill but become increasingly trained, credentialed, financed, and organized around a narrow set of dominant infrastructures. The reserve shrinks even while nominal talent increases.

A healthy talent reserve is therefore not merely a surplus population waiting to be employed. It is a distributed capacity to instantiate alternative continuations.

This is why the cultivation of seemingly impractical skills, independent technical knowledge, local institutions, open standards, artistic experimentation, and heterogeneous forms of expertise can have political value beyond their immediate output. They preserve ways of doing things that the dominant optimizer has not completely absorbed.

Sprezzatura returns here in its final sense. The appearance of effortless excess, the capacity not fully displayed, the card not played, and the skill not immediately monetized all represent forms of reserve. They deny the present evaluator complete jurisdiction over the future use of capacity. The reserve can later become the basis of refusal, adaptation, or invention.

The demand for total legibility is therefore politically dangerous even when motivated by efficiency. A completely legible population is easier to administer, predict, segment, price, and optimize. It is also easier to exhaust. If every capacity has already been assigned its highest-valued current use, every relationship incorporated into a platform, every preference inferred, every cultural experiment tested against present demand, and every institutional redundancy eliminated, then little remains from which an unanticipated continuation can be built.

Hope requires slack.

Not infinite slack, not romantic waste, and not indiscriminate preservation of everything that exists. It requires enough uncommitted structure that the future is not merely an extrapolation of the present optimizer.

This gives a final interpretation to Doctorow's refusal of both optimism and pessimism. The pessimist mistakes present constraints for destiny. The optimist mistakes favourable extrapolation for security. Hope refuses both errors because it treats constraints themselves as objects of intervention.

The appropriate political question is therefore not simply, ``Will platforms get better?'' It is: ``What would have to remain reachable for their deterioration not to determine everyone else's future?''

That question directs attention away from promises and toward architecture, away from benevolent leadership and toward portability, away from exceptional winners and toward denominators, away from nominal choice and toward transition costs, away from prediction and toward admissibility, and away from the quality of the current ruler toward the survivability of refusal.

A society capable of asking that question has already recovered part of the outside.

\section{Final Thesis}

The argument of this essay can now be compressed into a single claim. Enshittification is the institutional process by which an intermediary converts successful coordination into control over the counterfactual conditions of continued participation.

The process begins with value. Without genuine usefulness, there is little reason for enough people to gather within the intermediary for dependence to become consequential. Successful coordination creates relations; relations create accumulated state; accumulated state creates switching costs; switching costs weaken refusal; weakened refusal permits extraction; extraction finances further control over routing, representation, and information; and control over those structures makes alternatives progressively less visible and less reachable.

At the same time, asymmetric legibility permits the intermediary to learn the reservation thresholds of participants while participants remain uncertain about the rules governing their treatment. Vocational awe lowers the monetary cost of retaining people who care intrinsically about their work. A large talent reserve makes individuals substitutable. Spectacular winners obscure the denominator of unsuccessful attempts. Artificial scarcity concentrates power at the level of routing even when production becomes abundant. Algorithmic allocation shapes the behaviour later interpreted as evidence of preference. Representational narrowing makes human practice increasingly compatible with the metrics through which it is evaluated, and automation can then reproduce the selected projection while appearing to reproduce the practice itself.

The result is not merely monopoly or monopsony, although it can contain both. It is a system in which the intermediary increasingly controls the evidence by which its own necessity is established.

The corrective principle is therefore not simply competition, regulation, transparency, or public ownership considered in isolation. Each can fail if it reproduces the same closure at another level. The deeper requirement is counterfactual plurality under conditions of effective reachability. Participants must retain continuations whose existence does not depend upon the permission of the institution they are meant to constrain.

This requirement explains why portability without interoperability is often insufficient, why transparency without bargaining power can become one-sided legibility, why competition without migration paths can remain nominal, why generosity without democratic warrant can preserve dependency, why automation evaluated solely against platform metrics can mistake projection for equivalence, and why exit without preservation of accumulated relations can amount to little more than permission to disappear.

The political economy of platforms is therefore inseparable from the topology of continuation.

A system is free in the relevant sense not because every participant can imagine another state, but because enough paths toward other states remain materially, legally, technically, and socially executable. A system is corrigible not because criticism is permitted, but because criticism can connect to transitions that alter the institution's incentives. A system is plural not because many alternatives can be named, but because some can survive long enough to become genuine competitors, comparisons, and refuges.

This is what enshittification progressively removes.

And this is what disenshittification must restore.

The central institutional question is consequently not whether an intermediary presently deserves our trust. Trust is necessarily contingent, and institutions outlive the people whose conduct initially earned it. The question is whether trusting the institution today destroys our capacity to stop trusting it tomorrow.

Where the answer is yes, dependence has already begun to exceed coordination.

Where the answer is no, another continuation remains possible.

That remainder is not a defect in optimization. It is the condition under which optimization remains answerable to a world beyond itself.

It is the outside that keeps the system corrigible, the reserve that prevents efficiency from becoming exhaustion, the unchosen branch that keeps the realized branch from becoming necessity, and the counterfactual against which power can still be measured.

It is, finally, the difference between a platform we use and a world we are no longer permitted to leave.

\clearpage
\section*{Bibliography}
\addcontentsline{toc}{section}{Bibliography}

\begin{thebibliography}{99}

\bibitem[Doctorow and Giblin(2022)]{doctorowgiblin2022chokepoint}
Cory Doctorow and Rebecca Giblin (2022). \emph{Chokepoint Capitalism: How Big Tech and Big Content Captured Creative Labor Markets and How We'll Win Them Back}. Boston: Beacon Press.

\bibitem[Doctorow(2023)]{doctorow2023enshittification}
Cory Doctorow (2023). ``The `Enshittification' of TikTok.'' \emph{Wired}. Provisional citation.

\bibitem[Doctorow(2024)]{doctorow2024enshittification}
Cory Doctorow (2024). ``My McLuhan Lecture on Enshittification.'' Expanded account of platform enshittification; publication details to verify.

\bibitem[Doctorow and Fox(2026)]{doctorowfox2026}
Cory Doctorow and Jesse David Fox (2026). ``How Tech Is Making the Life of Artists Impossible.'' \emph{Good One} podcast, hosted by Jesse David Fox (Vulture / Vox Media Podcast Network). \url{https://youtu.be/c0pHEEDGv4U} (accessed 2026-07-25). Episode discussing Doctorow's \emph{Enshittification}; exact release date to verify.

\bibitem[Doctorow(2003)]{doctorow2003magickingdom}
Cory Doctorow (2003). \emph{Down and Out in the Magic Kingdom}. New York: Tor Books. Novel; source of the Whuffie reputation economy and the Bitchun Society, discussed here as fiction rather than as an empirical claim.

\bibitem[Doctorow(2025)]{doctorow2025picksandshovels}
Cory Doctorow (2025). \emph{Picks and Shovels}. New York: Tor Books. Martin Hench series; publication month and edition details to verify.

\bibitem[Chesterton(1904)]{chesterton1904napoleon}
G. K. Chesterton (1904). \emph{The Napoleon of Notting Hill}. London: John Lane.

\bibitem[Strugatsky and Strugatsky(1977)]{strugatsky1971prisoners}
Arkady Strugatsky and Boris Strugatsky (1977). \emph{Prisoners of Power}. Translated by Helen Saltz Jacobson. Introduction by Theodore Sturgeon. First edition. Best of Soviet Science Fiction series. New York: Macmillan Publishing Co., Inc. 286 pp. ISBN 0-02-615160-X. First English-language edition, hardcover, first printing; jacket illustration by Richard Powers. Originally serialized in Russian 1969--1971 as \emph{Obitaemyi Ostrov} (\emph{The Inhabited Island}).

\bibitem[Lewis(1945)]{lewis1945greatdivorce}
C. S. Lewis (1945). \emph{The Great Divorce}. London: Geoffrey Bles.

\bibitem[van Vogt(1945)]{vanvogt1945worldofnulla}
A. E. van Vogt (1945). \emph{The World of Null-A}. New York: Simon \& Schuster. Serialized in \emph{Astounding Science Fiction}, 1945; first book publication 1948; revised edition 1970. Sequel: \emph{The Players of Null-A} (also published as \emph{The Pawns of Null-A}), 1948. Edition consulted to verify.

\bibitem[Stephenson(1995)]{stephenson1995diamondage}
Neal Stephenson (1995). \emph{The Diamond Age: Or, A Young Lady's Illustrated Primer}. New York: Bantam Spectra.

\bibitem[Pachniewski(2023)]{pachniewski2023intersubjectivity}
Pawel Pachniewski (2023). ``The Intersubjectivity Collapse: A Catastrophic or Existential Risk for Any Civilization That Radically Self-Modifies Their Minds and Develops {AGI}.'' Mental Contractions (Substack). \url{https://mentalcontractions.substack.com/p/the-intersubjectivity-collapse} (accessed 2026-07-25).

\bibitem[Ettarh(2018)]{ettarh2018vocationalawe}
Fobazi Ettarh (2018). ``Vocational Awe and Librarianship: The Lies We Tell Ourselves.'' \emph{In the Library with the Lead Pipe}. Source of the concept ``vocational awe''.

\bibitem[Rosenblat(2018)]{rosenblat2018uberland}
Alex Rosenblat (2018). \emph{Uberland: How Algorithms Are Rewriting the Rules of Work}. Oakland: University of California Press.

\bibitem[Gray and Suri(2019)]{gray2019ghostwork}
Mary L. Gray and Siddharth Suri (2019). \emph{Ghost Work: How to Stop Silicon Valley from Building a New Global Underclass}. Boston: Houghton Mifflin Harcourt.

\bibitem[Wood(2020)]{wood2020goodgig}
Alex J. Wood (2020). \emph{Despotism on Demand: How Power Operates in the Flexible Workplace}. Ithaca: Cornell University Press.

\bibitem[Posner and Weyl(2018)]{posner2018radicalmarkets}
Eric A. Posner and E. Glen Weyl (2018). \emph{Radical Markets: Uprooting Capitalism and Democracy for a Just Society}. Princeton: Princeton University Press.

\bibitem[Naidu(2024)]{naidu2024monopsony}
Suresh Naidu (2024). \emph{Monopsony and Labor Market Power}. Provisional background reference; replace or supplement with the most appropriate primary survey.

\bibitem[Azar et al.(2022)]{azarmarinescu2022monopsony}
Jos{\'e} Azar, Ioana Marinescu, and Marshall Steinbaum (2022). ``Labor Market Concentration.'' \emph{Journal of Human Resources}. Bibliographic details to verify.

\bibitem[Wu(2018)]{wu2018curse}
Tim Wu (2018). \emph{The Curse of Bigness: Antitrust in the New Gilded Age}. New York: Columbia Global Reports.

\bibitem[Stoller(2019)]{stoller2019goliath}
Matt Stoller (2019). \emph{Goliath: The 100-Year War Between Monopoly Power and Democracy}. New York: Simon \& Schuster.

\bibitem[Zuboff(2019)]{zuboff2019surveillance}
Shoshana Zuboff (2019). \emph{The Age of Surveillance Capitalism: The Fight for a Human Future at the New Frontier of Power}. New York: PublicAffairs.

\bibitem[Pasquale(2015)]{pasquale2015blacksociety}
Frank Pasquale (2015). \emph{The Black Box Society: The Secret Algorithms That Control Money and Information}. Cambridge, MA: Harvard University Press.

\bibitem[Eubanks(2018)]{eubanks2018automating}
Virginia Eubanks (2018). \emph{Automating Inequality: How High-Tech Tools Profile, Police, and Punish the Poor}. New York: St.\ Martin's Press.

\bibitem[Fourcade and Healy(2024)]{fourcade2024ordinal}
Marion Fourcade and Kieran Healy (2024). \emph{The Ordinal Society}. Cambridge, MA: Harvard University Press.

\bibitem[Scott(1998)]{scott1998seeing}
James C. Scott (1998). \emph{Seeing Like a State: How Certain Schemes to Improve the Human Condition Have Failed}. New Haven: Yale University Press.

\bibitem[Dubal(2023)]{dubal2023algorithmic}
Veena Dubal (2023). ``On Algorithmic Wage Discrimination.'' Provisional citation; verify final publication details.

\bibitem[Calo(2014)]{calo2015digitalmarket}
Ryan Calo (2014). ``Digital Market Manipulation.'' \emph{George Washington Law Review} 82(4), 995--1051.

\bibitem[United States District Court for the District of Columbia(2024)]{usgoogleantitrust2024}
United States District Court for the District of Columbia (2024). ``United States et al.\ v.\ Google LLC.'' Antitrust proceedings and findings. Primary materials concerning Google's search monopoly; insert specific opinion/exhibit citations during citation pass.

\bibitem[United States v.\ Google LLC(2023)]{googlequalityrevenue}
United States v.\ Google LLC (2023). ``Internal Google Search Quality and Revenue Materials.'' Trial exhibits discussed in the Doctorow interview; identify individual exhibits before final publication.

\bibitem[Hesmondhalgh(2019)]{hesmondhalgh2019cultural}
David Hesmondhalgh (2019). \emph{The Cultural Industries}. 4th edition. London: SAGE.

\bibitem[Eriksson et al.(2019)]{eriksson2019spotify}
Maria Eriksson, Rasmus Fleischer, Anna Johansson, Pelle Snickars, and Patrick Vonderau (2019). \emph{Spotify Teardown: Inside the Black Box of Streaming Music}. Cambridge, MA: MIT Press.

\bibitem[Prey(2020)]{prey2020musicians}
Robert Prey (2020). \emph{Knowing Me, Knowing You: Datafication on Music Streaming Platforms}. Provisional reference; replace with precise article/book-chapter citation during verification.

\bibitem[Sandel(2020)]{sandel2020tyranny}
Michael J. Sandel (2020). \emph{The Tyranny of Merit: What's Become of the Common Good?} New York: Farrar, Straus and Giroux.

\bibitem[Frank(2016)]{frank2016success}
Robert H. Frank (2016). \emph{Success and Luck: Good Fortune and the Myth of Meritocracy}. Princeton: Princeton University Press.

\bibitem[Merton(1968)]{merton1968matthew}
Robert K. Merton (1968). ``The Matthew Effect in Science.'' \emph{Science} 159(3810), 56--63.

\bibitem[Hirschman(1970)]{hirschman1970exit}
Albert O. Hirschman (1970). \emph{Exit, Voice, and Loyalty: Responses to Decline in Firms, Organizations, and States}. Cambridge, MA: Harvard University Press.

\bibitem[Olson(1965)]{olson1965collective}
Mancur Olson (1965). \emph{The Logic of Collective Action: Public Goods and the Theory of Groups}. Cambridge, MA: Harvard University Press.

\bibitem[Schelling(1978)]{schelling1978micromotives}
Thomas C. Schelling (1978). \emph{Micromotives and Macrobehavior}. New York: W. W. Norton.

\bibitem[Ostrom(1990)]{ostrom1990governing}
Elinor Ostrom (1990). \emph{Governing the Commons: The Evolution of Institutions for Collective Action}. Cambridge: Cambridge University Press.

\bibitem[Lessig(2006)]{lessig2006code}
Lawrence Lessig (2006). \emph{Code: Version 2.0}. New York: Basic Books.

\bibitem[Wu(2010)]{wu2010master}
Tim Wu (2010). \emph{The Master Switch: The Rise and Fall of Information Empires}. New York: Alfred A. Knopf.

\bibitem[United States Congress(1998)]{dmca1201}
United States Congress (1998). ``17 U.S.C. \S 1201: Circumvention of Copyright Protection Systems.'' Digital Millennium Copyright Act.

\bibitem[Castiglione(1528)]{castiglione1528courtier}
Baldassare Castiglione (1528). \emph{Il libro del cortegiano}. Cite the translation/edition actually consulted in final version.

\bibitem[Lewis(1973)]{lewis1973counterfactuals}
David Lewis (1973). \emph{Counterfactuals}. Cambridge, MA: Harvard University Press.

\bibitem[Pearl(2009)]{pearl2009causality}
Judea Pearl (2009). \emph{Causality: Models, Reasoning, and Inference}. 2nd edition. Cambridge: Cambridge University Press.

\end{thebibliography}

\end{document}
